\documentclass[11pt]{amsart}
\usepackage{amsmath,amssymb,amsthm,mathtools,mathrsfs}
\usepackage[margin=1.08in]{geometry}
\usepackage{booktabs}
\usepackage{array}
\usepackage{longtable}
\usepackage{tabularx}
\usepackage{hyperref}
\hypersetup{hidelinks}
\usepackage[final,expansion=false]{microtype}
\setlength{\parskip}{4pt plus 1pt}
\emergencystretch=2.5em

\newtheorem{theorem}{Theorem}[section]
\newtheorem{proposition}[theorem]{Proposition}
\newtheorem{lemma}[theorem]{Lemma}
\newtheorem{corollary}[theorem]{Corollary}
\theoremstyle{definition}
\newtheorem{definition}[theorem]{Definition}
\newtheorem{problem}[theorem]{Problem}
\newtheorem{conjecture}[theorem]{Conjecture}
\theoremstyle{remark}
\newtheorem{remark}[theorem]{Remark}

\newcommand{\C}{\mathbb C}
\newcommand{\R}{\mathbb R}
\newcommand{\N}{\mathbb N}
\newcommand{\Z}{\mathbb Z}
\newcommand{\Q}{\mathbb Q}
\newcommand{\A}{\mathbb A}
\newcommand{\F}{\mathbb F}
\newcommand{\Rel}{\operatorname{Re}}
\newcommand{\Iml}{\operatorname{Im}}
\newcommand{\Spec}{\operatorname{Spec}}
\newcommand{\Frob}{\operatorname{Frob}}
\newcommand{\Sh}{\operatorname{Sh}}
\newcommand{\Prim}{\operatorname{Prim}}
\newcommand{\supp}{\operatorname{supp}}
\newcommand{\vol}{\operatorname{vol}}
\newcommand{\Hom}{\operatorname{Hom}}
\newcommand{\Cone}{\operatorname{Cone}}
\newcommand{\degdet}{\operatorname{deg}_{\mathrm{det}}}
\newcommand{\Ch}{\operatorname{Ch}}
\newcommand{\Knat}{K^{\natural}}
\newcommand{\Drad}{\mathcal D_A^{\mathrm{rad}}}
\newcommand{\DA}{\mathcal D_A}
\newcommand{\Doff}{\mathcal D_A^{\mathrm{off}}}
\newcommand{\Wfin}{W_{\mathrm{fin}}}
\newcommand{\cfin}{\mathrm{c}_{\mathrm{fin}}}
\newcolumntype{Y}{>{\raggedright\arraybackslash}X}

\begin{document}

\title[An arithmetic Frobenius family and its integral shadow]
{An Arithmetic Frobenius Family and its Integral Shadow\\
on a Correspondence Algebra of the Idele Class Group\\[4pt]
\large Inverse semigroup structure, two Gram forms,
and the von Mangoldt function as a shell volume}

\author{David Alejandro Trejo Pizzo}
\thanks{Independent Researcher. \texttt{dtrejopizzo@gmail.com}}

\date{\today}

\begin{abstract}
For an integer $n\ge2$, choose a semilocal idele class group whose finite
places contain the prime divisors of $n$. Convolution of the normalized
valuation-shell measures defines a Frobenius correspondence $\Frob_n$ of
bidegree $(n^{1/2},n^{-1/2})$. These correspondences multiply according to
the product of their indices. Their finite-place contact with the diagonal
has coefficient $\Lambda(n)/\sqrt n$: a shell for $n=p^k$ has multiplicative
Haar volume $\log p$, whereas an index with two distinct prime divisors
meets no single-place shell.

The integral model is the complex $[L_S\xrightarrow{n}L_S]$. Its
M\"obius cofiber records the cyclotomic value $\Phi_n(1)$. We compute the
Gram form both on the Frobenius family and after averaging over the unit
group. The first calculation gives the finite-place kernel; the radial
average also recovers the archimedean term.
\end{abstract}

\maketitle

\tableofcontents

\section{Introduction}

Let $S$ be a finite set of places of $\Q$ containing the archimedean one, let
$S_f=S\smallsetminus\{\infty\}$, and let $C_S$ be the associated idele class
group.  In \cite{Companion} we constructed on $C_S$ an algebra of
correspondences with a composition $\star$, an involution $\dagger$, a
canonical half-density $|\cdot|_S^{-1/2}$, and an analytic trace $T$ with
explicit local terms: at a finite place $p$ the trace integrates a function
along the slice $j_p:\Q_p^\times\to C_S$ against
$\tau_p(x)=|x|_p^{1/2}/|1-x|_p$ and the multiplicative Haar measure, and at
$\infty$ it integrates against the density $K(t)=e^{-t/2}/(1-e^{-2t})$ on
tests even under the sign (and against its two halves $K_s$, $K_c$ on the two
real half-lines in general).  In
this paper we construct arithmetic objects inside that algebra and compute
their invariants.

The objects are elementary.  For a prime $p$ and $k\ge1$ the set
$p^k\Z_p^\times$ of elements of valuation $k$ is compact; its image under the
slice map is a compact subset of $C_S$, and we let
\[
   \Frob_{p^k}:=\text{the normalized Haar measure of } j_p(p^k\Z_p^\times),
\]
a probability measure on $C_S$.  For $n=\prod p^{k_p}$ we set
$\Frob_n:=\star_{p\mid n}\Frob_{p^{k_p}}$, and $\Frob_1:=\Delta$, the point
mass at the identity class.  The family has the following properties.
\begin{itemize}
\item It is strong monoidal with a genuine unit: $\Frob_m\star\Frob_n
=\Frob_{mn}$ as an equality of measures, and $\Frob_1$ is the diagonal
itself, a point mass and not a projector (Theorem~\ref{thm:frobmonoidal}).
\item It remembers the arithmetic.  The support of $\Frob_n$ determines $n$
and the set $P(n)$ of primes dividing $n$ (Theorem~\ref{thm:frobmonoidal}(iv)),
and no weak-$*$ continuous one-parameter family interpolates the $\Frob_n$
(Remark~\ref{rem:nointerp}).  By contrast, the \emph{radial} objects of
Section~\ref{sec:radial} factor through $(\R_{>0},\times)$ and forget
factorization entirely (Lemma~\ref{lem:notideal}).
\item Its closure under $\dagger$ is an inverse semigroup, not a group.  The
smallest $\star$- and $\dagger$-closed family containing the $\Frob_n$ is
indexed by pairs $(r,F)$ with $r\in\Q_{>0}^\times$ and $F$ a finite set of
primes containing those dividing $r$; its idempotents are the Haar measures
$\Delta^{(F)}$ of compact subgroups $K_F$, which form the semilattice of finite
sets of primes under union (Theorem~\ref{thm:invsemigroup}).  The assignment
$a/b\mapsto\Frob_a\star\Frob_b^\dagger$ does not descend to $\Q_{>0}^\times$.
\item Its contact with the diagonal is the von Mangoldt function.  For a
smoothing of $\Frob_n$ at any level $N$ one has, with no regularization and
no limit at the finite places,
\[
   \sum_{q\in S_f}W_q\bigl(\Frob_n^{[h,N]}\bigr)=\Lambda(n)\,n^{-1/2}h(1/n)
\]
(Theorem~\ref{thm:frobcontact}).  Here $\Lambda(p^k)=\log p$ is the
multiplicative Haar volume $\vol_{d^*x}(p^k\Z_p^\times)$ of the support of the
object, and the vanishing at indices $n$ with at least two distinct prime divisors is the disjointness of that support
from every valuation shell (Theorem~\ref{thm:compositedisjoint}).
\end{itemize}

Over $\C$ the arithmetic torsion is invisible: in a $\C$-linear additive
category $p\cdot\mathrm{id}$ is invertible, so no additive functor produces a
complex with cohomology $\Z/p$.  What can be constructed is an integral shadow
defined directly on the objects of the Frobenius monoid, using the lattice
$L_S=\prod_{p\in S_f}\Z_p$:
\[
  \Sh_\Z(\Frob_n):=\Cone\bigl(L_S\xrightarrow{\ n\ }L_S\bigr).
\]
It depends only on the measure $\Frob_n$ (Lemma~\ref{lem:shwelldef}), it is a
strict functor on the divisor poset (Lemma~\ref{lem:shfunctor}), its
transitivity data form an exact triangle that does not split, and its
M\"obius cofiber over the divisor poset is $[\Z\xrightarrow{\Phi_n(1)}\Z]$
(Theorem~\ref{thm:shstructure}).  The resulting identities of degrees,
\[
  \cfin(\Frob_n)=n^{-1/2}\degdet\Prim(\Sh_\Z)_n,\qquad
  \degdet\Sh_\Z(\Frob_n)=\log n=\sum_{d\mid n}d^{1/2}\cfin(\Frob_d),
\]
are Theorem~\ref{thm:shdegrees}; both sides of each are invariants of the same
objects.  The complexification $[\C\xrightarrow{n}\C]$ of the shadow is
acyclic, so these are identities of degrees and not of complexes.

The last part of the paper computes Gram forms.  Let
$W_{\mathrm{fin}}=\sum_{q,k}(\log q)q^{-k/2}(\delta_{k\log q}+\delta_{-k\log
q})$ and $\omega_\infty=K(|t|)\,dt$; their sum is the kernel of Weil's
quadratic form.  The Gram form of the Frobenius family carries
$W_{\mathrm{fin}}$ off the diagonal and has no archimedean part: the
archimedean entry is multiplied by the congruence indicator
$\prod_{s\notin F}\mathbf 1[s^N\mid (n-m)]$, which vanishes for all large $N$
(Theorem~\ref{thm:gramoff}, Corollary~\ref{cor:onlyfinite}), because
$\supp\Frob_n$ does not meet the archimedean slice
(Corollary~\ref{cor:frobnoarch}).  The Gram form of the unit-averaged family
$\Frob_n\star\Delta^{\mathrm{rad}}=\Gamma_n^{\mathrm{rad}}$ is the whole kernel
$W_{\mathrm{fin}}+\omega_\infty$, at every level and with no side conditions
(Theorem~\ref{thm:fullweil}).  So the unit-averaging projector converts the
finite part of the kernel into the whole of it, at the cost of forgetting
factorization (Theorem~\ref{thm:twohalves}).

Three further results complete the picture.  For an $L^2$-normalized radial
approximate identity of width $\varepsilon$, the self-pairing is
$-\log(1/\varepsilon)+O(1)$ (Theorem~\ref{thm:infgenus}); in Weil's dictionary
$\Delta\cdot\Delta=2-2g$ this corresponds to $g=\infty$.  The finite-place
form alone is strictly positive on a primitive class, explicit up to one
sufficiently large parameter, for every finite
$S$ (Theorem~\ref{thm:qfinpositive}).  And the normalization
$\vol_{d^*x}(\Z_p^\times)=\log p$, which makes $\Lambda$ a volume, is the only
local normalization, up to one global scale, for which the local Tate residue
is the same at every place (Theorem~\ref{thm:tateresidue}).

No statement about the zeros of the zeta function is used.  The only statement
about them is Lemma~\ref{lem:singlefreq}, a self-contained computation of the
contribution of a single resonant quadruple of zeros to a Hermitian zero sum
against a single-frequency test; it is not used elsewhere.  No sign of the
Gram forms is asserted.

\subsection*{Place in the series}
We use the algebra, involution and trace of \cite{Companion}. The radial comparison theorem \cite{P45} uses the contact coefficients computed here. The sign problem in \cite{P47} concerns the complete radial form, including the archimedean term.

Section~\ref{sec:prelim} fixes notation and recalls what we need from
\cite{Companion}.  Sections~\ref{sec:radial} and \ref{sec:radtrace} treat the
radial family and its trace, Sections~\ref{sec:frob} and \ref{sec:frobcontact}
the Frobenius family and its contact, Section~\ref{sec:integral} the
$\Z$-linear models, and Section~\ref{sec:shadow} the integral shadow.
Sections~\ref{sec:gram}--\ref{sec:normalization} contain the Gram forms, the
self-intersection of the diagonal and the normalization theorem, and
Section~\ref{sec:approx} collects the approximation lemmas and the
single-frequency computation.

\section{Notation and background}\label{sec:prelim}

$\Lambda$ denotes the von Mangoldt function: $\Lambda(n)=\log p$ if $n=p^k$
with $k\ge1$, and $\Lambda(n)=0$ otherwise, so $\Lambda(1)=0$. $\mu$ is the
M\"obius function, $P(n)$ the set of distinct primes dividing $n$, and
$V(m)\subset\Spec\Z$ the closed set of primes dividing $m$. $\Phi_n$ is the
$n$-th cyclotomic polynomial. Absolute values are normalized so that the
product formula $\prod_v|x|_v=1$ holds for $x\in\Q^\times$, with
$|p|_p=1/p$.

We call a finite set $S$ of places of $\Q$ with $\infty\in S$ a
\emph{chart}; we write $S_f=S\smallsetminus\{\infty\}$ and
\[
  L_S:=\prod_{p\in S_f}\Z_p .
\]
When an index $n$ is under discussion the chart is always assumed to contain
every prime dividing $n$, unless the contrary is said explicitly. On $\Q_p^\times$
the multiplicative Haar measure $d^*x$ is normalized by
\begin{equation}\label{eq:shellvol}
   \vol_{d^*x}(\Z_p^\times)=\log p ,
\end{equation}
and on $\R^\times$ by $d^*x_\infty=dx/(2|x|)$; \S\ref{sec:normalization} proves
that \eqref{eq:shellvol} is forced, up to one global scale, by the requirement
that the local Tate residue be the same at every place. By translation
invariance $\vol_{d^*x}(p^k\Z_p^\times)=\log p$ for every $k\in\Z$.

We now recall the facts from \cite{Companion} that are used below, in the
form in which they are used.

\begin{proposition}[The valuation splitting; {\cite[Lemma 2.2]{Companion}}]
\label{prop:splitting}
Let $S$ be a chart and $\A_S^\times$ the group of $S$-ideles. Let
$\varphi(x)=\prod_{p\in S_f}p^{v_p(x_p)}$ and let
$\Gamma_S=\Z_S^\times$ be the group of positive rationals supported on $S_f$.
Then $\A_S^\times=\Gamma_S\times\ker\varphi$ and the idele class group is
\[
  C_S\;\cong\;\ker\varphi\;=\;\R^\times\times\prod_{p\in S_f}\Z_p^\times .
\]
The modulus $|\cdot|_S$ restricted to $C_S$ is $|x|_S=|x_\infty|$, its kernel
is the maximal compact subgroup
\begin{equation}\label{eq:Knat}
   \Knat:=\{\pm1\}\times\prod_{p\in S_f}\Z_p^\times\subset C_S ,
\end{equation}
and $C_S/\Knat\cong\R_{>0}$. Haar measure on $C_S$ disintegrates as
$d^*w=dt\otimes du$ in the coordinate $t=\log|w|_S$, with $du$ the probability
Haar measure of $\Knat$.
\end{proposition}

\begin{proposition}[Slice maps; {\cite[\S6, (32)]{Companion}}]\label{prop:slices}
For $q\in S_f$ and $k\in\Z$, $w\in\Z_q^\times$, and for $x\in\R^\times$,
\begin{align*}
  j_q(q^kw)&=\bigl(q^{-k}\ \text{at}\ \infty;\ w\ \text{at}\ q;\
     q^{-k}\ \text{at every other place of}\ S\bigr),\\
  j_\infty(x)&=\bigl(x\ \text{at}\ \infty;\ 1\ \text{at every finite place of}\
     S\bigr).
\end{align*}
Each $j_v:\Q_v^\times\to C_S$ is an injective continuous homomorphism; indeed
$v_r(q^{-k})=0$ for $r\ne q$, so $j_q$ lands in $\ker\varphi$.
\end{proposition}

\begin{proposition}[The test algebra and the trace;
{\cite[\S6, Def.~6.5, Conv.~6.6, Thm.~6.7, \S11, Thm.~11.6]{Companion}}]\label{prop:trace}
Let $\DA:=C_c^\infty(C_S)$, meaning compactly supported functions which are
smooth in the archimedean coordinate and locally constant in the compact
coordinates. Convolution $\star$ on $C_S$ makes $\DA$ an algebra, the
involution is $F^\dagger(x)=\overline{F(x^{-1})}$, and $F\mapsto E_F$ realizes
$\DA$ inside the correspondence algebra. There is a linear functional
$T:\DA\to\C$, the analytic trace, with local terms
\begin{equation}\label{eq:localterms}
  W_q(F)=\sum_{k\ne0}\int_{q^k\Z_q^\times}(F\circ j_q)\,\tau_q\,d^*x
   +(\log q)\int_{\Z_q^\times}\frac{(F\circ j_q)(u)-(F\circ j_q)(1)}{|1-u|_q}
     \,du,
\end{equation}
\begin{equation}\label{eq:archterm}
\begin{aligned}
  W_\infty(F)={}&(\gamma+\log4\pi)\,F(1)
   +\int_0^\infty\Bigl[K_s(t)\bigl(b_+(t)+b_+(-t)\bigr)
     -2K(t)F(1)e^{-t/2}\Bigr]dt\\
   &+\int_0^\infty K_c(t)\bigl(b_-(t)+b_-(-t)\bigr)dt,
\end{aligned}
\end{equation}
where $b_\pm(t):=F(j_\infty(\pm e^{t}))$ are the profiles of $F$ on the two
real half-lines, $\tau_v(x)=|x|_v^{1/2}/|1-x|_v$,
$K_s(t)=1/(4\sinh(t/2))$, $K_c(t)=1/(4\cosh(t/2))$ and
$K=K_s+K_c=e^{-t/2}/(1-e^{-2t})$ (\cite[Def.~6.5 and Convention~6.6]{Companion}). For $F$ even under the sign, $b_+=b_-=:b$ and
\eqref{eq:archterm} is $(\gamma+\log4\pi)F(1)+\int_0^\infty
K(t)[b(t)+b(-t)-2F(1)e^{-t/2}]dt$; and $T(F)=\sum_{v\in S}W_v(F)+c\,F(1)$; throughout
this paper the constant $c$ is taken to be $0$. We extend the notation $E$ to
points: for $x\in C_S$ we write $E_x$ for the point mass $\delta_x$, a
multiplier of $\DA$, and $E_a:=E_{[a]}$ for an idele $a$ with class $[a]$;
since push-forward of point masses along the multiplication of $C_S$ is
multiplicative, $E_x\star E_y=E_{xy}$ and $E_1=\Delta$. Moreover:
\begin{enumerate}
\item[(a)] \emph{(Exact orbital values.)} $\tau_q\equiv q^{-|k|/2}$ on the
shell $q^k\Z_q^\times$, $k\ne0$.
\item[(b)] \emph{(Absolute convergence off the identity.)} If $F$ vanishes on
a neighbourhood of the identity class --- we write $F\in\Doff$ --- then
$T(F)=\sum_{v\in S}\int_{\Q_v^\times}F(j_v(x))\tau_v(x)\,d^*x$, an absolutely
convergent sum with no counterterm and no dependence on $c$; the archimedean
term reduces to $\int_0^\infty K_s(t)[b_+(t)+b_+(-t)]\,dt
+\int_0^\infty K_c(t)[b_-(t)+b_-(-t)]\,dt$, which is
$\int_0^\infty K(t)[b(t)+b(-t)]\,dt$ when $F$ is even under the sign
(see the paragraph after Proposition~\ref{prop:trace}).
\item[(c)] \emph{(Chart independence and the colimit.)} The charts form a
directed system under inflation, and there is a chart-independent trace
$T_\infty$ on the colimit $\DA^\infty$, for which the finite-place sum runs
over all primes; for a test whose support radius is smaller than
$\log\ell_{\mathrm{next}}(S)$, $T_S$ and $T_\infty$ agree, and in general they
differ by the finite sum of the missing prime terms.
\item[(d)] \emph{(The diagonal is a multiplier.)} $\Delta=E_1$ is a multiplier
of the algebra and not an element of it; $T(\Delta)$ is undefined.
\end{enumerate}
\end{proposition}

\medskip\noindent\phantomsection\label{rem:signodd}%
\emph{Tests that are not even under the sign.}
The positive and negative real half-lines are weighted differently by
$\tau_\infty$: $\tau_\infty(e^{\pm t})=2K_s(t)$ and
$\tau_\infty(-e^{\pm t})=2K_c(t)$ for $t>0$, with $d^*x_\infty=dt/2$ on each
half-line. Formula \eqref{eq:archterm} is the definition of $W_\infty$ on all
of $\DA$ in \cite{Companion}, with no parity assumption, and (b) is the
statement that off the identity it equals $\int F(j_\infty(x))\tau_\infty(x)
\,d^*x_\infty$. The radial classes of \S\ref{sec:radial} are even under the
sign, so only $K$ enters there. The smeared Frobenii of
\S\ref{sec:frobcontact} and their convolutions are supported in
$x_\infty>0$, so $b_-=0$ and only $K_s$ enters their archimedean terms.
\medskip

\begin{proposition}[The Tate functionals and the pairing;
{\cite[Thms.~7.1, 7.2, 7.4, 8.4]{Companion}}]\label{prop:tate}
Define $\lambda_\mp(\mu):=\int_{C_S}|x|_S^{\mp1/2}\,d\mu(x)$ and
$M:=(\lambda_-,\lambda_+)$. Then $\lambda_\mp(F^\dagger)
=\overline{\lambda_\pm(F)}$. The bilinear form
$I_A(F,G):=T(F\star G^\dagger)$ is nondegenerate on $\DA$, and in Fourier
terms
\begin{equation}\label{eq:planch}
  I_A(F,G)=\int_{\widehat{C_S}}\widehat F(\chi)\,\overline{\widehat G(\chi)}\,
    \widehat T(\chi)\,d\chi .
\end{equation}
On the principal component $\mathfrak c_0$ of $\widehat{C_S}$ --- the
characters trivial on $\Knat$, parametrized by $\chi(w)=|w|_S^{i\tau}$ --- the
symbol is
\begin{equation}\label{eq:symbol}
  \widehat W_\infty(0,\tau)=\log\pi
    -\Rel\psi\!\left(\tfrac14+\tfrac{i\tau}2\right),
  \qquad
  \widehat W_p(\mathfrak c_0,\tau)=(\log p)\cdot2\Rel\!\left(
    \frac{\zeta_p}{1-\zeta_p}\right),
\end{equation}
\begin{equation}\label{eq:zetap}
  \zeta_p=p^{-1/2}e^{i(\tau\log p+\theta_p)},
\end{equation}
with the constant phase $\theta_p=\arg\prod_{q\in S_f,\ q\ne p}\kappa_q(p)$
determined by the finite components $\kappa_q$ of the character; on the
principal component every $\kappa_q$ is trivial, so
\begin{equation}\label{eq:nophase}
  \theta_p=0\quad\text{and}\quad\zeta_p=p^{-1/2+i\tau}\ \text{ on }
  \mathfrak c_0 .
\end{equation}
$\psi$ the digamma function; $\widehat T$ is real-analytic and even on
$\mathfrak c_0$, and on each component $|\widehat T(\mathfrak c,\tau)|\le
A_{\mathfrak c}+B_S\log(2+|\tau|)$ with $A_{\mathfrak c}=A_S+\sum_{p\in S_f}
c(\kappa_p)\log p$ depending on the conductors of the component
(\cite[Lemma~8.12, (76)]{Companion}): the bound is uniform on finite sets of
components and not over all of $\widehat{C_S}$. Furthermore $\max_\tau\widehat W_\infty(0,\tau)
=\log\pi-\psi(1/4)\approx5.372$.
\end{proposition}

\begin{proposition}[The completion and its splitting;
{\cite[Lemma~8.12, Thms.~8.11, 8.14]{Companion}}]\label{prop:completion}
Write $J(F):=(I_A(F,\cdot),M(F))\in\DA'\oplus B_T$ and let
$\mathcal N_A^{\mathrm{dist}}:=\overline{J(\DA)}$. There is a seminorm
$\mathcal N$ on $\DA$, bounded on bounded sets, with
$|I_A(F,G)|\le\sup_\chi|\widehat F|\cdot\mathcal N(G)$; and
$\mathcal N_A^{\mathrm{dist}}$ splits canonically as
$\mathcal N_A^{\mathrm{dist,prim}}\oplus\sigma(B_T)$, the splitting sending
$x=(\ell,m)$ to $(\ell,0)+\sigma(m)$. For every bounded $B\subset\DA$ there
are $N_B$, $d_{N_B}$, $C''_B$ such that every $G\in B$ is
$U_{N_B}$-invariant, $\widehat G$ is supported on $d_{N_B}$ components, and
$|\widehat G(\mathfrak c,\tau)|\le C''_B\min(1,\tau^{-2})$. We put
$A_B:=\max A_{\mathfrak c}$ over these $d_{N_B}$ components
(Proposition~\ref{prop:tate}); then $|\widehat T|\le A_B+B_S\log(2+|\tau|)$
wherever some $\widehat G$, $G\in B$, is nonzero.
\end{proposition}

\begin{proposition}[The shadow at a non-unit shell;
{\cite[Lemma~5.7, Thm.~5.12(c), Prop.~6.13(a)]{Companion}}]\label{prop:shadow}
At a dilation $x\in\Q_p^\times$ with $v_p(x)\ne0$ (the shell $p^k\Z_p^\times$,
$k\ne0$) the local twisted Hochschild complex $C(E_p,(E_p)_x)$ is homotopy
equivalent to the rank-one class $\C[0]$ of the fixed point $0$, independent
of $x$; the fixed-point weight is $1/|1-x|_p$ and the half-density is
$|x|_p^{1/2}$, these appearing separately in
$\tau_p(x)=|x|_p^{1/2}|1-x|_p^{-1}=p^{-|v_p(x)|/2}$; and the local weight of
the shell in $W_p$ is its volume $\log p$. In particular the local factor at
a finite place $p$ of an external class with $v_p(w_p)\ne0$ is rank one, so no
integral torsion survives the passage to the $\C$-linear category. The trace
$T$ is defined by the densities $\tau_v$ directly and is \emph{not} derived
from this complex; \cite[Prop.~6.13]{Companion} is a comparison only.
\end{proposition}

\begin{proposition}[No $\Z/p$ from a $\C$-linear additive functor]
\label{prop:noZp}
Let $\mathcal A$ be a $\C$-linear additive category and $\mathcal F$ an
additive functor to complexes of abelian groups. Then no object of
$\mathcal A$ has $H^0\mathcal F(X)\cong\Z/p$.
\end{proposition}

\begin{proof}
In a $\C$-linear additive category $p\cdot\mathrm{id}_X$ is invertible, with
inverse $p^{-1}\cdot\mathrm{id}_X$. An additive functor carries
$p\cdot\mathrm{id}_X$ to $p\cdot\mathrm{id}_{\mathcal F(X)}$ and preserves
invertibility, so $p$ acts bijectively on $H^0\mathcal F(X)$. But $p$ acts as
$0$ on $\Z/p$, which is nonzero.
\end{proof}

This is why every comparison in this paper between an analytic contact and an
integral complex is a comparison of degrees and not of objects.

\section{The radial family}\label{sec:radial}

Two families of objects on $C_S$ must be distinguished.  This section
constructs the first, the \emph{radial} family, and shows that as a family of
objects it carries no arithmetic information.

\begin{lemma}\label{lem:classunit}
Let $n\ge1$ and $S\supseteq\{\infty\}\cup P(n)$, and let
$a_S(n)=(1\ \text{at}\ \infty;\ n\ \text{at each}\ p\in S_f)$. In the
coordinates of Proposition~\ref{prop:splitting} the class of $a_S(n)$ is
\begin{equation}\label{eq:deltaS}
  \delta_S(1/n):=a_S(n)/\varphi(a_S(n))
   =\bigl(1/n\ \text{at}\ \infty;\ 1\ \text{at every}\ p\in S_f\bigr),
\end{equation}
and $|\delta_S(1/n)|_S=1/n$.
\end{lemma}

\begin{proof}
$\varphi(a_S(n))=\prod_{p\in S_f}p^{v_p(n)}=n$, which is a positive $S$-unit
because $S\supseteq P(n)$; so $n\in\Gamma_S$ and $a_S(n)/n$ is the canonical
representative of the class furnished by the splitting
$\A_S^\times=\Gamma_S\times\ker\varphi$. Its archimedean coordinate is $1/n$
and its $p$-coordinate is $n/n=1$ for every $p\in S_f$. Finally
$|x|_S=|x_\infty|$ on $\ker\varphi$.
\end{proof}

The finite coordinates of the class are $1$, not $n$: the integer survives only
in the archimedean coordinate and in the modulus.  In other words, the class
of an arithmetic idele lies on the archimedean slice, and this controls much
of what follows.

\begin{proposition}\label{prop:noslice}
Let $n\ge2$ and $|S_f|\ge2$. With the slice maps of
Proposition~\ref{prop:slices}:
\begin{enumerate}
\item[(i)] $\delta_S(1/n)=j_\infty(1/n)\in\Iml\,j_\infty$;
\item[(ii)] $\delta_S(1/n)\notin\Iml\,j_q$ for \emph{every} $q\in S_f$, both
for $q\mid n$ and for $q\nmid n$;
\item[(iii)] there is a neighbourhood basis of $\delta_S(1/n)$ in $C_S$
consisting of open sets disjoint from $\bigcup_{q\in S_f}\Iml\,j_q$.
\end{enumerate}
\end{proposition}

\begin{proof}
(i) is the formula for $j_\infty$ with $x=1/n$.

(ii) Suppose $\delta_S(1/n)=j_q(q^kw)$. Pick $r\in S_f\smallsetminus\{q\}$,
possible since $|S_f|\ge2$. The $r$-coordinate of $j_q(q^kw)$ is $q^{-k}$ and
that of $\delta_S(1/n)$ is $1$, so $q^{k}=1$ in $\Z_r^\times$, i.e. $k=0$.
Then the archimedean coordinate of $j_q(w)$ is $q^0=1$, while that of
$\delta_S(1/n)$ is $1/n\ne1$. Contradiction.

(iii) Fix $q\in S_f$ and $r\in S_f\smallsetminus\{q\}$. Consider the basic
open sets
\[
  N_{\varepsilon,m}:=\{x:|x_\infty-1/n|<\varepsilon\}\times
     \prod_{p\in S_f}(1+p^m\Z_p),
\]
a neighbourhood basis of $\delta_S(1/n)$, the classes being written
multiplicatively against the representative \eqref{eq:deltaS}. A point
$j_q(q^kw)\in N_{\varepsilon,m}$ must have
$q^{-k}\in(1/n-\varepsilon,1/n+\varepsilon)$; for $\varepsilon<1/(2n)$ at most
one integer $k$ qualifies, and there is such a $k$ only if $q^{-k}$ is that
close to $1/n$. If no such $k$ exists the intersection is empty. If $k$
exists, the $r$-coordinate of the point is $q^{-k}$, and we need
$q^{-k}\in1+r^m\Z_r$. Now $q^{-k}\ne1$: otherwise $k=0$ and
$|1-1/n|<\varepsilon<1/(2n)$, which is false for $n\ge2$. Since $r\ne q$ we
have $q^{-k}\in\Z_r^\times$, so $v_r(q^{-k}-1)$ is a finite integer $m_0$, and
taking $m>m_0$ makes the condition fail. Doing this for the finitely many
$q\in S_f$ and taking the largest $m$ gives (iii).
\end{proof}

\begin{remark}\label{rem:minimalchart}
Proposition~\ref{prop:noslice}(ii)--(iii) \emph{fails} for $|S_f|=1$: if
$S_f=\{p\}$ and $n=p^k$ then $j_p(p^k\cdot1)=(p^{-k}\ \text{at}\ \infty;\
1\ \text{at}\ p)=\delta_S(1/n)$ exactly. Since $S=\{\infty,p\}$ is the minimal
admissible chart for $n=p^k$, this case is not vacuous, and the hypothesis
$|S_f|\ge2$ is carried in every statement that uses the proposition. The
conclusion which matters, however, holds on every chart: even at $|S_f|=1$ the
\emph{point} class has finite contact $0$, because by \eqref{eq:localterms} the
shell term integrates over $p^k\Z_p^\times$ against $d^*x$, whose restriction
to the shell is $\log p$ times the \emph{probability} Haar measure of
$\Z_p^\times$, and a point mass at the single unit $u=1$ is a null set for it.
Only the unit-spread object --- the norm class of
Definition~\ref{def:normclass}, which is Haar on the whole shell --- receives
the shell volume. So on every chart the point class has no finite contact while its unit
average does, although for $|S_f|=1$ the reason is different.
\end{remark}

\begin{corollary}\label{cor:liftnotcontact}
Let $n\ge2$, $|S_f|\ge2$, and let $F\in\DA$ be supported in one of the
neighbourhoods of Proposition~\ref{prop:noslice}(iii). Then $W_q(F)=0$ for
every $q\in S_f$, and $T(F)=W_\infty(F)$.
\end{corollary}

\begin{proof}
By \eqref{eq:localterms}, $W_q(F)$ is an integral of $F\circ j_q$ over
$\Q_q^\times$. If $\supp F$ misses $\Iml\,j_q$ then $F\circ j_q\equiv0$ and
$W_q(F)=0$. Proposition~\ref{prop:noslice}(iii) provides such supports. With
$c=0$ the trace then reduces to $W_\infty(F)$.
\end{proof}

The consequence for the local factors of the \emph{idele} $a_S(n)$ is recorded
in Theorem~\ref{thm:separation} below: the numbers $\prod_{q\nmid n}
|1-n|_q^{-1}=n-1$ and $\prod_{p\mid n}|n|_p^{1/2}=n^{-1/2}$ are canonical
invariants of a chosen \emph{lift}, and the first of them is not a contribution
of the class to any local term of $T$.

\begin{lemma}\label{lem:notideal}
Let $S\supseteq\{\infty\}\cup P(n)$. Then $n\in\Gamma_S$, so $n\Z_S=\Z_S$ is
the unit ideal. Moreover the assignment $n\mapsto\Gamma_n^{\mathrm{rad}}$ of
Definition~\ref{def:normclass} below factors as
\begin{equation}\label{eq:factorsthrough}
  (\N_{\ge1},\times)\ \hookrightarrow\ (\R_{>0},\times)\ \cong\
  \{\text{level-set Haar measures on }C_S\},
\end{equation}
the second map an isomorphism of topological monoids; so as an \emph{object}
$\Gamma_n^{\mathrm{rad}}$ forgets the factorization of $n$ completely, and
$\Gamma_6^{\mathrm{rad}}$ is structurally indistinguishable from
$\Gamma^{\mathrm{rad}}_{(t)}$ for any real $t>0$.
\end{lemma}

\begin{proof}
$\Gamma_S$ is the group of positive rationals supported on $S_f$
(Proposition~\ref{prop:splitting}) and $P(n)\subseteq S_f$, so $n\in\Gamma_S$
and $n$ is invertible in $\Z_S$. For \eqref{eq:factorsthrough}: by
Proposition~\ref{prop:splitting} the map $|\cdot|_S:C_S\to\R_{>0}$ is a
surjective homomorphism with kernel $\Knat$, so $C_S/\Knat\cong\R_{>0}$, and
$r\mapsto\mu_{r^{-1}\Knat}$ (normalized Haar measure of the fibre) is a monoid
isomorphism onto the level measures by the coset computation in the proof of
Theorem~\ref{thm:radmonoidal}(i).
\end{proof}

We therefore call these objects \emph{norm classes}; there is no ideal of which
they are the class.  The arithmetic is not in these objects;
Section~\ref{sec:radtrace} shows where it is.

\begin{definition}\label{def:normclass}
For $t\in\R$ define the \emph{norm class}
\begin{equation}\label{eq:normclass}
  \Gamma^{\mathrm{rad}}_{(t)}:=\text{the normalized Haar probability measure
   of }\{x\in C_S:|x|_S=e^{-t}\},
\end{equation}
a compactly supported Radon probability measure on $C_S$, and write
\begin{equation}\label{eq:normclassn}
  \Gamma_n^{\mathrm{rad}}:=\Gamma^{\mathrm{rad}}_{(\log n)}
   =\int_{\Knat}\delta_S(1/n)\cdot u\,du
\end{equation}
for its value at an integer $n\ge1$, $du$ the probability Haar measure of
$\Knat$; the coset $\delta_S(1/n)\Knat$ is exactly the level set
$\{|x|_S=1/n\}$. The family $\{\Gamma^{\mathrm{rad}}_{(t)}\}_{t\in\R}$ is the
\emph{radial family}.
\end{definition}

The averaging group must be $\Knat$, including the archimedean sign factor
$\{\pm1\}$: only then is the support of $\Gamma_n^{\mathrm{rad}}$ the whole
level set, which is what makes composition a one-line Haar computation, and
only then are both archimedean densities $\tau_\infty(\pm1/n)$ present.

\begin{remark}\label{rem:whatthecosetis}
The ideles $a\in\A_S^\times$ with $\varphi(a)=n$ and $|a|_S=1/n$ are exactly
$a=(u_\infty;(u_pn)_p)$ with $u_\infty\in\{\pm1\}$, $u_p\in\Z_p^\times$, and
their classes are exactly $\delta_S(1/n)\cdot u$ with $u\in\Knat$. So
$\Gamma_n^{\mathrm{rad}}$ is the unit average of the classes of all lifts of
$a_S(n)$ with the given modulus, and the coset depends only on that modulus.
The averaging kills the unit part of the lift; what remains is the level set.
By Lemma~\ref{lem:notideal} what remains does not remember $n$ as an object,
only $|a_S(n)|_S$.
\end{remark}

\begin{remark}\label{rem:multiplier}
The norm classes are \emph{multipliers} of $\DA$, not elements of it. Indeed:
\begin{enumerate}
\item[(i)] $\Gamma_n^{\mathrm{rad}}\notin\DA$, since \eqref{eq:normclass} is a
point mass in the archimedean coordinate, hence not a smooth function of
$x_\infty$;
\item[(ii)] convolution with $\Gamma_n^{\mathrm{rad}}$ maps $\DA$ continuously
into the radial subalgebra $\Drad:=\{F\in\DA:F\ \Knat\text{-invariant}\}$,
because $(\Gamma_n^{\mathrm{rad}}\star G)(x)=\int_{\Knat}
G(u^{-1}\delta_S(n)x)\,du$ is smooth in $x_\infty$ by differentiation under a
finite measure, $\Knat$-invariant, and compactly supported;
\item[(iii)] $\Gamma_n^{\mathrm{rad}}$ \emph{is} an element of the completion
$\mathcal N_A^{\mathrm{dist}}$ of Proposition~\ref{prop:completion}, which is
where these classes live; this is Theorem~\ref{thm:radincompletion}.
\end{enumerate}
This is the same type distinction that Proposition~\ref{prop:trace}(d) makes
for the diagonal. Note that $\Gamma_1^{\mathrm{rad}}$ is the Haar measure of
$\Knat$, the unit-averaging projector: the identity of $\Drad$ and \emph{not}
the identity of $\DA$.
\end{remark}

\begin{theorem}\label{thm:radmonoidal}
Let $m,n\ge1$ and $s,t\in\R$.
\begin{enumerate}
\item[(i)] \emph{(Monoidality.)}
$\Gamma_m^{\mathrm{rad}}\star\Gamma_n^{\mathrm{rad}}=\Gamma_{mn}^{\mathrm{rad}}$,
and $\Gamma_1^{\mathrm{rad}}=\Delta^{\mathrm{rad}}$ is the unit-averaging
projector, the identity of $\Drad$. Hence $n\mapsto\Gamma_n^{\mathrm{rad}}$ is
a strong monoidal functor from $(\N_{\ge1},\times)$ to the convolution monoid
of $\Knat$-invariant compactly supported probability measures on $C_S$.
\item[(ii)] \emph{(Dagger.)}
$(\Gamma^{\mathrm{rad}}_{(t)})^\dagger=\Gamma^{\mathrm{rad}}_{(-t)}$ and
$(\Gamma^{\mathrm{rad}}_{(s)}\star\Gamma^{\mathrm{rad}}_{(t)})^\dagger
=\Gamma^{\mathrm{rad}}_{(-s)}\star\Gamma^{\mathrm{rad}}_{(-t)}$. In particular
$(\Gamma_n^{\mathrm{rad}})^\dagger=\Gamma^{\mathrm{rad}}_{(-\log n)}$ is
\emph{not} in $\{\Gamma_m^{\mathrm{rad}}:m\in\N_{\ge1}\}$ for $n\ge2$: the
realized family is not $\dagger$-stable, and the smallest $\dagger$-stable
submonoid of the radial family containing all $\Gamma_n^{\mathrm{rad}}$ is
indexed by $\Q_{>0}^\times$. The only $\dagger$-fixed member of the family is
$\Gamma_1^{\mathrm{rad}}=\Delta^{\mathrm{rad}}$.
\item[(iii)] \emph{(Bidegree.)} With the Tate functionals of
Proposition~\ref{prop:tate},
\begin{equation}\label{eq:radbidegree}
  M(\Gamma_n^{\mathrm{rad}})=\bigl(n^{1/2},n^{-1/2}\bigr)
   =n^{-1/2}\cdot(n,1),
\end{equation}
multiplicative in $n$ and exchanged by $\dagger$.
\item[(iv)] \emph{(Independence of all choices.)} $\Gamma_n^{\mathrm{rad}}$ is
independent of the lift, of the gauge $\gamma\in\Gamma_S$, of the uniformizers
and of the chart; it is defined for every chart, including charts with
$P(n)\not\subseteq S_f$, and the smeared classes of
Definition~\ref{def:smeared} form a family over the colimit $\DA^\infty$.
\end{enumerate}
\end{theorem}

\begin{proof}
(i) Let $K:=\Knat$ and write $\mu_{aK}$ for the normalized Haar measure of a
coset. Then $\mu_{aK}\star\mu_{bK}$ is the image of $\mu_K\otimes\mu_K$ under
$(u,v)\mapsto abuv$; since $u,v$ are independent Haar variables on the compact
group $K$, the product $uv$ is Haar on $K$, so the image is $\mu_{abK}$. With
$a=\delta_S(1/m)$, $b=\delta_S(1/n)$ and
$\delta_S(1/m)\delta_S(1/n)=\delta_S(1/mn)$ --- immediate from
\eqref{eq:deltaS} --- this is the claim. For $n=1$ the coset is $K$ itself and
$\mu_K\star F$ is the $K$-average of $F$, the identity on $K$-invariant $F$ and
a projector in general. Strong monoidality is the statement that the structural
isomorphisms are equalities, and the coherences are those of convolution of
measures on the locally compact abelian group $C_S$.

(ii) The involution on measures is push-forward along $x\mapsto x^{-1}$
composed with complex conjugation; the level measures are real, and
$\{|x|_S=e^{-t}\}^{-1}=\{|x|_S=e^{t}\}$, with Haar measure carried to Haar
measure by the inversion automorphism. Anti-multiplicativity is commutativity
of $C_S$, and $(\Delta^{\mathrm{rad}})^\dagger=\Delta^{\mathrm{rad}}$ because
$\Knat$ is a group. Non-$\dagger$-stability:
$\Gamma^{\mathrm{rad}}_{(-\log n)}=\Gamma^{\mathrm{rad}}_{(\log m)}$ would
force $\log m=-\log n$, i.e. $mn=1$. The $\dagger$-stable submonoid generated
by the $\Gamma_n^{\mathrm{rad}}$ is the image of $\Q_{>0}^\times$ under
\eqref{eq:factorsthrough}, because the group generated by $\{\log n\}$ inside
$(\R,+)$ is $\log\Q_{>0}^\times$.

(iii) $\lambda_\mp(\mu)=\int|x|_S^{\mp1/2}d\mu$ and $|x|_S\equiv1/n$ on the
support of the probability measure $\Gamma_n^{\mathrm{rad}}$, so
$\lambda_\mp=n^{\pm1/2}$. Multiplicativity follows from (i), and the exchange
under $\dagger$ from $\lambda_\mp(F^\dagger)=\overline{\lambda_\pm(F)}$.

(iv) \emph{Lift and gauge.} By Remark~\ref{rem:whatthecosetis} the support of
$\Gamma_n^{\mathrm{rad}}$ is the set of classes of \emph{all} admissible lifts
and the measure is the Haar measure of that set, so no lift is preferred; and a
gauge change $a\mapsto\gamma a$ with $\gamma\in\Gamma_S$ does not move the
class at all. \emph{Charts.} $\Gamma^{\mathrm{rad}}_{(t)}$ is \emph{defined} by
the single equation $|x|_S=e^{-t}$, and by Proposition~\ref{prop:splitting} the
modulus is the modulus of the archimedean coordinate in every chart, with Haar
measure $dt\otimes(\text{probability})$ in every chart. Hence the defining
condition is chart-invariant and the family is defined in every chart by the
same equation, with no comparison map needed; the inflation used by the colimit
carries $\Gamma_h$ to $\Gamma_h$ for every profile $h$, because
$\Gamma_h=h\circ(-\log|\cdot|_S)$ is a function of the modulus. Finally
\eqref{eq:normclass} refers to $n$ only through the real number $1/n$, so it
makes sense in every chart.
\end{proof}

\begin{remark}\label{rem:whatisrealized}
Theorem~\ref{thm:radmonoidal} is true with $n$ replaced by any positive real,
because by Lemma~\ref{lem:notideal} the realization factors through
$(\R_{>0},\times)$. It therefore proves nothing arithmetic on its own. What is
realized as objects is the monoid $(\N_{\ge1},\times)$ embedded in the modulus
group, with composition, unit, dagger on the group completion, and
half-density; the objects forget factorization. For indices $n$ with at least two distinct prime divisors the object
$\Gamma_n^{\mathrm{rad}}$ is structurally indistinguishable from a generic
$\Gamma^{\mathrm{rad}}_{(t)}$. The divisor arrows of the divisor poset have no
image at all: morphisms between distinct external classes vanish in the
ambient category, and averaging creates none.
\end{remark}

\begin{definition}\label{def:smeared}
For $h\in C_c^\infty(\R)$ put
\begin{equation}\label{eq:smeared}
   \Gamma_h(x):=h(-\log|x|_S),\qquad x\in C_S .
\end{equation}
\end{definition}

\begin{lemma}\label{lem:smeared}
$\Gamma_h\in\Drad$; $\Gamma_h=\int_\R h(t)\Gamma^{\mathrm{rad}}_{(t)}\,dt$;
$\Gamma_h\star\Gamma_g=\Gamma_{h*g}$ with $*$ additive convolution on $\R$;
$\Gamma_h^\dagger=\Gamma_{h^\vee}$ with $h^\vee(t)=\overline{h(-t)}$;
$\lambda_\mp(\Gamma_h)=\int_\R h(s)e^{\pm s/2}ds$; and if $0\notin\supp h$ then
$\Gamma_h\in\Doff$.
\end{lemma}

\begin{proof}
Membership: $|x|_S=|x_\infty|$, so $\Gamma_h$ is a smooth function of
$x_\infty$ on $\R^\times$, is constant in the compact coordinates, and is
supported in $\{e^{-b}\le|x_\infty|\le e^{-a}\}\times\prod\Z_p^\times$ for
$\supp h\subseteq[a,b]$, a compact subset of $C_S$. The integral formula and
the formula for $\lambda_\mp$ are the disintegration $d^*w=dt\otimes du$ of
Proposition~\ref{prop:splitting} in the coordinate $t=\log|w|_S$:
$\lambda_\mp(\Gamma_h)=\int_\R h(-t)e^{\mp t/2}dt=\int_\R h(s)e^{\pm s/2}ds$.
Convolution:
\[
  (\Gamma_h\star\Gamma_g)(x)=\int h(-\log|y|_S)\,g(-\log|x|_S+\log|y|_S)\,d^*y
   =\int_\R h(-\sigma)\,g(-\log|x|_S+\sigma)\,d\sigma,
\]
which is $(h*g)(-\log|x|_S)$.
The dagger formula is the definition of the involution. Last assertion: $|x|_S$
near $1$ means $-\log|x|_S$ near $0$, where $h$ vanishes, so $\Gamma_h$
vanishes on a neighbourhood of the identity class.
\end{proof}

\section{The trace on the radial family}
\label{sec:radtrace}

The arithmetic of the radial family lies in its \emph{contact locus} with the
diagonal: the support and the atom masses of the finite-place part of the
trace, disintegrated along the family.  Since the point classes are not test
functions, the disintegration is taken through the smeared classes.  The
finite part turns out to be purely atomic and the archimedean part a
continuous density, and the two play different roles.

\begin{theorem}\label{thm:radtrace}
Let $h\in C_c^\infty(\R)$ with $0\notin\supp h$. Then $T(\Gamma_h)$ is an
absolutely convergent finite sum, with no regularization and no counterterm,
and
\begin{align}
  T(\Gamma_h)&=\sum_{p\in S_f}\sum_{k\ge1}(\log p)\,p^{-k/2}
    \bigl[h(k\log p)+h(-k\log p)\bigr]+\int_\R K(|t|)h(t)\,dt
     \label{eq:radtrace}\\
  &=\sum_{\substack{n\ge2,\ n=p^k\\ p\in S_f}}\Lambda(n)\,n^{-1/2}
    \bigl[h(\log n)+h(-\log n)\bigr]+\int_\R K(|t|)h(t)\,dt .
     \label{eq:radtraceprime}
\end{align}
Here
\begin{equation}\label{eq:lambdavol}
  \Lambda(p^k)=\log p=\vol_{d^*x}\bigl(p^k\Z_p^\times\bigr),
\end{equation}
the multiplicative Haar volume of the valuation shell, and $\Lambda(n)=0$ for
indices $n$ with at least two distinct prime divisors because no single place has a shell of modulus $1/n$. On the
colimit with the trace $T_\infty$ the restriction $p\in S_f$ disappears and the
sums run over all prime powers $n$ with $|\log n|$ at most the support radius
of $h$.
\end{theorem}

\begin{proof}
By Lemma~\ref{lem:smeared}, $\Gamma_h\in\Doff$, so
Proposition~\ref{prop:trace}(b) applies:
\[
  T(\Gamma_h)=\sum_{v\in S}\int_{\Q_v^\times}\Gamma_h(j_v(x))\,\tau_v(x)\,d^*x,
\]
absolutely convergent, with no $c$ term and no counterterm.

\emph{Finite places.} By Proposition~\ref{prop:slices},
$j_p(p^kw)=(p^{-k}\ \text{at}\ \infty;\ w\ \text{at}\ p;\ p^{-k}$ at every
other place of $S)$. This lies in $\ker\varphi$, its valuation being $0$ at
every place of $S_f$, so $|j_p(p^kw)|_S=p^{-k}$ and
\[
  a(p^kw):=\Gamma_h(j_p(p^kw))=h(-\log p^{-k})=h(k\log p),
\]
\emph{independent of $w$}. Two consequences. First, the unit packet of
\eqref{eq:localterms} vanishes identically: for $k=0$ we get $a(u)=h(0)=a(1)$
for every $u\in\Z_p^\times$, so $\int_{\Z_p^\times}(a(u)-a(1))/|1-u|_p\,du=0$.
This is a property of radial sections; it is pointwise, and needs no smallness
hypothesis. Second, the shell part is, by
Proposition~\ref{prop:trace}(a) ($\tau_p\equiv p^{-|k|/2}$, constant on the
shell) and \eqref{eq:shellvol},
\[
  \sum_{k\ne0}(\log p)\,p^{-|k|/2}h(k\log p)
   =\sum_{k\ge1}(\log p)p^{-k/2}\bigl[h(k\log p)+h(-k\log p)\bigr],
\]
a finite sum because $\supp h$ is compact.

\emph{The archimedean place.} $j_\infty(x)=(x$ at $\infty$; $1$ at every finite
place$)$, so $|j_\infty(x)|_S=|x|$ and, in the coordinate $x=\pm e^t$,
$b(t):=\Gamma_h(j_\infty(\pm e^t))=h(-t)$, the same for both signs. By
Proposition~\ref{prop:trace}(b) with the measure $dx/(2|x|)$,
\[
  \int_{\R^\times}\Gamma_h(j_\infty(x))\,\tau_\infty(x)\,d^*x_\infty
   =\int_0^\infty K(t)[b(t)+b(-t)]\,dt
   =\int_0^\infty K(t)[h(-t)+h(t)]\,dt,
\]

which is $\int_\R K(|t|)h(t)\,dt$.

\emph{Grouping.} Setting $n=p^k$ turns $(\log p)p^{-k/2}$ into
$\Lambda(n)n^{-1/2}$, which is \eqref{eq:radtraceprime}; and
\eqref{eq:lambdavol} is the identification of $\log p$ as the shell volume just
used. indices $n$ with at least two distinct prime divisors contributes nothing because $n=p^k$ has no solution when
$n$ has two prime divisors.

\emph{Colimit.} By Proposition~\ref{prop:trace}(c) the passage to $T_\infty$
adds, for each prime $\ell\notin S_f$ with $\log\ell$ inside the support
radius, exactly the term $(\log\ell)\sum_{k\ne0}\ell^{-|k|/2}
\Gamma_h(\delta_S(\ell^{-k}))$, and $\Gamma_h(\delta_S(\ell^{-k}))=h(k\log\ell)$
since $|\delta_S(\ell^{-k})|_S=\ell^{-k}$; this is the missing part of the
sum, and nothing else is added.
\end{proof}

\begin{theorem}[The representing measures]\label{thm:repmeasures}
There are a unique Radon measure $\omega_{\mathrm{fin}}$ and a unique
continuous density $\omega_\infty$ on $\R\smallsetminus\{0\}$ with
\[
  \sum_{p\in S_f}W_p(\Gamma_h)=\int h\,d\omega_{\mathrm{fin}},\qquad
  W_\infty(\Gamma_h)=\int h\,\omega_\infty,
\]
for all $h\in C_c^\infty(\R\smallsetminus\{0\})$, namely
\begin{equation}\label{eq:repmeasures}
  \omega_{\mathrm{fin}}=\sum_{n=p^k,\ p\in S_f}\Lambda(n)n^{-1/2}
    \bigl(\delta_{\log n}+\delta_{-\log n}\bigr),\qquad
  \omega_\infty=K(|t|)\,dt .
\end{equation}
Define the \emph{finite contact} of the norm class $\Gamma_n^{\mathrm{rad}}$,
$n\ge2$, as the mass of the atom of $\omega_{\mathrm{fin}}$ at $t=\log n$, and
the \emph{archimedean density of the radial family at the index $\log n$} as
the value of $\omega_\infty/dt$ there:
\begin{align}
  \cfin(\Gamma_n^{\mathrm{rad}})&:=\omega_{\mathrm{fin}}(\{\log n\})
     =\Lambda(n)\,n^{-1/2},\label{eq:radcontact}\\
  \omega_\infty(\log n)&:=\frac{\omega_\infty}{dt}\Big|_{t=\log n}
     =K(\log n)=\tfrac12\bigl[\tau_\infty(1/n)+\tau_\infty(-1/n)\bigr]
     =\frac{n^{3/2}}{n^2-1}=n^{-1/2}\sum_{j\ge0}n^{-2j}.\label{eq:archdensity}
\end{align}
Since $\omega_{\mathrm{fin}}$ and $\omega_\infty$ exist and are unique, the atom
mass and the density value at $\log n$ are invariants determined by $T$ and the
radial family alone.
\end{theorem}

\begin{proof}
Existence and the displayed forms are \eqref{eq:radtrace}, read as the statement
that the linear functionals $h\mapsto\sum_pW_p(\Gamma_h)$ and
$h\mapsto W_\infty(\Gamma_h)$ on $C_c^\infty(\R\smallsetminus\{0\})$ are given
by integration against the displayed measures. Both are positive and locally
finite --- the first is a locally finite sum of atoms because the prime powers
are discrete in $\R_{>0}$, the second has continuous density $K(|t|)$ away from
$0$ --- hence Radon. Uniqueness is the Riesz representation theorem: a Radon
measure on $\R\smallsetminus\{0\}$ is determined by its integrals against
$C_c^\infty$ tests. For \eqref{eq:archdensity},
$\tau_\infty(x)=|x|^{1/2}/|1-x|$ gives
\[
  \tau_\infty(1/n)=\frac{n^{-1/2}}{1-1/n}=\frac{n^{1/2}}{n-1},\qquad
  \tau_\infty(-1/n)=\frac{n^{-1/2}}{1+1/n}=\frac{n^{1/2}}{n+1},
\]
whose half-sum is $\tfrac12n^{1/2}\cdot\frac{2n}{n^2-1}=n^{3/2}/(n^2-1)$; and
$K(\log n)=n^{-1/2}/(1-n^{-2})=n^{3/2}/(n^2-1)$, whose geometric expansion is
$n^{-1/2}\sum_{j\ge0}n^{-2j}$. The half-sum form, with the factor $\tfrac12$
in front, comes from the archimedean measure $dx/(2|x|)$.
\end{proof}

\begin{remark}\label{rem:typewarning}
The quantity $T(\Gamma_n^{\mathrm{rad}})$ is \emph{undefined}: $T$ is a
functional on $\DA$ and $\Gamma_n^{\mathrm{rad}}\notin\DA$
(Remark~\ref{rem:multiplier}(i)), and we do not extend $T$ to multipliers.
Moreover, the naive place-wise extension
$W_p(\mu):=\int(\mu\circ j_p)\tau_p\,d^*x$ returns $0$, not
$\Lambda(n)n^{-1/2}$: for $|S_f|\ge2$ the set
$\Iml\,j_p\cap\{|x|_S=1/n\}$ is, in each coordinate $\Z_r^\times$ with
$r\ne p$, the single point $n^{-1}$, hence a $\Gamma_n^{\mathrm{rad}}$-null
set; this is the computation of Proposition~\ref{prop:noslice}. So
$\Lambda(n)n^{-1/2}$ is obtained \emph{only} as the atom mass of the unique
representing measure of the radial disintegration, and every number attached to
$\Gamma_n^{\mathrm{rad}}$ in this paper is to be read that way.
\end{remark}

\begin{theorem}[The contact locus]\label{thm:contactlocus}
Let $\omega_{\mathrm{fin}}$ be as in Theorem~\ref{thm:repmeasures}. Then:
\begin{enumerate}
\item[(i)] $\omega_{\mathrm{fin}}$ is purely atomic and
\[
  \supp\omega_{\mathrm{fin}}=\{\pm k\log p:p\in S_f,\ k\ge1\}
   \subset\R\smallsetminus\{0\},
\]
the condition $p\in S_f$ being dropped in the colimit;
\item[(ii)] the atom masses are the shell volumes against the half-densities:
$\omega_{\mathrm{fin}}(\{\pm k\log p\})=(\log p)p^{-k/2}$, with $\log p$ the
multiplicative Haar volume of $p^k\Z_p^\times$, independent of $k$;
\item[(iii)] consequently, for every $t>0$,
\[
  \omega_{\mathrm{fin}}(\{t\})\ne0\iff e^t=p^k\quad(p\in S_f,\ k\ge1),
\]
with all primes allowed in the colimit. At such a point $\omega_{\mathrm{fin}}(\{t\})=\Lambda(e^t)e^{-t/2}$. In particular
$\Lambda$, the set of prime powers, and the vanishing of $\Lambda$ at integers with at least two distinct prime divisors are read off the contact locus of the one-parameter family
$\{\Gamma^{\mathrm{rad}}_{(t)}\}$, whose members by Lemma~\ref{lem:notideal}
carry no arithmetic information themselves. For $t=\log n$ with $n\ge2$ supported on the chart, there is no finite contact exactly when $n$ has at least two distinct prime divisors.
\end{enumerate}
\end{theorem}

\begin{proof}
All three parts are \eqref{eq:radtrace} read as an identity of measures. The
finite-place functional is
$h\mapsto\sum_{p\in S_f}\sum_{k\ge1}(\log p)p^{-k/2}[h(k\log p)+h(-k\log p)]$,
which is integration against
$\sum_{p,k}(\log p)p^{-k/2}(\delta_{k\log p}+\delta_{-k\log p})$; the
coefficients are strictly positive, so the support is exactly the index set and
the measure is purely atomic. Local finiteness is the discreteness of
$\{k\log p\}$ in $\R_{>0}$: in a compact window $[a,b]$ with $a>0$, only
finitely many pairs $(p,k)$ satisfy $a\le k\log p\le b$, since
$k\le b/\log2$ and $p\le e^b$. The shell volume is \eqref{eq:shellvol} and its
$k$-independence is translation invariance of $d^*x$ on $\Q_p^\times$. For
(iii), $e^t=p^k$ for some prime $p$ and $k\ge1$ is exactly the condition that
$t$ lies in the support, and then $\Lambda(p^k)(p^k)^{-1/2}=(\log p)p^{-k/2}$.
\end{proof}

\begin{remark}\label{rem:typesdiffer}
The finite contact and the archimedean density are of different type and are
not added. $\cfin(\Gamma_n^{\mathrm{rad}})$ is the mass of an \emph{atom}: a
number attached to the single norm class, zero unless $n$ is a prime power.
The quantity $\omega_\infty(\log n)$ is the value of a \emph{continuous
density} at a point of its argument: it is an invariant of $\omega_\infty$ and
of the real number $\log n$, it is defined for every real $t$, and nothing
about the object $\Gamma_n^{\mathrm{rad}}$ enters it --- the family has no
archimedean atom at all.  It is therefore not a contact, the two are not
added, and $\omega_\infty(\log n)$ is no part of $\Lambda(n)$.  The archimedean
contribution to the radial disintegration is the continuous density
$K(|t|)\,dt$.
\end{remark}

\begin{remark}\label{rem:lambdaderived}
Formula \eqref{eq:lambdavol} exhibits $\Lambda(p^k)=\log p$ as the
multiplicative Haar volume of a valuation shell, and the
vanishing of $\Lambda$ at indices $n$ with at least two distinct prime divisors as the statement that no single place
has a shell of modulus $1/n$. Nothing in the derivation selects an arithmetic
family from outside or uses the von Mangoldt function as an input: the only
ingredients are the local terms \eqref{eq:localterms}, the exact orbital values
of Proposition~\ref{prop:trace}(a), and the shell volume \eqref{eq:shellvol}.
That $\Lambda$ is a volume depends on the normalization \eqref{eq:shellvol};
Section~\ref{sec:normalization} shows that this normalization is, up to a
global scale, the only one compatible with a common Tate residue at all
places.
\end{remark}

\begin{theorem}\label{thm:radincompletion}
Let $n\ge2$. Choose $h_\varepsilon\in C_c^\infty(\R)$, $h_\varepsilon\ge0$,
$\int h_\varepsilon=1$, $\supp h_\varepsilon\subseteq[\log n-\varepsilon,
\log n+\varepsilon]$ with $\varepsilon<\tfrac12\log2$, and put
$F_\varepsilon:=\Gamma_{h_\varepsilon}\in\DA$. Then:
\begin{enumerate}
\item[(i)] $\Gamma_n^{\mathrm{rad}}\star G^\dagger\in\DA$ for every
$G\in\DA$, so $\ell_n(G):=T(\Gamma_n^{\mathrm{rad}}\star G^\dagger)
=I_A(\Gamma_n^{\mathrm{rad}},G)$ is defined;
$\widehat{\Gamma_n^{\mathrm{rad}}}$ is supported on the principal component
and equals $n^{i\tau}$ there, of modulus $1$, whence
$|\ell_n(G)|\le\mathcal N(G)$ and $\ell_n\in\DA'$;
\item[(ii)] $I_A(F_\varepsilon,\cdot)\to\ell_n$ in the strong dual $\DA'$,
i.e. uniformly on every bounded subset of $\DA$;
\item[(iii)] $M(F_\varepsilon)\to(n^{1/2},n^{-1/2})
=M(\Gamma_n^{\mathrm{rad}})$;
\item[(iv)] consequently $J(\Gamma_n^{\mathrm{rad}}):=(\ell_n,
(n^{1/2},n^{-1/2}))\in\mathcal N_A^{\mathrm{dist}}$;
\item[(v)] in the canonical splitting of Proposition~\ref{prop:completion},
\[
  J(\Gamma_n^{\mathrm{rad}})=\underbrace{(\ell_n,0)}
    _{\in\mathcal N_A^{\mathrm{dist,prim}}}
   +\underbrace{\sigma(n^{1/2},n^{-1/2})}_{\in\sigma(B_T)},
\]
so the ruling component is nonzero for every $n$ and
$\Gamma_n^{\mathrm{rad}}$ is not a primitive class.
\end{enumerate}
\end{theorem}

\begin{proof}
(i) $\Gamma_n^{\mathrm{rad}}$ is a multiplier of $\DA$
(Remark~\ref{rem:multiplier}(ii)), so
$\Gamma_n^{\mathrm{rad}}\star G^\dagger\in\Drad\subset\DA$ and $T$ of it is
defined. For the Fourier transform: $\Gamma_n^{\mathrm{rad}}$ is the Haar
measure of the coset $\delta_S(1/n)\Knat$, so for a character $\chi$ of $C_S$,
\[
  \widehat{\Gamma_n^{\mathrm{rad}}}(\chi)
   =\int\overline\chi\,d\Gamma_n^{\mathrm{rad}}
   =\overline{\chi(\delta_S(1/n))}\int_{\Knat}\overline{\chi(u)}\,du,
\]
which vanishes unless $\chi|_{\Knat}=1$, that is, unless $\chi$ lies on the
principal component, where $\chi(w)=|w|_S^{i\tau}$; and then it equals
$\overline{(1/n)^{i\tau}}=n^{i\tau}$. The bound is
Proposition~\ref{prop:completion} with $\sup_\chi|\widehat F|=1$, and the
formula for $I_A$ is \eqref{eq:planch} applied to
$\Gamma_n^{\mathrm{rad}}\star G^\dagger$, legitimate because that element lies
in $\DA$.

(ii) $F_\varepsilon$ is $\Knat$-invariant, so $\widehat F_\varepsilon$ is also
supported on the principal component, where
$\widehat F_\varepsilon(\tau)=\int_\R h_\varepsilon(s)e^{i\tau s}\,ds$. Write
$\Delta_\varepsilon(\tau):=\widehat F_\varepsilon(\tau)-n^{i\tau}
=\int h_\varepsilon(s)(e^{i\tau s}-e^{i\tau\log n})\,ds$. Then
$|\Delta_\varepsilon(\tau)|\le2$ always, and
$|\Delta_\varepsilon(\tau)|\le|\tau|\varepsilon$ since
$|e^{i\theta}-e^{i\theta'}|\le|\theta-\theta'|$ and $|s-\log n|\le\varepsilon$
on $\supp h_\varepsilon$. Let $B\subseteq\DA$ be bounded and take $d_N,C''_B$
as in Proposition~\ref{prop:completion}. Then for $R>0$
\[
  \sup_{G\in B}\bigl|I_A(F_\varepsilon,G)-\ell_n(G)\bigr|
   \le\sup_{G\in B}\int_{\mathfrak c_0}|\Delta_\varepsilon|\,|\widehat G|\,
     |\widehat T|\,d\tau
   \le R\varepsilon\,\mathcal N_B+\eta(R),
\]
with $\mathcal N_B:=\sup_{G\in B}\mathcal N(G)<\infty$ and
\[
  \eta(R):=2C''_B\int_{|\tau|>R}\tau^{-2}\bigl(A_B+B_S\log(2+|\tau|)\bigr)
    \frac{d\tau}{2\pi},
\]
which is finite and tends to $0$ as $R\to\infty$, uniformly on $B$. Given
$\delta>0$ pick $R$ with $\eta(R)<\delta/2$ and then
$\varepsilon<\delta/(2R\mathcal N_B)$. This is uniform convergence on $B$.

(iii) By Lemma~\ref{lem:smeared},
$\lambda_\mp(F_\varepsilon)=\int h_\varepsilon(s)e^{\pm s/2}ds\to n^{\pm1/2}$.

(iv) $J(F)=(I_A(F,\cdot),M(F))$; combine (ii) and (iii). Each
$J(F_\varepsilon)$ lies in $J(\DA)$, so the limit lies in its closure.

(v) is the canonical splitting of Proposition~\ref{prop:completion} applied to
$x=(\ell_n,M(\Gamma_n^{\mathrm{rad}}))$; the ruling component does not vanish
because $M(\Gamma_n^{\mathrm{rad}})=(n^{1/2},n^{-1/2})\ne(0,0)$.
\end{proof}

\section{The Frobenius family}\label{sec:frob}

We now construct objects that are not radial and that do carry the
arithmetic: the Haar measures of the valuation shells of the local dilations.
They compose with a genuine unit, their supports remember $n$ and $P(n)$, and
in Section~\ref{sec:frobcontact} we compute their contact to be
$\Lambda(n)n^{-1/2}$, which vanishes at indices $n$ with at least two distinct prime divisors for a reason visible in
the support.

Throughout, $S$ is a chart and $n\ge1$ with $P(n)\subseteq S_f$.

\begin{definition}\label{def:localfrob}
For $p\in S_f$ and $k\ge1$ let
\[
  \Z_p^{(S)}:=j_p(\Z_p^\times)
   =\{(1\ \text{at}\ \infty;\ w\ \text{at}\ p;\ 1\ \text{elsewhere}):
      w\in\Z_p^\times\},
\]
a compact subgroup of $C_S$, and define the \emph{local Frobenius}
\begin{equation}\label{eq:localfrob}
  \Frob_{p^k}:=\int_{\Z_p^\times}E_{j_p(p^ku)}\,du
   =\text{the normalized Haar probability measure of }
     j_p(p^k\Z_p^\times),
\end{equation}
$du$ the probability Haar measure of $\Z_p^\times$. Thus $\Frob_{p^k}$ is the
Haar measure of the valuation shell $p^k\Z_p^\times$, transported to $C_S$ by
the slice map. For a finite $F\subseteq S_f$ put
\[
  K_F:=\prod_{p\in F}\Z_p^{(S)}\subseteq C_S,\qquad
  \Delta^{(F)}:=\text{normalized Haar measure of }K_F,
\]
and for $n\ge1$ write $\Delta^{(n)}:=\Delta^{(P(n))}$. Thus
$\Delta^{(\emptyset)}=\delta_1=\Delta$, the point mass at the identity class,
i.e. the diagonal.
\end{definition}

\begin{definition}\label{def:frob}
For $n=\prod_p p^{k_p}\ge1$,
\begin{equation}\label{eq:frob}
  \Frob_n:=\mathop{\star}_{p\mid n}\Frob_{p^{k_p}},\qquad \Frob_1:=\Delta .
\end{equation}
Equivalently, with the integral idele
$\iota(n):=(1\ \text{at}\ \infty;\ p^{v_p(n)}$ at every finite place of $S)$
and $c_n$ the class of $\iota(n)$,
\begin{equation}\label{eq:frobcoset}
  \Frob_n=\int_{\prod_{p\mid n}\Z_p^\times}E_{c_n\cdot u}\,du
   =\text{the normalized Haar measure of the coset } c_nK_{P(n)} .
\end{equation}
\end{definition}

\begin{lemma}\label{lem:frobsupport}
$c_n=\prod_{p\mid n}j_p(p^{v_p(n)})$, and
\begin{equation}\label{eq:frobsupp}
  \supp\Frob_n=c_nK_{P(n)}
  =\bigl\{(1/n\ \text{at}\ \infty;\ u_p\in\Z_p^\times\ \text{at each}\ p\mid n;
    \ 1/n\ \text{at each}\ r\nmid n)\bigr\}.
\end{equation}
In particular $|x|_S=1/n$ on the support, and $\Frob_n$ is supported in the
positive archimedean component.
\end{lemma}

\begin{proof}
$\varphi(\iota(n))=\prod_{p\in S_f}p^{v_p(n)}=n\in\Gamma_S$, so by
Proposition~\ref{prop:splitting} the class of $\iota(n)$ is
$\iota(n)/n=(1/n$ at $\infty$; $p^{v_p(n)}/n$ at $p)$. At $p\mid n$ the
coordinate $p^{v_p(n)}/n=(n/p^{v_p(n)})^{-1}$ is a unit of $\Z_p$; at $r\nmid
n$ it is $1/n\in\Z_r^\times$. Separately,
$\prod_{p\mid n}j_p(p^{v_p(n)})$ has archimedean coordinate
$\prod_{p\mid n}p^{-v_p(n)}=1/n$, $r$-coordinate
$\prod_{p\mid n,\,p\ne r}p^{-v_p(n)}$ for $r\mid n$, a unit of $\Z_r$, and
$1/n$ for $r\nmid n$. The two agree at $\infty$ and at every $r\nmid n$, and at
$p\mid n$ they differ by a unit of $\Z_p$, which the coset $c_nK_{P(n)}$
absorbs. Hence \eqref{eq:frobcoset} and \eqref{eq:frobsupp} describe the same
coset, and $|x|_S=|x_\infty|=1/n$ on it.
\end{proof}

\begin{theorem}\label{thm:frobmonoidal}
For all $m,n\ge1$:
\begin{enumerate}
\item[(i)] $\Frob_m\star\Frob_n=\Frob_{mn}$, an \emph{equality} of measures;
\item[(ii)] $\Frob_1=\Delta=E_1$, the diagonal itself --- a point mass, not a
projector;
\item[(iii)] $n\mapsto\Frob_n$ is a strong monoidal functor from
$(\N_{\ge1},\times)$ to the convolution monoid $M_c(C_S)$ of compactly
supported Radon measures, preserving the unit strictly;
\item[(iv)] it is injective, and $\supp\Frob_n$ determines both $n$ and
$P(n)$.
\end{enumerate}
\end{theorem}

\begin{proof}
(i) Write $\mu_{cK}$ for the normalized Haar measure of a coset of a compact
subgroup $K\le C_S$. If $K,K'$ are compact subgroups then
$\mu_{cK}\star\mu_{c'K'}=\mu_{cc'KK'}$: the convolution is the image of
$\mu_K\otimes\mu_{K'}$ under $(u,v)\mapsto cc'uv$, and the image of
$\mu_K\otimes\mu_{K'}$ under multiplication is a $KK'$-invariant probability
measure on the compact group $KK'$, hence $\mu_{KK'}$ by uniqueness of Haar
measure. Now $K_{P(m)}K_{P(n)}=K_{P(m)\cup P(n)}=K_{P(mn)}$ and
\[
  c_mc_n=\prod_p j_p(p^{v_p(m)})j_p(p^{v_p(n)})
   =\prod_p j_p(p^{v_p(m)+v_p(n)})=c_{mn},
\]
using that each $j_p$ is a homomorphism and that distinct $j_p$ commute. Hence
$\Frob_m\star\Frob_n=\mu_{c_{mn}K_{P(mn)}}=\Frob_{mn}$.

(ii) For $n=1$ the product \eqref{eq:frob} is empty, $c_1=1$ and
$K_{P(1)}=K_\emptyset=\{1\}$, so $\Frob_1=\delta_1$, the point mass at the
identity class, which is $E_1=\Delta$; and $\delta_1$ is the unit of
$(M_c(C_S),\star)$.

(iii) is (i) and (ii); the coherences are those of convolution of measures on
the locally compact abelian group $C_S$.

(iv) By \eqref{eq:frobsupp}, $\sup_{x\in\supp\Frob_n}|x|_S=1/n$ recovers $n$,
and the set of finite places at which the support projects \emph{onto} the
whole of $\Z_p^\times$, rather than to a single point, is exactly $P(n)$.
\end{proof}

\begin{remark}\label{rem:nointerp}
By Lemma~\ref{lem:notideal} the radial realization factors through
$(\R_{>0},\times)$ and forgets factorization. The Frobenius realization does
not: by Theorem~\ref{thm:frobmonoidal}(iv) the object determines the pair
$(n,P(n))$, and it is defined only for integer indices, because the datum
$K_{P(n)}$ --- which finite coordinates are averaged --- has no meaning for a
real parameter. Concretely there is \emph{no} one-parameter family
$\{\Frob_{(t)}\}_{t\in\R}$ of compactly supported measures with
$\Frob_{(\log n)}=\Frob_n$ continuous in the weak-$*$ topology: the support of
$\Frob_n$ has the discrete invariant $|P(n)|$, the number of coordinates on
which the support is a full $\Z_p^\times$, and $|P(n)|$ is not locally constant
along any ordering of the integers --- $|P(2)|=|P(4)|=1$ but $|P(6)|=2$ ---
while any weak-$*$ limit of measures supported in $\{x_r=1/n\}$ for $r\nmid n$
is again supported in such a set.
\end{remark}

\begin{proposition}\label{prop:avgtorad}
$\Frob_n\star\Gamma_1^{\mathrm{rad}}=\Gamma_n^{\mathrm{rad}}$, where
$\Gamma_1^{\mathrm{rad}}=\Delta^{\mathrm{rad}}$ is the Haar measure of
$\Knat$. So the norm class is the image of the Frobenius under the total
unit-averaging projector, and the loss of factorization in
Lemma~\ref{lem:notideal} is the kernel of that projector.
\end{proposition}

\begin{proof}
$\mu_{c_nK_{P(n)}}\star\mu_{\Knat}=\mu_{c_nK_{P(n)}\Knat}=\mu_{c_n\Knat}$ by
the coset formula in the proof of Theorem~\ref{thm:frobmonoidal}(i), since
$K_{P(n)}\subseteq\Knat$. And $c_n\Knat=\{|x|_S=1/n\}$ because $|c_n|_S=1/n$
and $\Knat=\ker(|\cdot|_S)$.
\end{proof}

\subsection{The involution and the inverse semigroup}

\begin{theorem}\label{thm:frobdagger}
For $m,n\ge1$:
\begin{enumerate}
\item[(i)] $\Frob_n^\dagger=\mu_{c_n^{-1}K_{P(n)}}$, supported on
$\{|x|_S=n\}$;
\item[(ii)] $\Frob_n\star\Frob_n^\dagger=\Delta^{(n)}$, the Haar measure of
$K_{P(n)}$, an idempotent:
$\Delta^{(n)}\star\Delta^{(n)}=\Delta^{(n)}$;
\item[(iii)] $M(\Frob_n)=(n^{1/2},n^{-1/2})$ and
$M(\Delta^{(n)})=(1,1)=M(\Delta)$;
\item[(iv)] $\Delta^{(n)}\ne\Delta$ for every $n>1$.
\end{enumerate}
\end{theorem}

\begin{proof}
(i) The involution is push-forward along inversion composed with conjugation;
$\Frob_n$ is a real positive measure, and $(c_nK_{P(n)})^{-1}
=c_n^{-1}K_{P(n)}$ because $K_{P(n)}$ is a group; and $|c_n^{-1}|_S=n$.

(ii) By the coset formula,
$\mu_{c_nK}\star\mu_{c_n^{-1}K}=\mu_{KK}=\mu_K$ with $K=K_{P(n)}$;
idempotence is the same formula with $c=1$.

(iii) $|x|_S\equiv1/n$ on $\supp\Frob_n$, a probability measure, so
$\lambda_-=n^{1/2}$, $\lambda_+=n^{-1/2}$; on $K_{P(n)}$ the modulus is $1$.

(iv) $\Delta^{(n)}$ is the Haar measure of the compact group $K_{P(n)}$, which
is infinite for $P(n)\ne\emptyset$ since it contains an infinite $\Z_p^\times$;
the Haar measure of an infinite compact group has no atoms, whereas
$\Delta=\delta_1$ is an atom.
\end{proof}

\begin{remark}\label{rem:weildegree}
Weil's degree relation for the Frobenius correspondence on a curve over $\F_q$
reads $\Gamma_F\star\Gamma_F^t=q\cdot\Delta$ \cite{Weil1948}.
Theorem~\ref{thm:frobdagger}(ii) is its averaged, half-density normalized
analogue: the factor $q$ is absent because the canonical half-density
$|\cdot|_S^{-1/2}$ has already been applied --- the bidegree
$(n^{1/2},n^{-1/2})$ has product $1$, which is where the degree went --- and
$\Delta$ is replaced by the idempotent $\Delta^{(n)}$, which by (iv) is
not $\Delta$.  Part (iv) records how the two relations differ; the analogy is
not used in any proof.
\end{remark}

\begin{theorem}\label{thm:invsemigroup}
For $r\in\Q_{>0}^\times$ with $P(r)\subseteq S_f$ and a finite $F\subseteq S_f$
with $P(r)\subseteq F$, put
\begin{equation}\label{eq:gradedfrob}
  c_r:=\prod_p j_p(p^{v_p(r)}),\qquad \Frob_{r,F}:=\mu_{c_rK_F},
\end{equation}
so that $\Frob_n=\Frob_{n,P(n)}$. Then:
\begin{enumerate}
\item[(i)] $\Frob_{r,F}\star\Frob_{s,G}=\Frob_{rs,\,F\cup G}$ and
$\Frob_{r,F}^\dagger=\Frob_{r^{-1},F}$;
\item[(ii)] $\mathfrak F_S:=\{\Frob_{r,F}\}$ is a commutative \emph{inverse
semigroup} under $\star$, with inverse $\Frob_{r,F}^\dagger$, whose idempotents
are exactly the $\Delta^{(F)}$, forming a semilattice isomorphic to the finite
subsets of $S_f$ under union;
\item[(iii)] $\mathfrak F_S$ is $\dagger$-closed, contains every $\Frob_n$ and
every $\Frob_n^\dagger$, and is the smallest such set closed under $\star$ and
$\dagger$;
\item[(iv)] the assignment $a/b\mapsto\Frob_a\star\Frob_b^\dagger$ does
\emph{not} descend to $\Q_{>0}^\times$.
\end{enumerate}
\end{theorem}

\begin{proof}
(i) is the coset formula together with $c_rc_s=c_{rs}$ and
$K_FK_G=K_{F\cup G}$, exactly as in Theorem~\ref{thm:frobmonoidal}(i).

(ii) $\mathfrak F_S$ is closed under $\star$ by (i), and commutative because
$C_S$ is abelian. Each element is regular:
\[
  x\star x^\dagger\star x=\Frob_{r,F}\star\Frob_{r^{-1},F}\star\Frob_{r,F}
   =\Delta^{(F)}\star\Frob_{r,F}=\Frob_{r,F},
\]
the last step because $K_FK_F=K_F$ and $c_rK_FK_F=c_rK_F$; and symmetrically
$x^\dagger\star x\star x^\dagger=x^\dagger$. A commutative regular semigroup is
an inverse semigroup with unique inverses \cite{Lawson}. Idempotents:
$\Frob_{r,F}\star\Frob_{r,F}=\Frob_{r^2,F}$, which equals $\Frob_{r,F}$ iff
$c_{r^2}K_F=c_rK_F$ iff $c_r\in K_F$ iff $|c_r|_S=1$ and $v_p(r)=0$ for
$p\notin F$, i.e. iff $r=1$. So the idempotents are the
$\Frob_{1,F}=\Delta^{(F)}$, and $\Delta^{(F)}\star\Delta^{(G)}
=\Delta^{(F\cup G)}$ gives the semilattice.

(iii) $\Frob_n=\Frob_{n,P(n)}$ and $\Frob_n^\dagger=\Frob_{n^{-1},P(n)}$, and
by (i) any $\star$/$\dagger$ word in these lies in $\mathfrak F_S$.
Conversely every element of $\mathfrak F_S$ is such a word: for a finite
$F=\{p_1,\dots,p_k\}\subseteq S_f$ we have
$\Delta^{(F)}=\Frob_{p_1}\star\Frob_{p_1}^\dagger\star\cdots\star
\Frob_{p_k}\star\Frob_{p_k}^\dagger$ by Theorem~\ref{thm:frobdagger}(ii) and
(i); and for $r=a/b$ in lowest terms with $P(r)\subseteq F$,
\[
  \Frob_{r,F}=\Frob_a\star\Frob_b^\dagger\star\Delta^{(F)},
\]
since the right-hand side is $\mu_{c_ac_b^{-1}K_{P(a)\cup P(b)\cup F}}
=\mu_{c_rK_F}$ by (i) and $P(a)\cup P(b)=P(r)\subseteq F$. So
$\mathfrak F_S$ is the smallest $\star$- and $\dagger$-closed set containing
the $\Frob_n$.

(iv) We exhibit the failure. We have $2/1=6/3$, but
\[
  \Frob_2\star\Frob_1^\dagger=\mu_{j_2(2)K_{\{2\}}},\qquad
  \Frob_6\star\Frob_3^\dagger=\mu_{c_6c_3^{-1}K_{\{2,3\}}}
   =\mu_{j_2(2)K_{\{2,3\}}},
\]
and these differ: the first is supported on a coset of $\Z_2^{(S)}$, the second
on a coset of $\Z_2^{(S)}\times\Z_3^{(S)}$ --- the second projects onto all of
$\Z_3^\times$ at the place $3$, the first onto the single point $1/2$. Both are
probability measures on those cosets and a measure determines its support, so
they are distinct. The correct index set is therefore the one of
\eqref{eq:gradedfrob}: $\Q_{>0}^\times$ \emph{graded by the semilattice of
finite sets of primes}.  The idempotents $\Delta^{(F)}$ are genuine extra
data, and this is what makes $\mathfrak F_S$ an inverse semigroup rather than
a group.
\end{proof}

\section{The contact of the Frobenius family}\label{sec:frobcontact}

The smearing is chosen so that the finite contact is exact at every level, with
no limit taken at the finite places.

\begin{definition}\label{def:smearedfrob}
Let $n\ge2$ and $P:=P(n)$. Fix $h\in C_c^\infty(\R_{>0})$ with
\begin{equation}\label{eq:hsupport}
  \supp h\subset\bigl(n^{-1}e^{-\varepsilon},\ n^{-1}e^{\varepsilon}\bigr),
\end{equation}
and an integer $N\ge1$. Define $\Frob_n^{[h,N]}:C_S\to\C$ by
\begin{equation}\label{eq:smearedfrob}
  \Frob_n^{[h,N]}(x):=h(x_\infty)\cdot\prod_{r\in S_f\smallsetminus P}
   \mathbf 1_{\,n^{-1}(1+r^N\Z_r)}(x_r),
\end{equation}
with no condition at the places $p\in P$ and with the indicators
unnormalized. Write $\Frob_n^{[h,N],\flat}:=\bigl(\prod_{r\notin P}
\vol(1+r^N\Z_r)\bigr)^{-1}\Frob_n^{[h,N]}$ for the normalized version.
\end{definition}

\begin{lemma}\label{lem:smearedfrobinDA}
$\Frob_n^{[h,N]}\in\DA$, and if $\varepsilon<\log n$ then
$\Frob_n^{[h,N]}\in\Doff$. Moreover, if $h\ge0$ is normalized by $\int_{\R_{>0}}h\,d^*x_\infty=1$
(the measure $d^*x_\infty=dx/(2|x|)$ of \S\ref{sec:prelim}), then
$\Frob_n^{[h,N],\flat}\to\Frob_n$ weak-$*$ as $\varepsilon\to0$ and
$N\to\infty$; without this normalization the limit is
$(\int h\,d^*x_\infty)\cdot\Frob_n$ whenever that integral converges as
$\varepsilon\to0$.
\end{lemma}

\begin{proof}
The function is smooth in $x_\infty$, being a $C_c^\infty$ function of the
archimedean coordinate; it is locally constant in the compact coordinates,
being a product of indicators of cosets of the open subgroups $1+r^N\Z_r$; and
it is supported in a compact subset of $C_S$, compact in $x_\infty$ by
\eqref{eq:hsupport} and compact in the remaining coordinates automatically. For
$\Doff$: the identity class has $x_\infty=1$, and $1\notin\supp h$ because
$n^{-1}e^\varepsilon<1$ when $\varepsilon<\log n$; since $\supp h$ is closed in
$\R^\times$ the function vanishes on a neighbourhood of the identity class.
Weak-$*$ convergence: with $\int h\,d^*x_\infty=1$ and $h\ge0$, the
normalized function is a probability density (for the Haar measure of $C_S$,
which is $d^*x_\infty$ times the probability Haar measures of the
$\Z_r^\times$) on a
neighbourhood basis of the coset $c_nK_P$ intersected with the level sets
$\{x_r\in n^{-1}(1+r^N\Z_r)\}$, and these shrink to $c_nK_P$.
\end{proof}

\begin{theorem}\label{thm:frobcontact}
Let $n\ge2$, and let $h$ and $N$ be as in Definition~\ref{def:smearedfrob}
with $\varepsilon$ small enough that $\supp h$ contains no point of
$\{q^{-m}:q\in S_f,\ m\in\Z\}$ other than $1/n$ itself. Then
\begin{equation}\label{eq:frobcontacteq}
  \sum_{q\in S_f}W_q\bigl(\Frob_n^{[h,N]}\bigr)
   =\Lambda(n)\,n^{-1/2}\,h(1/n)
\end{equation}
for every $N\ge1$, with no regularization, counterterm or limit at the
finite places; and in particular the left-hand side is $0$ for
every index $n$ with at least two distinct prime divisors. The archimedean term is
\begin{equation}\label{eq:frobarch}
  W_\infty\bigl(\Frob_n^{[h,N]}\bigr)
   =\Bigl(\int_0^\infty K_s(t)\,h(e^{-t})\,dt\Bigr)\cdot
     \prod_{r\in S_f\smallsetminus P(n)}\mathbf 1\bigl[\,r^N\mid n-1\,\bigr];
\end{equation}
in particular it vanishes for every sufficiently large $N$ as soon as
$S_f\supsetneq P(n)$.
\end{theorem}

\begin{proof}
Put $F:=\Frob_n^{[h,N]}\in\Doff$ by Lemma~\ref{lem:smearedfrobinDA}, so
Proposition~\ref{prop:trace}(b) applies and
$T(F)=\sum_{v\in S}\int F(j_v(x))\tau_v(x)\,d^*x$, absolutely convergent, with
no $c$ term. Write $a_q:=F\circ j_q$.

\emph{Shells at $q\in S_f$, $m\ne0$.} By Proposition~\ref{prop:slices} the
archimedean coordinate of $j_q(q^mw)$ is $q^{-m}$, so
$a_q(q^mw)=h(q^{-m})\cdot(\text{indicators})$. By hypothesis $h(q^{-m})=0$
unless $q^{-m}=1/n$, i.e. unless $n=q^m$, so $m\ge1$ and $n$ is a power of $q$.

If $n$ is not a prime power, no pair $(q,m)$ qualifies and every shell integral
vanishes.

If $n=p^k$, the only surviving term is $q=p$, $m=k$. Then $P=\{p\}$, and for
$r\in S_f\smallsetminus\{p\}$ the $r$-coordinate of $j_p(p^kw)$ is
$p^{-k}=1/n$, which lies in $n^{-1}(1+r^N\Z_r)$ for \emph{every} $N$, being the
centre of that coset. So every indicator equals $1$ and
\[
  a_p(p^kw)=h(p^{-k})=h(1/n),\quad\text{independent of }w\in\Z_p^\times .
\]
By Proposition~\ref{prop:trace}(a), $\tau_p\equiv p^{-k/2}$ on the shell, and
$\vol_{d^*x}(p^k\Z_p^\times)=\log p$ by \eqref{eq:shellvol}. Hence
\[
  \int_{p^k\Z_p^\times}a_p\,\tau_p\,d^*x=(\log p)\,p^{-k/2}\,h(1/n)
   =\Lambda(n)\,n^{-1/2}h(1/n).
\]

\emph{Unit packets at $q\in S_f$, $m=0$.} By Proposition~\ref{prop:slices},
$j_q(w)=(1$ at $\infty$; $w$ at $q$; $1$ elsewhere$)$ for $w\in\Z_q^\times$,
whose archimedean coordinate is $1$. Since $1\notin\supp h$, we get $a_q(w)=0$
for every $w\in\Z_q^\times$ and $a_q(1)=0$; hence
\[
  (\log q)\int_{\Z_q^\times}\frac{a_q(u)-a_q(1)}{|1-u|_q}\,du=0
\]
for \emph{every} $q$, whether $q\mid n$ or $q\nmid n$. The unit slice sits at
archimedean coordinate $1$, and the Frobenius is supported at archimedean
coordinate $1/n\ne1$, so no cancellation or principal value is involved.

Summing gives \eqref{eq:frobcontacteq}, and $\Lambda(n)=0$ for indices $n$ with at least two distinct prime divisors
matches the vanishing of all shell terms.

\emph{Archimedean.} $j_\infty(x)=(x$ at $\infty$; $1$ at every finite place$)$,
so $b:=F\circ j_\infty$ is
$b(x)=h(x)\prod_{r\in S_f\smallsetminus P}\mathbf 1_{n^{-1}(1+r^N\Z_r)}(1)$, and
the $r$-th indicator equals $1$ iff $n\in1+r^N\Z_r$, i.e. iff $r^N\mid n-1$.
By Proposition~\ref{prop:trace}(b) with the measure $dx/(2|x|)$, and
writing $b_\pm(t):=F(j_\infty(\pm e^t))$ in the logarithmic variable,
\[
  W_\infty(F)=\int_0^\infty K_s(t)\bigl[b_+(t)+b_+(-t)\bigr]dt
   +\int_0^\infty K_c(t)\bigl[b_-(t)+b_-(-t)\bigr]dt .
\]
Since $F$ is supported in $x_\infty>0$, $b_-\equiv0$; and since
$e^t>1>n^{-1}e^\varepsilon$ for $t>0$, $b_+(t)$ vanishes for $t>0$, so only
$b_+(-t)=h(e^{-t})$ survives,
giving \eqref{eq:frobarch}. For each $r\in S_f\smallsetminus P$ the condition
$r^N\mid n-1$ fails once $r^N>n-1$, so the whole term is $0$ for all large $N$
whenever $S_f\smallsetminus P(n)\ne\emptyset$.
\end{proof}

\begin{corollary}\label{cor:frobnoarch}
Let $n\ge2$ and $S_f\supsetneq P(n)$. Then
$\supp\Frob_n\cap\Iml\,j_\infty=\emptyset$, so the Frobenius object has no
archimedean contact. Dually, by Corollary~\ref{cor:liftnotcontact} the element
class $\delta_S(1/n)$ lies on $\Iml\,j_\infty$ and on no $\Iml\,j_q$ for
$|S_f|\ge2$, so it has an archimedean density and no finite contact.  The two
families are supported on complementary slices of $C_S$.
\end{corollary}

\begin{proof}
A point of $\Iml\,j_\infty$ has every finite coordinate equal to $1$; a point
of $\supp\Frob_n$ has $r$-coordinate $1/n$ for $r\in S_f\smallsetminus P(n)$ by
Lemma~\ref{lem:frobsupport}, and $1/n\ne1$ in $\Z_r^\times$ because $n\ge2$ and
$r\nmid n$, so $v_r(n-1)<\infty$ and $n^{-1}\notin1+r^m\Z_r$ for
$m>v_r(n-1)$; in particular $n^{-1}\ne1$.
\end{proof}

\begin{remark}\label{rem:profilevalue}
The shell integral in the proof of Theorem~\ref{thm:frobcontact} evaluates the
archimedean profile at the single point $x_\infty=p^{-k}$, because $a_p$ is
\emph{constant} on the shell. What appears is therefore the \emph{value}
$h(1/n)$, not the integral $\int h$. With an $L^1$-normalized bump
$h=\rho_\varepsilon$ one has $h(1/n)\asymp1/\varepsilon\to\infty$, so the
quantity $\Lambda(n)n^{-1/2}\int\rho_\varepsilon$ is not the contact of
$\Frob_n^{[\rho_\varepsilon,N]}$. The normalization-free invariant is recovered
as in \S\ref{sec:radtrace}, through the representing measure; this is
Corollary~\ref{cor:frobinvariant}.
\end{remark}

\begin{corollary}\label{cor:frobinvariant}
Fix $N$. The map $h\mapsto\sum_qW_q(\Frob_n^{[h,N]})$ is, by
\eqref{eq:frobcontacteq}, integration of $h$ against the measure
$\Lambda(n)n^{-1/2}\delta_{1/n}$ on $\R_{>0}$, independently of $N$. Define
\begin{equation}\label{eq:frobcontactdef}
  \cfin(\Frob_n):=\Lambda(n)\,n^{-1/2},
\end{equation}
the mass of that atom. It is the unique such invariant, a Radon measure being
determined by its $C_c^\infty$ integrals, and for $n=p^k$
\begin{equation}\label{eq:frobvolume}
  \cfin(\Frob_{p^k})=\vol_{d^*x}\bigl(j_p(p^k\Z_p^\times)\bigr)\cdot p^{-k/2}
   =(\log p)\,p^{-k/2},
\end{equation}
the multiplicative Haar volume of the support of the object itself, times its
half-density.
\end{corollary}

\begin{remark}\label{rem:lambdaisownvolume}
Formula \eqref{eq:frobvolume} is a statement about a named object.
$\Frob_{p^k}$ \emph{is} the Haar measure of the stratum $j_p(p^k\Z_p^\times)$,
and \eqref{eq:frobvolume} says that its finite contact is the volume of that
stratum against the half-density. Note the two measures on one set:
$\Frob_{p^k}$ is the \emph{probability} Haar measure of the stratum, of total
mass $1$, while $\Lambda(p^k)=\log p$ is the $d^*x$-volume of the same set.
They are related by the scalar $\log p$, and the probability normalization of
$\Frob_{p^k}$ is what makes the object a correspondence of mass one.
\end{remark}

\begin{theorem}\label{thm:compositedisjoint}
Let $n\ge2$. Then
\[
  \cfin(\Frob_n)\ne0
  \iff n\ \text{is a prime power}
  \iff \supp\Frob_n\ \text{meets some shell}
\]
of some $\Iml\,j_q$ in a set of positive Haar measure; and in that case the intersection is the \emph{whole} shell, that is,
$\supp\Frob_{p^k}=j_p(p^k\Z_p^\times)\subseteq\Iml\,j_p$. For indices $n$ with at least two distinct prime divisors the
support meets $\Iml\,j_q$ for no $q$ at all.
\end{theorem}

\begin{proof}
For $n=p^k$ the support is $c_nK_{\{p\}}=j_p(p^k)\cdot j_p(\Z_p^\times)
=j_p(p^k\Z_p^\times)$ by Lemma~\ref{lem:frobsupport}, the coordinates at
$r\ne p$ being $p^{-k}$ exactly as Proposition~\ref{prop:slices} prescribes;
this is the whole shell. For $n$ with at least two prime divisors, a point of
$\supp\Frob_n$ has archimedean coordinate $1/n$, while a point of
$\Iml\,j_q$ has archimedean coordinate $q^{-m}$ for some $m\in\Z$; and
$1/n=q^{-m}$ forces $n=q^m$, contradicting $|P(n)|\ge2$. So the supports are
disjoint.
\end{proof}

Theorem~\ref{thm:compositedisjoint} is a statement about the object, different
for every $n$ by Theorem~\ref{thm:frobmonoidal}(iv): the vanishing of the
contact at indices $n$ with at least two distinct prime divisors is the \emph{disjointness of the support from every
valuation stratum}. In particular $\Frob_6$ and $\Frob_4$ have different supports.

\subsection{The archimedean flow}

\begin{proposition}\label{prop:archflow}
For $t\in\R$ set $\Frob_\infty(t):=\delta_{j_\infty(e^{-t})}$, the point mass
at $(e^{-t}$ at $\infty$; $1$ at every finite place$)$. Then:
\begin{enumerate}
\item[(i)] $\Frob_\infty$ is a one-parameter \emph{group}:
$\Frob_\infty(s)\star\Frob_\infty(t)=\Frob_\infty(s+t)$,
$\Frob_\infty(0)=\Delta$, $\Frob_\infty(t)^\dagger=\Frob_\infty(-t)$;
\item[(ii)] $\Frob_\infty(\log n)=\delta_S(1/n)$ is the element class of
Lemma~\ref{lem:classunit};
\item[(iii)] its finite contact is $0$ for every $t\ne0$ when $|S_f|\ge2$, and
its archimedean invariant is the \emph{density} $\omega_\infty(t)=K(|t|)$, not
an atom.
\end{enumerate}
\end{proposition}

\begin{proof}
(i) $j_\infty$ is an injective homomorphism $\R^\times\to C_S$ and push-forward
of point masses along a homomorphism is multiplicative. (ii) is
Lemma~\ref{lem:classunit}. (iii) is Corollary~\ref{cor:liftnotcontact} together
with Theorem~\ref{thm:repmeasures}, restricted to the positive component.
\end{proof}

\begin{remark}\label{rem:twoflows}
The trace has exactly two kinds of local term, and the two families exhibit the
objects they measure:
\begin{center}
\begin{tabular}{ll}
\toprule
finite places & $\Frob_n$, $n$ a prime power: discrete, supported\\
 & on the valuation strata $j_p(p^k\Z_p^\times)$;\\
 & contact an \emph{atom} of mass $\Lambda(n)n^{-1/2}$.\\
\midrule
archimedean place & $\Frob_\infty(t)$, $t\in\R$: a continuous\\
 & one-parameter group on the slice $\Iml\,j_\infty$;\\
 & invariant a \emph{density} $K(|t|)$.\\
\bottomrule
\end{tabular}
\end{center}
The finite family is a monoid with a true unit and proper idempotents; the
archimedean family is a group. They meet only at $n=1$, $t=0$, where both are
$\Delta$.
\end{remark}

\begin{theorem}\label{thm:frobincompletion}
Let $n\ge2$ and let $F_{\varepsilon,N}:=\Frob_n^{[h_\varepsilon,N],\flat}$ with
$h_\varepsilon$ an $L^1(d^*x_\infty)$-normalized bump satisfying
\eqref{eq:hsupport}. Then:
\begin{enumerate}
\item[(i)] $\Frob_n\star G^\dagger\in\DA$ for every $G\in\DA$, so
$\ell_n^F(G):=T(\Frob_n\star G^\dagger)=I_A(\Frob_n,G)$ is defined, and
$|\widehat{\Frob_n}|\le1$ pointwise, whence
$|\ell_n^F(G)|\le\mathcal N(G)$;
\item[(ii)] $I_A(F_{\varepsilon,N},\cdot)\to\ell_n^F$ in the strong dual as
$\varepsilon\to0$, $N\to\infty$;
\item[(iii)] $M(F_{\varepsilon,N})\to(n^{1/2},n^{-1/2})=M(\Frob_n)$;
\item[(iv)] hence $J(\Frob_n):=(\ell_n^F,(n^{1/2},n^{-1/2}))
\in\mathcal N_A^{\mathrm{dist}}$, with primitive component $(\ell_n^F,0)$ and
nonzero ruling component; in particular the Frobenius classes are not
primitive.
\end{enumerate}
\end{theorem}

\begin{proof}
(i) $\Frob_n$ is a compactly supported measure, so $\Frob_n\star G^\dagger$ is
smooth in $x_\infty$; it is locally constant in the compact coordinates,
being an average of translates of $G^\dagger$ over a compact group, and the
average of a locally constant function over a compact subgroup is again locally
constant because the function is invariant under a fixed open subgroup; and it
is compactly supported. For the Fourier transform, for a character $\chi$ of
$C_S$,
\[
  \widehat{\Frob_n}(\chi)=\int\overline\chi\,d\Frob_n
   =\overline{\chi(c_n)}\int_{K_{P(n)}}\overline{\chi(u)}\,du
   =\overline{\chi(c_n)}\cdot\bigl[\chi|_{K_{P(n)}}=1\bigr],
\]
by orthogonality on the compact group $K_{P(n)}$; so $|\widehat{\Frob_n}|\le1$
and the bound is Proposition~\ref{prop:completion}.

(ii) Let $B\subseteq\DA$ be bounded, and take $N_B$, $d_{N_B}$, $C''_B$ as in
Proposition~\ref{prop:completion}; these depend on $B$ alone. Take $N\ge N_B$.
Since every $G\in B$ is $U_{N_B}$-invariant, $\widehat G$ is supported on the
same $d_{N_B}$ components of $\widehat{C_S}$ for every $G\in B$, and the
integrand below is supported there; so the component count enters only through
the constant $C_B:=d_{N_B}C''_B\max(A_B,B_S)$, which is independent of $N$.
Since the bound below holds for every $N\ge N_B$ with the same constant, the
limits $\varepsilon\to0$ and $N\to\infty$ may be taken in either order. On a component carrying a
character $\kappa_r$ trivial on $1+r^{N_B}\Z_r\supseteq1+r^N\Z_r$, the
$r$-factor of $\widehat F_{\varepsilon,N}$ is
\[
  \vol(1+r^N\Z_r)^{-1}\int_{n^{-1}(1+r^N\Z_r)}\overline{\kappa_r(x)}\,dx
   =\overline{\kappa_r(n^{-1})},
\]
which is exactly the $r$-factor of $\widehat{\Frob_n}$; and the $p$-factors for
$p\in P(n)$ agree for the same reason, both being
$[\kappa_p=1]\cdot\overline{\kappa_p(\cdot)}$ computed on $\Z_p^\times$. Hence
on every component met by $B$,
\[
  \widehat F_{\varepsilon,N}(\chi)-\widehat{\Frob_n}(\chi)
   =\widehat{\Frob_n}(\chi)\Bigl(\frac{\widehat h_\varepsilon(\tau)}
     {(1/n)^{-i\tau}}-1\Bigr),
\]
whose modulus is at most $\min(2,|\tau|\varepsilon)$ by
$|e^{i\theta}-e^{i\theta'}|\le|\theta-\theta'|$ and
$|\log x_\infty+\log n|\le\varepsilon$ on $\supp h_\varepsilon$. Combining this
with $|\widehat G|\le C''_B\min(1,\tau^{-2})$ and
$|\widehat T|\le A_B+B_S\log(2+|\tau|)$ (Proposition~\ref{prop:completion}) gives, for $R>0$,
\[
  \sup_{G\in B}\bigl|I_A(F_{\varepsilon,N},G)-\ell_n^F(G)\bigr|
   \le C_B\,\varepsilon\,R^2\log R+\eta(R),
\]
where
\[
  \eta(R):=2C''_B\int_{|\tau|>R}\tau^{-2}
    \bigl(A_B+B_S\log(2+|\tau|)\bigr)\frac{d\tau}{2\pi},
\]
uniformly on $B$; choose $R$ first and then $\varepsilon$.

(iii) $\lambda_\mp(F_{\varepsilon,N})=\int F_{\varepsilon,N}(w)|w|_S^{\mp1/2}
d^*w=\int h_\varepsilon(x)|x|^{\mp1/2}d^*x_\infty\to n^{\pm1/2}$, since
$h_\varepsilon$ concentrates at $x=1/n$ with unit mass.

(iv) $J(F)=(I_A(F,\cdot),M(F))$; combine (ii) and (iii), and each
$J(F_{\varepsilon,N})$ lies in $J(\DA)$, so the limit lies in its closure. The
splitting is Proposition~\ref{prop:completion}, and $M(\Frob_n)\ne0$.
\end{proof}

\section{The integral models}\label{sec:integral}

This section is $\Z$-linear.  We construct two derived contacts attached to
the index $n$, show that they have coprime support on $\Spec\Z$, and identify
primitive extraction with M\"obius inversion over the divisor poset, an
operation which we show cannot be monoidal.  Nothing in this section takes
place in the correspondence algebra; the comparison with it, in
Section~\ref{sec:shadow}, is an identity of degrees.

\subsection{The two contacts}

Let $n\ge2$. Over $R=\Z$ or $\Z_p$ we have the following two complexes.

\emph{The ramification contact.} Let $\pi_n:\mathbb G_{m,\Z}\to
\mathbb G_{m,\Z}$ be the power map $z\mapsto z^n$ and let
$e:\Spec\Z\to\mathbb G_{m,\Z}$ be the identity section. Since $\pi_n$ is a
morphism of smooth affine $\Z$-schemes, its cotangent complex is the two-term
complex $L_{\pi_n}=[\pi_n^*\Omega_{\mathbb G_m/\Z}\to\Omega_{\mathbb G_m/\Z}]$
in degrees $-1,0$ \cite[II.1.2]{Illusie}, \cite[Tag 08UQ]{Stacks}. Now
$\Omega_{\mathbb G_m/\Z}$ is free of rank one on $dz/z$, and
$\pi_n^*(dz/z)=d(z^n)/z^n=n\,dz/z$; pulling back along $e$, which is flat over
the base in the sense that $e^*$ of a free module is free, we obtain
\begin{equation}\label{eq:cram}
  C_{\mathrm{ram}}(n):=e^*L_{\pi_n}=[\Z\xrightarrow{\ n\ }\Z],\qquad
  H^{-1}=0,\quad H^0=\Z/n .
\end{equation}

\emph{The fixed-graph contact.} By Lemma~\ref{lem:deformation} below, the fixed
equation of the correspondence, written in the coordinate $w$ in which the
morphism is $w\mapsto w^n$, is $w^{n-1}=1$. Linearising it at the fixed point
$w=1$, i.e. substituting $w=1+u$ and taking the linear term
$(n-1)u+O(u^2)$, we \emph{define} the fixed contact to be the Koszul complex
of that equation,
\begin{equation}\label{eq:cfix}
  C_{\mathrm{fix}}(n):=[R\xrightarrow{\ n-1\ }R],\qquad
  H^{-1}=0,\quad H^0=R/(n-1),
\end{equation}
over $R=\Z$ or $\Z_p$.

These are not quasi-isomorphic, decisively so at $n=p^k$ localized at $p$:
there $n-1\in\Z_p^\times$ makes \eqref{eq:cfix} contractible, while the
primitive extraction of \eqref{eq:cram} has cohomology $\Z/p$.  The following
lemma explains why both differentials occur.

\begin{lemma}\label{lem:deformation}
Consider the integral deformation $F_n(t,z)=((1+tz)^n-1)/t$, whose generic
fibre is $\pi_n$ in the coordinate $w=1+tz$ and whose special fibre $t=0$ is
$d_n$. In the coordinate $w$:
\[
  \text{fixed locus of the correspondence}:\ w^{n-1}=1,\ \text{i.e. }\mu_{n-1},
\]
\[
  \text{relative fibre over the unit}:\ w^n=1,\ \text{i.e. }\mu_n .
\]
\end{lemma}

\begin{proof}
$1+tF_n(t,z)=(1+tz)^n$, i.e. $F_n$ is $w\mapsto w^n$. The fixed equation
$F_n(t,z)=z$ reads $w^n=w$, so $w(w^{n-1}-1)=0$, and on the invertible locus
$w^{n-1}=1$. The relative fibre $F_n(t,z)=0$ reads $w^n=1$.
\end{proof}

So the fixed data are governed by $\mu_{n-1}$ and the ramification data by
$\mu_n$: the fixed group and the deck group of the same morphism, which have
coprime orders.

\begin{theorem}[Coprime support]\label{thm:separation}
Let $n\ge2$. Then
\begin{enumerate}
\item[(i)] $\supp H^0(C_{\mathrm{ram}}(n))=V(n)$ and
$\supp H^0(C_{\mathrm{fix}}(n))=V(n-1)$;
\item[(ii)] $V(n)\cap V(n-1)=\emptyset$;
\item[(iii)] for every finite place $p$,
\[
  p\mid n\implies|1-n|_p=1\ \text{and}\ \tau_p(n)=|n|_p^{1/2},\qquad
  p\nmid n\implies|n|_p=1\ \text{and}\ \tau_p(n)=|1-n|_p^{-1};
\]
\item[(iv)] consequently the product of the local factors of the idele
$a_S(n)$ over all finite places of $\Q$,
\[
  w(n):=\prod_{p\ \mathrm{finite}}\tau_p(n),
\]
factors place-wise and canonically as $w(n)=w_{\mathrm{ram}}(n)\cdot
w_{\mathrm{unit}}(n)$ with $w_{\mathrm{ram}}(n)=\prod_{p\mid n}\tau_p(n)$ and
$w_{\mathrm{unit}}(n)=\prod_{p\nmid n}\tau_p(n)$, and
\begin{equation}\label{eq:wramunit}
  w_{\mathrm{ram}}(n)=n^{-1/2}=\bigl(\det C_{\mathrm{ram}}(n)\bigr)^{-1/2},
  \qquad
  w_{\mathrm{unit}}(n)=n-1=\#H^0(C_{\mathrm{fix}}(n)).
\end{equation}
\end{enumerate}
\end{theorem}

\begin{proof}
(i) $H^0(C_{\mathrm{ram}}(n))=\Z/n$ and $H^0(C_{\mathrm{fix}}(n))=\Z/(n-1)$ by
\eqref{eq:cfix}--\eqref{eq:cram}; the support of a finite cyclic group $\Z/m$
in $\Spec\Z$ is $V(m)$.

(ii) $\gcd(n,n-1)=1$.

(iii) If $p\mid n$ then $n\equiv0\pmod p$, so $1-n\equiv1$ is a unit,
$|1-n|_p=1$, and $\tau_p(n)=|n|_p^{1/2}$. If $p\nmid n$ then $|n|_p=1$ and
$\tau_p(n)=|1-n|_p^{-1}$. The two cases are exhaustive and exclusive, and the
case distinction is a property of $n$ alone, so the factorization involves no
choice.

(iv) By (iii),
\[
  w_{\mathrm{ram}}(n)=\prod_{p\mid n}|n|_p^{1/2}
   =\Bigl(\prod_{p\ \mathrm{finite}}|n|_p\Bigr)^{1/2}
   =\bigl(|n|_\infty^{-1}\bigr)^{1/2}=n^{-1/2},
\]
using $|n|_p=1$ for $p\nmid n$ and the product formula for the rational integer
$n$; and $\det C_{\mathrm{ram}}(n)=n=\#H^0(C_{\mathrm{ram}}(n))$. Likewise
$w_{\mathrm{unit}}(n)=\prod_{p\nmid n}|1-n|_p^{-1}$, and since $|1-n|_p=1$ for
$p\mid n$ this equals $\prod_{p\ \mathrm{finite}}|1-n|_p^{-1}$, which by the
product formula applied to $n-1\in\Q^\times$ equals $|n-1|_\infty=n-1$; and
$\#H^0(C_{\mathrm{fix}}(n))=n-1$.
\end{proof}

Both products in (iv) must be taken over \emph{all} finite places of $\Q$: on
the minimal chart $S=\{\infty\}\cup P(n)$ the second product is empty, hence
$1$. This is legitimate and canonical because $\tau_p=1$ off the finite set
$V(n)\cup V(n-1)$, and a chart sees the whole quantity as soon as
$S\supseteq V(n)\cup V(n-1)$.  Note also that $a_S(n)$ has archimedean
component $1$, where the local factor is undefined; the identity class is
accounted for by the constant $c=0$.

\begin{remark}\label{rem:aboutthelift}
Parts (i)--(iii) are statements about the two integral complexes and about
$\tau_p$ as a function on $\Q_p^\times$. Parts (iv) are statements about the
\emph{idele} $a_S(n)$, a chosen lift, and about nothing else. They are not
statements about the class $\delta_S(1/n)$ in $C_S$, and they are not
values of, or contributions to, any local term $W_q$ of the trace: by
Proposition~\ref{prop:noslice} the class has unit coordinate $1$ at every
finite place and is in the image of no finite slice map, and
Corollary~\ref{cor:liftnotcontact} states the consequence.  In particular
$w_{\mathrm{unit}}(n)=n-1$ is a canonical number attached to the lift, but it
is not the contact of any object on $C_S$. That the number $n^{-1/2}$
is also the canonical half-density of the class, and the second coordinate of
the bidegree $M(\Gamma_n^{\mathrm{rad}})$, is an \emph{independent} fact, proved
in Theorem~\ref{thm:radmonoidal}(iii) from $|\delta_S(1/n)|_S=1/n$; the
coincidence is the product formula.
\end{remark}

\begin{remark}\label{rem:whatsep}
Theorem~\ref{thm:separation} produces no complex at $p\mid n$; it observes that
the number $|n|_p^{1/2}$ appears there. What is correct, and what is used, is
the \emph{support} statement: the two operations have coprime support, so they
never compete, and the numerical data of each appear at disjoint sets of places
with their own exponents. The factorization exists only place-wise: there is no
function of the modulus $|a_S(n)|_S$ alone that separates $w_{\mathrm{ram}}$
from $w_{\mathrm{unit}}$, because both are products of local factors of the
same idele. Classically this is the familiar dichotomy in the local Weil
distribution, where the integrand $1/|1-u|_v$ is bounded and equal to $1$ on
$|u|_v\ne1$ and divergent on average on the units \cite{Weil1952}; part (iii)
is that dichotomy read on a single arithmetic idele.
\end{remark}

\begin{lemma}\label{lem:tauinv}
For every place $v$ and every $\sigma_v\ne1$, $\tau_v(\sigma_v^{-1})
=\tau_v(\sigma_v)$.
\end{lemma}

\begin{proof}
$|1-\sigma^{-1}|_v=|\sigma^{-1}|_v|\sigma-1|_v=|\sigma|_v^{-1}|1-\sigma|_v$, so
\[
  \tau_v(\sigma^{-1})=\frac{|\sigma|_v^{-1/2}}{|\sigma|_v^{-1}|1-\sigma|_v}
   =\frac{|\sigma|_v^{1/2}}{|1-\sigma|_v}=\tau_v(\sigma). \qedhere
\]
\end{proof}

Place by place this is the symmetry behind the pairing of an index with its
dagger, and hence behind the two terms $h(\log n)+h(-\log n)$ of
\eqref{eq:radtrace}. It also removes an apparent ambiguity in the comparison of
the two contacts: the integral complexes $[\Z\to1-u\to\Z]$ and
$[\Z\to1-u^{-1}\to\Z]$ differ by the unit $u^{-1}$ when $u$ is invertible, hence
are isomorphic, and by Lemma~\ref{lem:tauinv} the analytic factor does not
distinguish them either.

\subsection{The index category and its realization}

\begin{definition}\label{def:indexcat}
The \emph{index category} $W$ has objects the integers $n\ge1$, exactly one
morphism $d\to n$ when $d\mid n$ and none otherwise, composition being
transitivity of divisibility, monoidal product $n\otimes m=nm$ on objects and
$(d\to n)\otimes(e\to m)=(de\to nm)$ on morphisms, and unit $1$. It is a
symmetric monoidal poset category.
\end{definition}

Two structures carried by $W$ must be kept apart, because the next subsection
depends on the distinction: the \emph{monoid} $(\N_{\ge1},\times)$ of objects, which
the realization preserves; and the \emph{divisor poset} $D(n)=\{d:d\mid n\}$
for each $n$, whose arrows are the maps $i_{d,n}=(\mathrm{id},n/d)$ on
two-term complexes and whose cube is the site of primitive extraction. These
maps satisfy $i_{b,n}i_{d,b}=i_{d,n}$ \emph{literally}, so $d\mapsto C_d$ is a
strict functor $D(n)\to\Ch(\Z)$, not a diagram commuting up to coherent
homotopy.

\begin{theorem}\label{thm:realization}
Fix a chart $S$ and, for $n\ge1$, the semilocal idele
$a_S(n)=(1$ at $\infty$; $n$ at each finite place of $S)$, so that
$a_S(m)a_S(n)=a_S(mn)$; set $\Gamma_n:=E_{a_S(n)}$. Then $n\mapsto\Gamma_n$ is
a strong monoidal functor from the monoid $(\N_{\ge1},\times)$ to the
correspondence algebra, with $\Gamma_m\star\Gamma_n\cong\Gamma_{mn}$ and
$\Gamma_1=\Delta$. There is, however, \emph{no} functor $W\to$ (correspondence
algebra): morphisms between distinct external classes vanish, so the divisor
arrows of $W$ have no image, and the symmetry of $W$ is likewise not realized.
Moreover $\Gamma_n$ is independent of the chosen lift of the external class, of
the uniformizers and of the chart: the finite arithmetic anchor
$b_S(a)=a_f/a_\infty$ is invariant under $a\mapsto\gamma a$ for a positive
rational $S$-unit $\gamma$, is multiplicative, and takes the value $n$ on
$a_S(n)$, so the invariant lattice $nL_S$ is attached to the class; and for
$S\subseteq T$ the lift transition carries $a_S(n)$ to $a_T(n)$ literally, the
coordinate projection $L_T\to L_S$ inducing a quasi-isomorphism of index
complexes which composes exactly on chart triples and commutes with all divisor
maps.
\end{theorem}

\begin{proof}
Monoidality is $a_S(m)a_S(n)=a_S(mn)$ together with $E_u\star E_v\cong E_{uv}$
and $E_1=\Delta$; the coherences hold because the multiplication
$\mu_{\beta,\eta}(x\otimes y)=x\beta(y)$ is associative on elementary tensors
and extends by density. The failure of functoriality on $W$ is the vanishing of
$\Hom$ between distinct external classes. Gauge invariance of $b_S$ is
immediate from $b_S(\gamma a)=\gamma a_f/(\gamma a_\infty)=a_f/a_\infty$; for
the chart statement, the kernel of the projection $L_T\to L_S$ is multiplication
by $n$ on the added $\Z_p$ with $p\nmid n$, hence contractible.
\end{proof}

\begin{theorem}\label{thm:unittransphalf}
\begin{enumerate}
\item[(i)] $\Gamma_1=\Delta$ is the monoidal unit and $\Delta^\dagger=\Delta$.
\item[(ii)] $\Gamma_n^\dagger=E_{a_S(n)^{-1}}$, and
$(\Gamma_m\star\Gamma_n)^\dagger\cong\Gamma_n^\dagger\star\Gamma_m^\dagger$.
\item[(iii)] The half-density of $\Gamma_n$ is canonical: $|a_S(n)|_S=1/n$, so
the half-density weight is $n^{-1/2}$, and the product formula
($|\gamma|_S=1$ for $\gamma\in\Gamma_S$) makes it independent of the lift.
There is no two-torsion ambiguity.
\end{enumerate}
\end{theorem}

\begin{proof}
(i) and (ii) are the definition of the involution applied to the objects
$E_{a_S(n)}$. For (iii): $a_S(n)$ has entry $1$ at $\infty$ and $n$ at every
$p\in S_f$, and $S$ contains every prime dividing $n$, so
$|a_S(n)|_S=\prod_{p\in S_f}|n|_p=\prod_{p\ \mathrm{finite}}|n|_p
=|n|_\infty^{-1}=1/n$.
\end{proof}

\subsection{M\"obius inversion}

\begin{lemma}\label{lem:mobiuscube}
Let $n\ge1$, let $P(n)$ be its set of distinct prime divisors, and let $f$ be
any function on the divisors of $n$ with values in an abelian group. Then
\begin{equation}\label{eq:mobiuscube}
  \sum_{J\subseteq P(n)}(-1)^{|J|}f\Bigl(n\big/\prod_{p\in J}p\Bigr)
   =\sum_{d\mid n}\mu(d)f(n/d).
\end{equation}
\end{lemma}

\begin{proof}
$\mu(d)=0$ unless $d$ is squarefree, so the right-hand side is supported on the
squarefree divisors of $n$. The map $J\mapsto\prod_{p\in J}p$ is a bijection
from subsets of $P(n)$ onto the squarefree divisors of $n$, and
$\mu(\prod_{p\in J}p)=(-1)^{|J|}$. Substituting gives the left-hand side.
\end{proof}

So the cube of radical divisors is a categorification of M\"obius inversion,
and its Euler-type invariants compute $\mu*f$. We make the object precise.

\begin{definition}\label{def:prim}
Let $N>1$ and let $P\subseteq P(N)$ be nonempty. The \emph{radical cube} of
$(N,P)$ is the $|P|$-dimensional commutative cube
\[
  J\longmapsto C_{\mathrm{ram}}\Bigl(N\big/\!\!\prod_{p\in J}p\Bigr),
  \qquad J\subseteq P,
\]
whose edge in the direction $q$, from the vertex at $J\cup\{q\}$ to the vertex
at $J$, is the chain map $i_{d,qd}=(\mathrm{id},q)$ with $d=N/(q\prod_{J}p)$.
Write $\Prim(N,P)$ for its total cofiber, the iterated cone taken along the
$|P|$ directions, and put
\[
  \Prim_n:=\Prim\bigl(n,P(n)\bigr)\qquad(n>1).
\]
\end{definition}

Two remarks on the definition. The cube is indexed by the \emph{radical}
divisors $\prod_{p\in J}p$, not by all divisors of $n$: for $n=p^k$ the divisor
poset is a chain of length $k$, whereas the radical cube is the single edge
$C_{\mathrm{ram}}(p^{k-1})\to C_{\mathrm{ram}}(p^k)$. And the squares commute
strictly: in degree $-1$ both composites are the identity, and in degree $0$
they are multiplication by the same integer, so no coherence data are needed.

\begin{lemma}[Structure of the radical cofiber]\label{lem:primstructure}
Let $N>1$ and let $\emptyset\ne P\subseteq P(N)$. Then
\[
  \Prim(N,P)\simeq
  \begin{cases}
    \Z/q\ \text{in degree }0\ \simeq\ [\Z\xrightarrow{\ q\ }\Z],
      & P=\{q\},\\[2pt]
    0, & |P|\ge2 .
  \end{cases}
\]
In particular, for $n>1$,
\[
  \Prim_n\simeq[\Z\xrightarrow{\ \Phi_n(1)\ }\Z],\qquad
  H^{-1}\Prim_n=0,\qquad H^0\Prim_n=\Z/\Phi_n(1),
\]
so $\Prim_n\cong\Z/p$ in degree $0$ when $n=p^k$ and $\Prim_n\simeq0$ when $n$
has at least two prime divisors.
\end{lemma}

\begin{proof}
We induct on $r=|P|$.

\emph{Base $r=1$, $P=\{q\}$.} The cube is the single edge
$i:C_{\mathrm{ram}}(N/q)\to C_{\mathrm{ram}}(N)$, $i=(\mathrm{id},q)$, and its
total cofiber is $\Cone(i)$. The map $i$ is injective in each degree --- the
identity in degree $-1$ and multiplication by $q$ in degree $0$ --- so
$\Cone(i)$ is quasi-isomorphic to the degreewise cokernel
$[\,0\to\Z/q\,]$, concentrated in degree $0$. Since
$[\Z\xrightarrow{q}\Z]$ is a free resolution of $\Z/q$, this is
$[\Z\xrightarrow{q}\Z]$ in $D(\Z)$.

\emph{Step $r\ge2$.} Fix $q\in P$ and put $P':=P\smallsetminus\{q\}$. Splitting
the cube along the direction $q$ exhibits it as a map of $(r-1)$-cubes: the
face $\{J:q\notin J\}$ is the radical cube of $(N,P')$, the face
$\{J:q\in J\}$ is the radical cube of $(N/q,P')$, and the direction-$q$ edges
give a map $\varphi$ of cubes from the second to the first. Total cofibers
being computed by iterated cones in any order,
\[
  \Prim(N,P)=\Cone\Bigl(\varphi_*:\Prim(N/q,P')\longrightarrow
   \Prim(N,P')\Bigr).
\]

If $r\ge3$ then $|P'|\ge2$, so both source and target are acyclic by the
inductive hypothesis and $\Cone(\varphi_*)$ is acyclic as well.

If $r=2$, write $P=\{q,q'\}$. By the base case, $\Prim(N,\{q'\})$ is the
degree-$0$ cokernel of $(\mathrm{id},q'):C_{\mathrm{ram}}(N/q')\to
C_{\mathrm{ram}}(N)$, namely $\Z/q'$ presented as the degree-$0$ term of
$C_{\mathrm{ram}}(N)$ modulo $q'$; likewise $\Prim(N/q,\{q'\})$ is $\Z/q'$
presented as the degree-$0$ term of $C_{\mathrm{ram}}(N/q)$ modulo $q'$. The
direction-$q$ edges are $(\mathrm{id},q)$ in both cases, so the induced map
$\varphi_*$ is multiplication by $q$ on $\Z/q'$. As $q\ne q'$ are distinct
primes, $q$ is invertible modulo $q'$ and $\varphi_*$ is an isomorphism;
hence $\Cone(\varphi_*)$ is acyclic.

\emph{The consequence for $\Prim_n$.} If $n=p^k$ then $P(n)=\{p\}$ and the
base case gives $\Z/p$ in degree $0$; and $\Phi_{p^k}(x)=\Phi_p(x^{p^{k-1}})$
gives $\Phi_{p^k}(1)=\Phi_p(1)=p$. If $|P(n)|\ge2$ then $\Prim_n\simeq0$ and
$\Phi_n(1)=1$ by Lemma~\ref{lem:cyclotomicvalue} below, so
$[\Z\xrightarrow{\Phi_n(1)}\Z]=[\Z\xrightarrow{1}\Z]$ is acyclic too.
\end{proof}

\begin{lemma}\label{lem:cyclotomicvalue}
For $n>1$ one has $\Phi_n(1)=\prod_{d\mid n}d^{\mu(n/d)}$, hence
$\log\Phi_n(1)=\Lambda(n)$; in particular $\Phi_n(1)=p$ if $n=p^k$ is a prime
power and $\Phi_n(1)=1$ otherwise.
\end{lemma}

\begin{proof}
From $\prod_{d\mid n}\Phi_d(x)=x^n-1$ and $\Phi_1(x)=x-1$ one gets
$\prod_{d\mid n,\,d>1}\Phi_d(1)=\lim_{x\to1}(x^n-1)/(x-1)=n$. (The restriction
$n>1$ is forced here and everywhere below by $\Phi_1(1)=0$.) Setting
$f(1):=1$ and $f(d):=\Phi_d(1)$ for $d>1$, we therefore have
$\prod_{d\mid n}f(d)=n$ for every $n\ge1$, the case $n=1$ being the empty
product. Multiplicative M\"obius inversion gives
$f(n)=\prod_{d\mid n}d^{\mu(n/d)}$, whose logarithm is
$\sum_{d\mid n}\mu(n/d)\log d=(\mu*\log)(n)=\Lambda(n)$ \cite[Ch.~1]{IK}, and the
last assertion is the definition of $\Lambda$.
\end{proof}

\begin{definition}\label{def:degdet}
Call a bounded complex $K$ of abelian groups \emph{of finite type} if every
$H^i(K)$ is finite, and set
\[
  \degdet K:=\sum_i(-1)^i\log\#H^i(K)\ \in\R .
\]
For a two-term complex $[\Z\xrightarrow{m}\Z]$ in degrees $-1,0$ with $m\ne0$
this is $\log\#(\Z/m)=\log|m|$, and for $m=\pm1$ it is $0$.
\end{definition}

\begin{lemma}\label{lem:degdetadd}
$\degdet$ is additive on finite direct sums, changes sign under a shift, and
is additive on exact triangles: if $A\to B\to C\xrightarrow{+1}$ is an exact
triangle of bounded complexes all of which are of finite type, then
$\degdet B=\degdet A+\degdet C$.
\end{lemma}

\begin{proof}
Additivity on direct sums and the behaviour under shifts are immediate from
the definition. For a triangle, the associated long exact sequence
\[
  \cdots\to H^{i}(A)\to H^{i}(B)\to H^{i}(C)\to H^{i+1}(A)\to\cdots
\]
is an exact sequence of finite abelian groups, bounded in both directions.
For any bounded exact sequence $0\to G_1\to G_2\to\cdots\to G_r\to0$ of finite
abelian groups one has $\prod_j\#G_j^{(-1)^j}=1$, by splitting it into short
exact sequences and using multiplicativity of orders. Taking logarithms and
sorting the terms by the complex they come from gives
$\degdet B-\degdet A-\degdet C=0$.
\end{proof}

\begin{theorem}\label{thm:degreelaw}
For every $n>1$,
\begin{equation}\label{eq:degreelaw}
  \degdet\Prim_n=\log\Phi_n(1)=\Lambda(n)=\sum_{d\mid n}\mu(d)\log(n/d)
   =\sum_{J\subseteq P(n)}(-1)^{|J|}\,
     \degdet C_{\mathrm{ram}}\Bigl(n\big/\!\!\prod_{p\in J}p\Bigr).
\end{equation}
For $n=1$ the complex $C_{\mathrm{ram}}(1)$ is acyclic and $\Gamma_1=\Delta$;
the index $1$ is treated as identity or excess by the term $c\cdot F(1)$ with
$c=0$, and not by \eqref{eq:degreelaw}.
\end{theorem}

\begin{proof}
The first equality is Lemma~\ref{lem:primstructure} together with
Definition~\ref{def:degdet}: $\Prim_n\simeq[\Z\xrightarrow{\Phi_n(1)}\Z]$ has
$H^{-1}=0$ and $H^0=\Z/\Phi_n(1)$, so
$\degdet\Prim_n=\log\#(\Z/\Phi_n(1))=\log\Phi_n(1)$, which is $0$ when
$\Phi_n(1)=1$. The second is Lemma~\ref{lem:cyclotomicvalue}. The third is the
identity $\Lambda=\mu*\log$ in the form $\Lambda(n)=-\sum_{d\mid n}\mu(d)\log
d$ together with $\sum_{d\mid n}\mu(d)=0$ for $n>1$, which turns
$\sum_{d\mid n}\mu(d)\log(n/d)$ into $-\sum_{d\mid n}\mu(d)\log d$. The fourth
is Lemma~\ref{lem:mobiuscube} applied to $f=\log$, together with
$\degdet C_{\mathrm{ram}}(m)=\log\#(\Z/m)=\log m$, which is
Definition~\ref{def:degdet} and \eqref{eq:cram}.
\end{proof}

\begin{corollary}[The ghost relation]\label{cor:ghost}
$\degdet C_{\mathrm{ram}}(n)=\sum_{d\mid n}\degdet\Prim_d$, i.e.
$\log n=\sum_{d\mid n}\Lambda(d)$: the lattice-index degree of $n$ is the sum
of the primitive degrees of its divisors.
\end{corollary}

\begin{proof}
M\"obius inversion of \eqref{eq:degreelaw}.
\end{proof}

\begin{theorem}\label{thm:notmonoidal}
Primitive extraction is not monoidal for the derived tensor product. In
$D(\Z)$,
\[
  \Prim_2\otimes^L\Prim_2=\Z/2\otimes^L\Z/2=\Z/2\oplus\Z/2[1]
   \ \not\cong\ \Z/2=\Prim_4 .
\]
Hence there is no natural isomorphism $\Prim_m\otimes^L\Prim_n
\cong\Prim_{mn}$.
\end{theorem}

\begin{proof}
Resolve $\Z/2$ by $[\Z\xrightarrow{2}\Z]$ and tensor with $\Z/2$: the complex
becomes $[\Z/2\xrightarrow{0}\Z/2]$ with zero differential, so $H_0=\Z/2$ and
$H_1=\mathrm{Tor}_1^\Z(\Z/2,\Z/2)=\Z/2$, and the complex splits. By
Lemma~\ref{lem:primstructure}, $\Prim_4\cong\Z/2$ concentrated in degree $0$.
\end{proof}

\begin{corollary}\label{cor:notadditive}
Since $\Lambda(6)=0\ne\Lambda(2)+\Lambda(3)=\log6$, no monoidal primitive
extraction with an additive degree invariant can satisfy
$\deg\Prim_n=\Lambda(n)$ either.
\end{corollary}

\begin{corollary}\label{cor:structurethatholds}
The structure that does hold is:
\begin{enumerate}
\item[(i)] the realization $n\mapsto\Gamma_n$ \emph{is} strong monoidal
(Theorem~\ref{thm:realization}), and its degree
$\degdet C_{\mathrm{ram}}(n)=\log n$ \emph{is} additive;
\item[(ii)] primitive extraction is the M\"obius operation of the divisor
\emph{poset}, not of the monoid (Lemma~\ref{lem:mobiuscube});
\item[(iii)] the two are compatible exactly through Corollary~\ref{cor:ghost}.
\end{enumerate}
\end{corollary}

This is also why the divisor cube must live on an auxiliary $\Z$-linear target
and not on correspondences: morphisms between distinct external classes vanish,
so the divisor arrows have no realization as morphisms of correspondences. They
are realized, strictly, on the index complexes of \S\ref{sec:shadow}.

\subsection{The index complex}

\begin{theorem}\label{thm:indexcomplex}
Let $n\ge2$ and $S\supseteq\{\infty\}\cup P(n)$. Put
$Q_{S,n}:=[L_S\xrightarrow{n}L_S]$ in degrees $-1,0$. Then:
\begin{enumerate}
\item[(i)] $Q_{S,n}$ is perfect and the diagonal $\Z\to L_S$ induces a
quasi-isomorphism $C_{\mathrm{ram}}(n)\to Q_{S,n}$, naturally on the whole
divisor cube;
\item[(ii)] $[L_S:nL_S]=n=|a_S(n)|_S^{-1}$, the inverse modulus of the
correspondence;
\item[(iii)] $-2\log w_{\mathrm{ram}}(n)=\log[L_S:nL_S]
=\degdet C_{\mathrm{ram}}(n)=\log n$.
\end{enumerate}
\end{theorem}

\begin{proof}
(i)--(ii): the cokernel of $n$ on $L_S$ is
$\prod_{p\in S_f}\Z_p/n\Z_p\cong\Z/n$ by primary decomposition, since $S$
contains every prime dividing $n$; multiplication by $n$ is injective on the
torsion-free $L_S$, so $H^{-1}=0$ and the diagonal induces an isomorphism on
$H^0$. The divisor maps $(\mathrm{id},n/d)$ commute strictly with the diagonal,
and the chart projections are quasi-isomorphisms. (iii): by
\eqref{eq:wramunit}, $w_{\mathrm{ram}}(n)=n^{-1/2}$, so
$-2\log w_{\mathrm{ram}}(n)=\log n$, and
$\degdet C_{\mathrm{ram}}(n)=\log n$ by Definition~\ref{def:degdet}.
\end{proof}

Part (iii) is essentially definitional: the half-density of the dilation is
the square root of the inverse index of that dilation on its invariant
lattice, and this index is the determinant of the ramification contact.  The
actual comparison with the analytic trace is Theorem~\ref{thm:degreematching}
below, whose two sides are computed independently, the left from $L_S/nL_S$
over $\Z$ and the right from the local terms \eqref{eq:localterms} and the
shell volumes.

\begin{remark}\label{rem:complexshadow}
Over $\C$ the torsion disappears: by Proposition~\ref{prop:shadow} the shadow
of an external class at a finite place $p\mid n$ is the rank-one class of the
fixed point, and complexification kills the finite torsion. The integral
information is carried by the lattice $L_S$, which is part of the coefficient
geometry, and the only trace of it in the analytic invariant is the
half-density exponent. Theorem~\ref{thm:indexcomplex}(iii) says that this trace
is complete: the half-density determines $\log n$, and $\log n$ determines
$C_{\mathrm{ram}}(n)$ up to quasi-isomorphism, a two-term complex
$[\Z\xrightarrow{m}\Z]$ with $m>0$ being determined by $m$.
\end{remark}

\begin{theorem}\label{thm:fourmodels}
For $n>1$ the primitive contact attached to the index $n$ is the perfect
complex
\begin{equation}\label{eq:Kn}
  K_n:=[\Z\xrightarrow{\Phi_n(1)}\Z],\qquad H^0=\Z/\Phi_n(1),\ H^{-1}=0,
\end{equation}
with no truncation. The constructions (a)--(c) below agree with $K_n$ in
$D(\Z)$, while (d) is a complex of $\Z[t]$-modules, namely
$K_n\otimes_\Z\Z[t]$, whose specialisation at every $t=c\in\Z$ is $K_n$:
\begin{enumerate}
\item[(a)] $\Prim_n$, the total cofiber of the radical cube of
$C_{\mathrm{ram}}(\cdot)$ (Definition~\ref{def:prim});
\item[(b)] the same radical cube taken on the index complexes
$Q_{S,\cdot}$, via the natural quasi-isomorphism of
Theorem~\ref{thm:indexcomplex}(i);
\item[(c)] $D_n=[V_n\xrightarrow{1-T}V_n]$ where $V_n=\Z[T]/\Phi_n(T)$ inside
the orbit module $\Z[L_S/nL_S]$ and $T$ is translation by the class of the
diagonal integer $1$, with $H^0(D_n)=\Z/\Phi_n(1)$ and
$\det(1-T\mid V_n)=\Phi_n(1)$;
\item[(d)] $K_{\mathrm{prim}}(n,t)=[\Z[t]\xrightarrow{\Phi_n(1)}\Z[t]]$, the
derived intersection of the primitive cyclotomic factor
$P_n(t,z)=\Phi_n(1+tz)$ of the relative fibre with the zero section.
\end{enumerate}
\end{theorem}

\begin{proof}
For (a), Lemma~\ref{lem:primstructure} gives $H^{-1}\Prim_n=0$ and
$H^0\Prim_n=\Z/\Phi_n(1)$. (b) follows from (a) by the natural
quasi-isomorphism of the whole cube in Theorem~\ref{thm:indexcomplex}(i). For
(c), the primitive submodule of $\Z[L_S/nL_S]\cong\Z[\Z/n]=\Z[T]/(T^n-1)$ is
$V_n=\Z[T]/\Phi_n(T)$, and $\det(1-T\mid V_n)=\prod_{\zeta}(1-\zeta)$ over the
primitive $n$-th roots of unity, which is $\Phi_n(1)$; the cokernel of $1-T$ on
$V_n$ is $\Z[T]/(\Phi_n(T),1-T)=\Z/\Phi_n(1)$, and $1-T$ is injective on $V_n$
because $\Phi_n(1)\ne0$. For (d), $P_n$ is a nonzerodivisor in $\Z[t,z]$ and
$P_n(0,z)=\Phi_n(1)\ne0$, so ordinary and derived specialization agree and the
Koszul resolution gives $K_{\mathrm{prim}}(n,t)=K_n\otimes_\Z\Z[t]$, whose
specialisation at $t=c$ is $K_n$; as a complex of $\Z[t]$-modules it has
$H^0=\Z[t]/\Phi_n(1)$ and is not an object of $D(\Z)$, which is why it is
listed separately. The complexes (a), (b), (c) all have $H^{-1}=0$ and
$H^0=\Z/\Phi_n(1)$, and \eqref{eq:Kn} is a free resolution of that group, so
their agreement is an isomorphism in $D(\Z)$ and not merely an equality of
orders.
\end{proof}

Of these, (b) is transported from (a) and (c) is built on the same finite
quotient, so there are two independent constructions, the cotangent--lattice
cube and the deformation Koszul complex. Combining
Theorem~\ref{thm:fourmodels} with Theorem~\ref{thm:degreelaw} gives
\begin{equation}\label{eq:degcontact}
  \degdet(\text{primitive contact of the index } n)=\log|\Phi_n(1)|=\Lambda(n),
  \qquad \text{analytic weight }=n^{-1/2},
\end{equation}
the second by Theorem~\ref{thm:unittransphalf}(iii) and \eqref{eq:wramunit},
unchanged by primitive extraction because extraction acts on the radical cube
of contacts and not on the index.

\begin{theorem}[The two contacts at $p$ and at $p^2$]\label{thm:discriminating}
Let $p$ be a prime.
\begin{enumerate}
\item[(i)] $C_{\mathrm{fix}}(p)=[\Z\xrightarrow{p-1}\Z]$ with
$H^0=\Z/(p-1)$, and $C_{\mathrm{fix}}(p^2)=[\Z\xrightarrow{p^2-1}\Z]$ with
$H^0=\Z/(p^2-1)$. Localized at $p$, both $p-1$ and $p^2-1$ are units in
$\Z_p$, so both complexes are contractible over $\Z_p$ and $v_p$ of both
determinants is $0$.
\item[(ii)] $K_p=[\Z\xrightarrow{p}\Z]$ and $K_{p^2}=[\Z\xrightarrow{p}\Z]$,
since $\Phi_p(1)=\Phi_{p^2}(1)=p$; both have $H^0=\Z/p$ and $v_p$ of both
determinants is $1$.
\item[(iii)] They differ, at $p$ and at $p^2$, already in the $p$-primary part:
$0$ versus $\Z/p$.
\item[(iv)] The cyclotomic factor is the determinant of the deck translation,
not of the correspondence: $\Phi_n(1)=\det(1-T\mid V_n)$, where $T$ is the
canonical order-$n$ automorphism of $L_S/nL_S\cong\Z/n$ given by translation by
the class of $1$. By Lemma~\ref{lem:deformation}, $T$ generates $\mu_n$, the
deck group of the relative fibre, while the fixed contact is governed by
$\mu_{n-1}$, and the two live on the disjoint sets $V(n)$ and $V(n-1)$.
\end{enumerate}
\end{theorem}

\begin{proof}
(i) $p-1$ and $p^2-1$ are prime to $p$, hence units in $\Z_p$, so the
differential is an isomorphism and the complex is contractible. (ii)
$\Phi_p(x)=1+x+\cdots+x^{p-1}$ gives $\Phi_p(1)=p$, and
$\Phi_{p^2}(x)=\Phi_p(x^p)$ gives $\Phi_{p^2}(1)=p$; then \eqref{eq:Kn}. (iii)
compares (i) and (ii). (iv) is Theorem~\ref{thm:fourmodels}(c) together with
Lemma~\ref{lem:deformation} and Theorem~\ref{thm:separation}.
\end{proof}

\begin{theorem}[The same comparison on $C_S$]\label{thm:carrierside}
Let $p\in S_f$, $k\ge1$, $n=p^k$. Then:
\begin{enumerate}
\item[(i)] the fixed-point weight $1/|1-x|_p$ of
Proposition~\ref{prop:shadow} is identically $1$ on the valuation shell
$p^k\Z_p^\times$, because $|1-p^ku|_p=1$ for every $u\in\Z_p^\times$ and
$k\ge1$. So the entire analytic weight of the shell is the half-density
$p^{-k/2}$, and no fixed-point data survive there. Dually, the places where the
fixed factor is nontrivial are the $q$ with $v_q(1-\sigma_q)>0$, and the norm
class never visits them: it is on no finite slice, and its finite local terms
come from shells only.
\item[(ii)] The surviving weight is a Haar volume: the shell $p^k\Z_p^\times$
is exactly the set of generators of the ideal $p^k\Z_p$, i.e. the
$\Z_p^\times$-orbit of $p^k$, and
\[
  \Lambda(p^k)=\log p=\vol_{d^*x}(p^k\Z_p^\times)=\vol_{d^*x}(\Z_p^\times),
\]
independent of $k$ by translation invariance of $d^*x$ on $\Q_p^\times$; the
dependence on $k$ sits entirely in the half-density $p^{-k/2}$.
\item[(iii)] Explicitly,
\[
  \cfin(\Gamma_p^{\mathrm{rad}})=(\log p)p^{-1/2},\qquad
  \cfin(\Gamma_{p^2}^{\mathrm{rad}})=(\log p)p^{-1},
\]
both with primitive degree $\Lambda=\log p$, differing only by the
half-density, whereas the fixed data of the same indices are $|1-p|_q^{-1}$ at
$q\mid p-1$ and $|1-p^2|_q^{-1}$ at $q\mid p^2-1$, numbers attached to the lift
which the class does not see.
\end{enumerate}
\end{theorem}

\begin{proof}
(i) For $k\ge1$, $v_p(p^ku)=k>0$, so $p^ku\in p\Z_p$ and
$1-p^ku\in1+p\Z_p\subset\Z_p^\times$, whence $|1-p^ku|_p=1$; then
Proposition~\ref{prop:trace}(a) and Proposition~\ref{prop:shadow} give rank one
and weight $\tau_p=p^{-k/2}$. The dual statement is
Proposition~\ref{prop:noslice} together with the vanishing of the unit packet
on radial sections established in the proof of Theorem~\ref{thm:radtrace}.
(ii) $p^k\Z_p^\times=\{x:v_p(x)=k\}$ is the coset $p^k\cdot\Z_p^\times$ and
$d^*x$ is translation-invariant with $\vol(\Z_p^\times)=\log p$; the generators
of $p^k\Z_p$ are exactly the elements of valuation $k$. (iii) is
Theorem~\ref{thm:radtrace} at $n=p$ and $n=p^2$, and
Theorem~\ref{thm:separation}(iv) for the lift data.
\end{proof}

\begin{remark}\label{rem:nomapsupplied}
Both $\vol_{d^*x}(p^k\Z_p^\times)=\log p$ and $\Phi_{p^k}(1)=p$ are independent
of $k$, the first by translation invariance of multiplicative Haar measure, the
second because $\Phi_{p^k}(x)=\Phi_p(x^{p^{k-1}})$. These are two unrelated
reasons, and we construct no map between the two objects.  What we prove is
that the two contacts differ already at $p$ and $p^2$
(Theorem~\ref{thm:discriminating}), and that the fixed-point weight is trivial
on the shell while the surviving weight is a Haar volume
(Theorem~\ref{thm:carrierside}).
\end{remark}

\begin{theorem}[Degree matching]\label{thm:degreematching}
Let $n\ge2$ and $S\supseteq\{\infty\}\cup P(n)$. Then
\begin{enumerate}
\item[(i)] $\log[L_S:nL_S]=\degdet Q_{S,n}=\degdet C_{\mathrm{ram}}(n)
=\log n$;
\item[(ii)] $\log n=\sum_{d\mid n}\Lambda(d)
=\sum_{d\mid n}d^{1/2}\cfin(\Gamma_d^{\mathrm{rad}})$;
\item[(iii)] $\degdet\Prim_n=\log\#(\Z/\Phi_n(1))=\log|\det(1-T\mid V_n)|
=\Lambda(n)=n^{1/2}\cfin(\Gamma_n^{\mathrm{rad}})$;
\item[(iv)] the passage from (i)--(ii) to (iii) is M\"obius inversion over the
divisor poset, categorified by the divisor cube.
\end{enumerate}
\end{theorem}

\begin{proof}
(i) is Theorem~\ref{thm:indexcomplex}(i)--(ii) and
Definition~\ref{def:degdet}. (ii): the first equality is
Corollary~\ref{cor:ghost}, the second is \eqref{eq:radcontact}, which gives
$d^{1/2}\cfin(\Gamma_d^{\mathrm{rad}})=\Lambda(d)$ for every $d\ge2$, together
with $\Lambda(1)=0$, the index $1$ being identity or excess and $t=0$ being
excluded from $\omega_{\mathrm{fin}}$ by construction. (iii) is
Lemma~\ref{lem:primstructure}, Theorem~\ref{thm:degreelaw} and
Theorem~\ref{thm:fourmodels}(c) together with
\eqref{eq:radcontact}. (iv) is Lemma~\ref{lem:mobiuscube} and
Theorem~\ref{thm:degreelaw}.
\end{proof}

\begin{remark}\label{rem:whatdegmatchdoesnotsay}
Theorem~\ref{thm:degreematching} does not say that $\Z/\Phi_n(1)$ is an object
of the correspondence algebra, or that a complex over $\Z$ is a shadow of
$\Gamma_n^{\mathrm{rad}}$, and it could not: by Proposition~\ref{prop:shadow} the
shadow of an external class at a finite place $p\mid n$ is the single rank-one
origin class, and by Proposition~\ref{prop:noZp} no additive contact functor
from a $\C$-linear category recovers nonzero integral torsion. The content is
an equality of \emph{numbers}, each independently computed on its own side, and
the reason the equality is not a tautology is that the left-hand sides are
lattice indices and determinants over $\Z$ while the right-hand sides are Haar
volumes of shells in $C_S$.
\end{remark}

\subsection*{Examples}

The following values illustrate Theorems~\ref{thm:separation}
and~\ref{thm:degreelaw}.

\begin{center}
\small
\begin{tabular}{lllllll}
\toprule
$n$ & $P(n)$ & $\Phi_n(1)$ & $\Lambda(n)$ & $\#H^0(C_{\mathrm{fix}}(n))$ &
 $V(n)$ & $V(n-1)$\\
\midrule
$2$  & $\{2\}$     & $2$ & $\log2$ & $1$  & $\{2\}$     & $\emptyset$\\
$4$  & $\{2\}$     & $2$ & $\log2$ & $3$  & $\{2\}$     & $\{3\}$\\
$8$  & $\{2\}$     & $2$ & $\log2$ & $7$  & $\{2\}$     & $\{7\}$\\
$9$  & $\{3\}$     & $3$ & $\log3$ & $8$  & $\{3\}$     & $\{2\}$\\
$6$  & $\{2,3\}$   & $1$ & $0$     & $5$  & $\{2,3\}$   & $\{5\}$\\
$12$ & $\{2,3\}$   & $1$ & $0$     & $11$ & $\{2,3\}$   & $\{11\}$\\
$30$ & $\{2,3,5\}$ & $1$ & $0$     & $29$ & $\{2,3,5\}$ & $\{29\}$\\
\bottomrule
\end{tabular}
\end{center}

\section{The integral shadow of the Frobenius family}\label{sec:shadow}

Proposition~\ref{prop:noZp} rules out a $\C$-linear additive functor producing
$\Z/p$, but not an integral construction whose complexification kills the
torsion while a determinant retains an arithmetic degree.  We construct such a
shadow on the Frobenius family.  It is not a functor on the $\C$-linear
category; it is defined on the monoid of Frobenius objects, using the integral
lattice $L_S=\prod_{p\in S_f}\Z_p$ of the coefficient algebra.

\begin{definition}\label{def:shadow}
For $n\ge1$ with $P(n)\subseteq S_f$, let
$\iota(n)_f=(p^{v_p(n)})_{p\in S_f}$ act on $L_S$ by coordinatewise
multiplication, and put
\begin{equation}\label{eq:shadow}
  \Sh_\Z(\Frob_n):=\Cone\bigl(L_S\xrightarrow{\ \iota(n)_f\ }L_S\bigr)
   =[\,L_S\xrightarrow{\ n\ }L_S\,]=Q_{S,n},
\end{equation}
in degrees $-1,0$.
\end{definition}

\begin{lemma}\label{lem:shwelldef}
$\Sh_\Z(\Frob_n)$ depends only on the measure $\Frob_n$, not on the chosen
idele representative:
\begin{enumerate}
\item[(i)] $\iota(n)_fL_S=nL_S$, so \eqref{eq:shadow} is
$[L_S\xrightarrow{n}L_S]$;
\item[(ii)] unit averaging acts trivially: for every $u\in K_{P(n)}$ --- every
point of the coset over which \eqref{eq:frobcoset} averages --- the idele
$\iota(n)u$ satisfies $(\iota(n)u)_fL_S=nL_S$, because the entries of $u$ are
units of the respective $\Z_p$;
\item[(iii)] gauge invariance: for $\gamma\in\Gamma_S$ the anchor
$b_S(a)=a_f/a_\infty$ satisfies $b_S(\gamma a)=b_S(a)$, and
$b_S(\iota(n)u)=n\cdot(\text{unit})$, so $b_S(\iota(n)u)L_S=nL_S$;
\item[(iv)] the object alone determines $n$
(Theorem~\ref{thm:frobmonoidal}(iv)), so \eqref{eq:shadow} may equivalently be
read as a definition in terms of $\supp\Frob_n$ with no lift at all.
\end{enumerate}
\end{lemma}

\begin{proof}
(i) $\prod_pp^{v_p(n)}\Z_p=\prod_pn\Z_p=nL_S$, since for $p\in S_f$ we have
$n\Z_p=p^{v_p(n)}\Z_p$, the cofactor $n/p^{v_p(n)}$ being a unit of $\Z_p$.
(ii) is the same computation with $u_p\in\Z_p^\times$. (iii)
$b_S(\gamma a)=\gamma a_f/(\gamma a_\infty)=a_f/a_\infty$, and
$b_S(\iota(n)u)=(p^{v_p(n)}u_p)_p$, which generates $nL_S$ by (i)--(ii). (iv)
is Theorem~\ref{thm:frobmonoidal}(iv).
\end{proof}

\begin{theorem}\label{thm:shstructure}
For $m,n\ge1$ with $P(mn)\subseteq S_f$:
\begin{enumerate}
\item[(i)] \emph{(Cohomology.)} $H^{-1}\Sh_\Z(\Frob_n)=0$ and
$H^0\Sh_\Z(\Frob_n)=L_S/nL_S\cong\Z/n$;
\item[(ii)] \emph{(Transitivity triangle.)} The chain map
$(\mathrm{id},n):\Sh_\Z(\Frob_m)\to\Sh_\Z(\Frob_{mn})$ is injective with
cokernel quasi-isomorphic to $\Sh_\Z(\Frob_n)$, giving an exact triangle
\begin{equation}\label{eq:triangle}
  \Sh_\Z(\Frob_m)\to\Sh_\Z(\Frob_{mn})\to\Sh_\Z(\Frob_n)\xrightarrow{+1};
\end{equation}
\item[(iii)] \emph{(Additivity of the degree.)}
$\degdet\Sh_\Z(\Frob_{mn})=\degdet\Sh_\Z(\Frob_m)+\degdet\Sh_\Z(\Frob_n)$,
i.e. $\log mn=\log m+\log n$;
\item[(iv)] \emph{($\Sh_\Z$ is not monoidal.)} \eqref{eq:triangle} is an exact
triangle and not a splitting, and
$\Sh_\Z(\Frob_m)\otimes^L\Sh_\Z(\Frob_n)\not\cong\Sh_\Z(\Frob_{mn})$ in
general: already for $m=n=2$, the complex
$[\Z\xrightarrow{2}\Z]\otimes^L[\Z\xrightarrow{2}\Z]$ has $H^{-1}=\Z/2\ne0$;
\item[(v)] \emph{(M\"obius cofiber.)} The total cofiber of the radical cube
$J\mapsto\Sh_\Z(\Frob_{n/\prod_{p\in J}p})$, $J\subseteq P(n)$, with edges
$(\mathrm{id},q)$, is quasi-isomorphic to
$\Prim_n\simeq[\Z\xrightarrow{\Phi_n(1)}\Z]$, with $H^{-1}=0$ and
$H^0=\Z/\Phi_n(1)$.
\end{enumerate}
\end{theorem}

\begin{proof}
(i) $n\ne0$ and $L_S$ is torsion-free, so multiplication by $n$ is injective
and $H^{-1}=0$; and $L_S/nL_S=\prod_{p\in S_f}\Z_p/n\Z_p\cong
\prod_{p\mid n}\Z/p^{v_p(n)}\cong\Z/n$ by primary decomposition, the factors at
$p\nmid n$ being $0$.

(ii) The square with rows $L_S\xrightarrow{m}L_S$ and
$L_S\xrightarrow{mn}L_S$, left vertical $\mathrm{id}$ and right vertical $n$,
commutes: $n\circ m=mn\circ\mathrm{id}$. Both verticals are injective. The
cokernel complex is $[0\to L_S/nL_S]$, concentrated in degree $0$, which is
quasi-isomorphic to $\Sh_\Z(\Frob_n)$ by (i), the latter being a free
resolution of $L_S/nL_S$. A short exact sequence of complexes gives the exact
triangle.

(iii) By Definition~\ref{def:degdet} and (i),
$\degdet[L_S\xrightarrow{n}L_S]=\log\#H^0=\log n$; alternatively apply the cone
law to \eqref{eq:triangle}.

(iv) The triangle is not split, and the displayed counterexample is the
computation of Theorem~\ref{thm:notmonoidal}: resolving $\Z/2$ by
$[\Z\xrightarrow{2}\Z]$ and tensoring gives
$H_1=\mathrm{Tor}_1^\Z(\Z/2,\Z/2)=\Z/2\ne0$, whereas $\Sh_\Z(\Frob_4)$ has
$H^{-1}=0$ by (i).

(v) By Theorem~\ref{thm:indexcomplex}(i) the diagonal $\Z\to L_S$ induces a
quasi-isomorphism $C_{\mathrm{ram}}(d)\to Q_{S,d}=\Sh_\Z(\Frob_d)$, naturally
on the whole radical cube of Definition~\ref{def:prim}, since both cubes have
the same edge maps $(\mathrm{id},q)$ and the comparison squares commute
degreewise. A quasi-isomorphism of cubes induces a quasi-isomorphism of total
cofibers, so the two total cofibers agree, and the one on the
$C_{\mathrm{ram}}$-side is $\Prim_n\simeq[\Z\xrightarrow{\Phi_n(1)}\Z]$ by
Definition~\ref{def:prim} and Lemma~\ref{lem:primstructure}.
\end{proof}

\begin{theorem}\label{thm:shdegrees}
Let $n\ge2$ with $P(n)\subseteq S_f$. Then
\begin{align}
  \cfin(\Frob_n)&=n^{-1/2}\cdot\degdet\Prim(\Sh_\Z)_n,\label{eq:shdeg1}\\
  \degdet\Sh_\Z(\Frob_n)&=\log n=\sum_{d\mid n}d^{1/2}\cfin(\Frob_d).
   \label{eq:shdeg2}
\end{align}
Both sides of each identity are attached to the \emph{objects} $\Frob_d$: the
left by Corollary~\ref{cor:frobinvariant}, an invariant of the measure, the
right by Definition~\ref{def:shadow} and Lemma~\ref{lem:shwelldef}, an
invariant of the same measure through its action on $L_S$. The identification
is of degrees; an identification at the level of objects, in a $\C$-linear
category, is impossible by Proposition~\ref{prop:noZp} and
Proposition~\ref{prop:shadow}.
\end{theorem}

\begin{proof}
For \eqref{eq:shdeg1}, the right-hand side is
$n^{-1/2}\log\#(\Z/\Phi_n(1))=n^{-1/2}\log|\Phi_n(1)|=n^{-1/2}\Lambda(n)$ by
Theorem~\ref{thm:shstructure}(v), Definition~\ref{def:degdet} and
Theorem~\ref{thm:degreelaw}; the left-hand side is $\Lambda(n)n^{-1/2}$ by
Corollary~\ref{cor:frobinvariant}. For \eqref{eq:shdeg2}, the first equality is
Theorem~\ref{thm:shstructure}(iii) and (i); the second is
$\log n=\sum_{d\mid n}\Lambda(d)$ (Corollary~\ref{cor:ghost}) together with
$d^{1/2}\cfin(\Frob_d)=\Lambda(d)$ for $d\ge2$ and $\Lambda(1)=0$, the index
$1$ being $\Frob_1=\Delta$, treated as identity or excess by the term
$c\cdot F(1)$ with $c=0$ and not by the degree formula.
\end{proof}

\begin{lemma}\label{lem:shfunctor}
For $d\mid e$ let $i_{d,e}:=(\mathrm{id},e/d):\Sh_\Z(\Frob_d)\to
\Sh_\Z(\Frob_e)$ be the chain map of Theorem~\ref{thm:shstructure}(ii). Then
\begin{equation}\label{eq:strictfunctor}
  i_{e,f}\circ i_{d,e}=i_{d,f}\quad\text{for all } d\mid e\mid f,
  \qquad i_{d,d}=\mathrm{id},
\end{equation}
\emph{literally} and not up to homotopy, so $d\mapsto\Sh_\Z(\Frob_d)$ is a
strict functor from the divisor poset to $\Ch(\Z)$. On the \emph{monoid}
$(\N_{\ge1},\times)$, viewed as a one-object category, $\Sh_\Z$ is a function
and not a functor, and it is not monoidal.
\end{lemma}

\begin{proof}
Degreewise: in degree $-1$ the composite is $\mathrm{id}\circ\mathrm{id}
=\mathrm{id}$; in degree $0$ it is $(f/e)\cdot(e/d)=f/d$. Both are the
components of $i_{d,f}$. Identity: $i_{d,d}=(\mathrm{id},1)=\mathrm{id}$. The
last sentence is Theorem~\ref{thm:shstructure}(iv).
\end{proof}

\subsection{The additive extension}

\begin{definition}\label{def:dirichletpolycorr}
Let $\Z[\mathfrak F^+]$ be the free abelian group on the Frobenius monoid
$\{\Frob_n\}_{n\ge1}$, whose elements $\sum_{n\in A}c_n\Frob_n$ with
$c_n\in\Z$ and $A$ finite we call \emph{integral Dirichlet-polynomial
correspondences}. Extend $\Sh_\Z$ additively,
\begin{equation}\label{eq:shadditive}
  \Sh_\Z\Bigl(\sum_{n\in A}c_n\Frob_n\Bigr):=
   \bigoplus_{n\in A}\Sh_\Z(\Frob_n)^{\oplus c_n}
   =\bigoplus_{n\in A}[L_S\xrightarrow{n}L_S]^{\oplus c_n},
\end{equation}
with $c_n<0$ read as a shift, and extend $\cfin$ linearly.
\end{definition}

\begin{theorem}\label{thm:additiveext}
\begin{enumerate}
\item[(i)] \eqref{eq:shadditive} is well defined on $\Z[\mathfrak F^+]$, and
$\degdet$ is additive on it, with
$\degdet\Sh_\Z(\sum_nc_n\Frob_n)=\sum_nc_n\log n$.
\item[(ii)] $\Z[\mathfrak F^+]$ is a commutative ring under the convolution
product induced by $\Frob_m\star\Frob_n=\Frob_{mn}$; cyclic invariance of any
trace on it is therefore automatic.
\item[(iii)] For every $\sum c_n\Frob_n\in\Z[\mathfrak F^+]$ with all $n\ge2$,
\begin{equation}\label{eq:linearcomparison}
  \sum_nc_n\,n^{1/2}\,\cfin(\Frob_n)=\sum_nc_n\,\degdet\Prim(\Sh_\Z)_n,
\end{equation}
both sides being linear functionals on $\Z[\mathfrak F^+]$, the left computed
from the trace on the smeared Frobenii, which are genuine elements of $\DA$,
and the right from the lattice $L_S$.
\item[(iv)] An object-level functor from the $\C$-linear correspondence
category producing these complexes does not exist, by
Proposition~\ref{prop:noZp}; the complex realization of $\Sh_\Z$ is acyclic.
\end{enumerate}
\end{theorem}

\begin{proof}
(i) Additivity of $\degdet$ on direct sums and the shift rule are
Definition~\ref{def:degdet}; the value on a generator is
Theorem~\ref{thm:shstructure}(iii). (ii) $\Z[\mathfrak F^+]$ is the monoid ring
of the commutative monoid $(\N_{\ge1},\times)$, hence commutative, and a trace
on a commutative algebra satisfies $T(ab)=T(ba)$ trivially. (iii) Both sides are
additive in the generator by construction, and on a generator the identity is
\eqref{eq:shdeg1} multiplied by $n^{1/2}$. (iv) is
Proposition~\ref{prop:noZp}; and the complex realization of
$\Sh_\Z(\Frob_n)$, namely $[\C\xrightarrow{n}\C]$, is contractible for $n\ne0$.
\end{proof}

\begin{remark}\label{rem:shadowlimits}
$\Sh_\Z$ is a construction on the objects $\Frob_n$
(Lemma~\ref{lem:shwelldef}), functorial for the divisor arrows
(Lemma~\ref{lem:shfunctor}), whose M\"obius cofiber is
$[\Z\xrightarrow{\Phi_n(1)}\Z]$. It is not a functor on the $\C$-linear
correspondence category, and it is not a shadow in the sense of
Proposition~\ref{prop:shadow}.  Its complexification is acyclic, as
Proposition~\ref{prop:noZp} predicts, so the integral object cannot be
recovered from the $\C$-linear category.  The identification \eqref{eq:shdeg1}--\eqref{eq:shdeg2} is therefore of degrees
only.
\end{remark}

\section{The two Gram forms}\label{sec:gram}

We now compute the Gram matrices of the two families, Frobenius and
unit-averaged, for the pairing $I_A(F,G)=T(F\star G^\dagger)$; they differ by
the archimedean part of the kernel.  Two conventions are used throughout.

\emph{The chart.}  At a fixed finite chart $S$ the finite part
of $T$ runs over $p\in S_f$ only. Whenever a statement below involves a profile
whose support radius exceeds the chart window, the trace meant is the colimit
trace $T_\infty$ of Proposition~\ref{prop:trace}(c), for which the prime sum is
complete at every fixed support radius. Accordingly, throughout this section
and the next two,
\[
  \Wfin:=\sum_{q\ \mathrm{prime},\ k\ge1}(\log q)q^{-k/2}
    \bigl(\delta_{k\log q}+\delta_{-k\log q}\bigr)
\]
is the \emph{colimit} kernel, over all primes; $T$ means $T_\infty$ and
$I_A(\cdot,\cdot)=T_\infty(\cdot\star\cdot^\dagger)$; and at a fixed finite
chart every displayed identity holds unchanged provided the support of the
relevant profile lies inside the chart window, the two traces differing outside
it by the finite sum of the missing prime terms. The one exception, explicitly
exempted, is Theorem~\ref{thm:infgenus}, which is a fixed-chart statement and
says so.  The values computed below do not depend on this choice: for a Gram
entry, $m/n=q^{\pm k}$ forces $q\mid mn$, hence $q\in S_f$ for every chart
containing $P(mn)$.

\emph{The normalization of the level factors.}
Definition~\ref{def:smearedfrob} offers two conventions: the \emph{bare} one,
with unnormalized indicators, and the flat one, divided by
$\prod_r\vol(1+r^N\Z_r)$. This section works in the bare convention, and every
statement says which normalization it is in. The conversion factor is
\begin{equation}\label{eq:VN}
  V_N(n):=\prod_{r\in S_f\smallsetminus P(n)}\vol(1+r^N\Z_r)
   =\prod_{r\in S_f\smallsetminus P(n)}\bigl((r-1)r^{N-1}\bigr)^{-1},
\end{equation}
which decays geometrically in $N$. Consequently a quantity that grows linearly
in $N$ in one convention tends to $0$ in the other, and no rate is asserted
below for the diagonal: what is convention-free, and what is stated, is that
the point-class pairing $I_A(\Frob_n,\Frob_n)$ does not exist.

\subsection{The Frobenius form}

\begin{theorem}\label{thm:gramoff}
Let $m,n\ge2$, $m\ne n$, with $P(mn)\subseteq S_f$, and let
$\Frob_m^{\mathrm{sm}},\Frob_n^{\mathrm{sm}}\in\DA$ be smeared Frobenii with
archimedean profiles $h_m,h_n$ and level $N$. Assume
$\varepsilon<\varepsilon_0(A,S)$, where $A\subset\N_{\ge2}$ is the finite set
of indices in play and $\varepsilon_0(A,S)>0$ is chosen so that for every
$m,n\in A$ the support of the convolved profile $h_m\star h_n^\dagger$ --- an
interval of width at most $2\varepsilon$ around $n/m$ --- contains no point of
the discrete set $\{q^{-k}:q\ \text{prime},\ k\in\Z\}$ other than $n/m$
itself when $n/m$ is such a point. (All primes, not only $q\in S_f$: in this
section $T=T_\infty$, whose prime sum samples the points $q^{-k}$ for every
prime $q$; since the set $\{q^{-k}\}\cap[\tfrac12 n/m,2n/m]$ is finite for
each pair, $\varepsilon_0(A,S)>0$ exists.) Then
\begin{equation}\label{eq:gramoffentry}
  I_A\bigl(\Frob_m^{\mathrm{sm}},\Frob_n^{\mathrm{sm}}\bigr)
   =T\bigl(\Frob^{\mathrm{sm}}_{m/n,\ P(m)\cup P(n)}\bigr),
\end{equation}
a smeared Frobenius at the rational index $r=m/n$ and prime set
$F=P(m)\cup P(n)$. Its finite part is
\begin{equation}\label{eq:gramfinite}
  \cfin(\Frob_{r,F})=
  \begin{cases}
    (\log q)\,q^{-k/2}, & r=q^{\pm k},\ q\ \text{prime},\ k\ge1,\\
    0, & \text{otherwise},
  \end{cases}
\end{equation}
nonzero exactly when $m/n$ is a positive or negative power of a single prime,
and symmetric under $r\mapsto r^{-1}$. The archimedean part of the same entry
carries the arithmetic indicator
\begin{equation}\label{eq:gramarch}
  W_\infty\bigl(\Frob_{r,F}^{\mathrm{sm}}\bigr)
   =\Bigl(\int_0^\infty K_s(t)\,\bigl[g_{mn}(e^{t})+g_{mn}(e^{-t})\bigr]\,dt\Bigr)
    \cdot\prod_{s\in S_f\smallsetminus F}
    \mathbf 1\bigl[\,s^N\mid(n-m)\,\bigr],
\end{equation}
where $g_{mn}:=h_m\star h_n^\dagger$ is supported in a $2\varepsilon$-window
around $n/m$, so that exactly one of the two terms survives: $g_{mn}(e^{-t})$
when $n<m$ and $g_{mn}(e^{t})$ when $n>m$. It
vanishes for every sufficiently large $N$ whenever $S_f\supsetneq F$.
Finally the convolution of the level-$N$ coset indicators contributes the
constant
\begin{equation}\label{eq:gramconst}
  c_N(m,n):=\prod_{s\in S_f\smallsetminus(P(m)\cap P(n))}\vol(1+s^N\Z_s),
\end{equation}
whose index set is the complement of the \emph{intersection} of the two prime
sets, not of their union. In the flat convention of
Definition~\ref{def:smearedfrob} the entry is divided by $V_N(m)V_N(n)$; since
$V_N(m)V_N(n)=c_N(m,n)\cdot\prod_{s\in S_f\smallsetminus F}\vol(1+s^N\Z_s)$,
the constant there is $\prod_{s\in S_f\smallsetminus F}\vol(1+s^N\Z_s)^{-1}$,
which is $1$ exactly when $F=S_f$.
\end{theorem}

\begin{proof}
$I_A(F,G)=T(F\star G^\dagger)$, and at the level of point classes
$\Frob_m\star\Frob_n^\dagger=\Frob_{m/n,\,P(m)\cup P(n)}$ by
Theorem~\ref{thm:invsemigroup}(i). The smeared version is the convolution of
that measure with the composed archimedean profile and the convolved coset
indicators, which is again of the form \eqref{eq:smearedfrob} with $n$ replaced
by $r$, $P$ by $F$, $h$ by $h_m\star h_n^\dagger$, and an extra scalar which
we now compute place by place. Write $U_s:=1+s^N\Z_s$. By
\eqref{eq:smearedfrob}, $\Frob_m^{\mathrm{sm}}$ carries the factor
$\mathbf 1_{m^{-1}U_s}$ at every $s\notin P(m)$ and the constant
$\mathbf 1_{\Z_s^\times}$ at every $s\in P(m)$, where there is no condition;
and $(\Frob_n^{\mathrm{sm}})^\dagger$ carries $\mathbf 1_{nU_s}$ at every
$s\notin P(n)$ --- the involution replaces $x_s$ by $x_s^{-1}$, and
$\mathbf 1_{n^{-1}U_s}(x_s^{-1})=\mathbf 1_{nU_s}(x_s)$ --- and
$\mathbf 1_{\Z_s^\times}$ at every $s\in P(n)$. Convolution on the compact
group $\Z_s^\times$, whose Haar measure is normalized to be a probability
measure, then gives four cases:
\begin{itemize}
\item $s\notin P(m)\cup P(n)$: both factors are indicators of cosets of the
 open subgroup $U_s$, and
 $\mathbf 1_{aU_s}\star\mathbf 1_{bU_s}=\vol(U_s)\mathbf 1_{abU_s}$, so the
 contribution is $\vol(U_s)$ and the surviving indicator is
 $\mathbf 1_{(n/m)U_s}$;
\item $s\in P(n)\smallsetminus P(m)$: here
 $(\mathbf 1_{m^{-1}U_s}\star\mathbf 1_{\Z_s^\times})(x)
  =\int_{m^{-1}U_s}\mathbf 1_{\Z_s^\times}(y^{-1}x)\,dy=\vol(U_s)$ for every
 $x\in\Z_s^\times$, so the factor is $\vol(U_s)\mathbf 1_{\Z_s^\times}$ and
 the contribution is a further $\vol(U_s)$;
\item $s\in P(m)\smallsetminus P(n)$: symmetrically, another $\vol(U_s)$;
\item $s\in P(m)\cap P(n)$:
 $\mathbf 1_{\Z_s^\times}\star\mathbf 1_{\Z_s^\times}
  =\mathbf 1_{\Z_s^\times}$, contributing $1$.
\end{itemize}
Multiplying over $s\in S_f$ gives \eqref{eq:gramconst}; and in every case the
surviving factor is of the shape \eqref{eq:smearedfrob} prescribes for the
index $r$ and the prime set $F=P(m)\cup P(n)$, namely a coset indicator at
$s\notin F$ and no condition at $s\in F$.

For \eqref{eq:gramfinite} repeat the proof of Theorem~\ref{thm:frobcontact}
with $c_rK_F$ in place of $c_nK_{P(n)}$. The support is
\[
  \bigl\{(1/r\ \text{at}\ \infty;\ u_q\in\Z_q^\times\ \text{at}\ q\in F;\
   1/r\ \text{at}\ s\notin F)\bigr\},
\]
with $1/r$ meaning the positive real $n/m$. A shell $q^{m'}\Z_q^\times$ of
$\Iml\,j_q$ meets it in a set of positive Haar measure only if the archimedean
coordinates agree, $q^{-m'}=1/r$, i.e. $r=q^{m'}$; and then $q\in F$, because
$r=q^{m'}$ forces $q\in P(r)\subseteq P(m)\cup P(n)=F$.  Note what excludes the
remaining shells: for the \emph{smeared} class
the level-$N$ coset at $s\notin F$ is \emph{open}, so a shell at $s\notin F$
could in principle meet the support in positive measure; what forbids it is the
archimedean coordinate together with the hypothesis
$\varepsilon<\varepsilon_0$, which forces $s^{-m'}$ to lie within $2\varepsilon$
of $n/m$, hence to equal $n/m$, hence $s\in P(r)\subseteq F$. The same
argument, with $\varepsilon_0$ taken over all primes, disposes of the prime
terms of $T_\infty$ at primes $\ell\notin S_f$ (which sample the test at
archimedean coordinate $\ell^{-k}$, \cite[Thm.~11.6(b)]{Companion}): none of them
lies in the support. On the surviving
shell $\tau_q=q^{-|m'|/2}$ by Proposition~\ref{prop:trace}(a) and the shell
volume is $\log q$, giving \eqref{eq:gramfinite}. The unit packets vanish for
the reason of Theorem~\ref{thm:frobcontact}, since $1\notin
\supp(h_m\star h_n^\dagger)$, which holds because $m\ne n$ and
$\varepsilon<\varepsilon_0$. Symmetry under $r\mapsto r^{-1}$ is
Lemma~\ref{lem:tauinv}, and amounts to the statement that
\eqref{eq:gramfinite} depends on $|m'|$ only.

For \eqref{eq:gramarch} apply the archimedean computation of
Theorem~\ref{thm:frobcontact} with $r$ in place of $n$ and $g_{mn}$ in place
of $h$; that computation used only that the profile is supported in
$x_\infty>0$ away from $1$, and when the support lies in $x_\infty>1$ (that is,
$n>m$) it is $b_+(t)$ rather than $b_+(-t)$ that survives. The numerator of
$r-1=(m-n)/n$ is $m-n$, so the condition $r\in1+s^N\Z_s$ reads
$s^N\mid(n-m)$.
\end{proof}

\begin{theorem}\label{thm:gramdiag}
Let $n\ge2$, $P:=P(n)$, and let $\Delta^{(n),[\varphi_\varepsilon,N]}\in\DA$ be
the smeared idempotent
\[
  \Delta^{(n),[\varphi_\varepsilon,N]}(x)
   :=\varphi_\varepsilon(\log|x_\infty|)\cdot
     \prod_{r\in S_f\smallsetminus P}\mathbf 1_{1+r^N\Z_r}(x_r),
\]
with $\varphi_\varepsilon\in C_c^\infty(\R)$ real, even, supported in
$[-\varepsilon,\varepsilon]$ and of \emph{unit height},
$\varphi_\varepsilon(0)=1$ --- write $\varphi_\varepsilon(t)=\varphi(t/\varepsilon)$
for a fixed even $\varphi\in C_c^\infty([-1,1])$ with $\varphi(0)=1$. Then:
\begin{enumerate}
\item[(i)] no shell contributes: for every $q\in S_f$ and $k\ne0$ the shell
term vanishes, because the archimedean coordinate of $j_q(q^kw)$ is $q^{-k}$
and $\log q^{-k}\notin\supp\varphi_\varepsilon$ for $\varepsilon<\tfrac12\log2$;
\item[(ii)] the unit packets at $p\mid n$ vanish, because
$a_p=\Delta^{(n),[\varphi_\varepsilon,N]}\circ j_p$ is constant on
$\Z_p^\times$, there being no condition at $p$;
\item[(iii)] the unit packets at $r\nmid n$ do not vanish, and are computed
exactly:
\begin{equation}\label{eq:unitpacket}
  (\log r)\int_{\Z_r^\times}\frac{a_r(u)-a_r(1)}{|1-u|_r}\,du
   =-(\log r)\Bigl[\frac{r-2}{r-1}+(N-1)\Bigr],
\end{equation}
so that
\begin{equation}\label{eq:unitpackettotal}
  \sum_{q\in S_f}W_q\bigl(\Delta^{(n),[\varphi_\varepsilon,N]}\bigr)
   =-\sum_{r\in S_f\smallsetminus P}(\log r)
     \Bigl[\frac{r-2}{r-1}+(N-1)\Bigr],
\end{equation}
exactly, for every $N\ge1$: a divergence linear in $N$ at each place
$r\notin P(n)$;
\item[(iv)] the archimedean term diverges logarithmically in $\varepsilon$:
\begin{equation}\label{eq:archdiverges}
  W_\infty\bigl(\Delta^{(n),[\varphi_\varepsilon,N]}\bigr)
   =-\log\tfrac1\varepsilon+O(1).
\end{equation}
\end{enumerate}
Consequently $T_\infty(\Delta^{(n),[\varphi_\varepsilon,N]})$ is finite for each
fixed $(\varepsilon,N)$ and has no limit as $\varepsilon\to0$ or
$N\to\infty$; both divergences are negative.
\end{theorem}

\begin{proof}
(i) $\varphi_\varepsilon$ is supported in $|\log|x_\infty||\le\varepsilon$ and
$|\log q^{-k}|=|k|\log q\ge\log2>2\varepsilon$.

(ii) By Proposition~\ref{prop:slices}, $j_p(u)=(1$ at $\infty$; $u$ at $p$; $1$
elsewhere$)$ for $u\in\Z_p^\times$; the archimedean coordinate is $1$, so
$\varphi_\varepsilon(0)=1$, and the indicators at $r\in S_f\smallsetminus P$
are evaluated at $1\in1+r^N\Z_r$, giving $1$. Since $p\in P$ carries no
condition, $a_p\equiv1$ on $\Z_p^\times$ and the packet vanishes.

(iii) For $r\in S_f\smallsetminus P$ the same computation gives
$a_r(u)=\mathbf 1_{1+r^N\Z_r}(u)$ and $a_r(1)=1$, so
\[
  (\log r)\int_{\Z_r^\times}\frac{a_r(u)-1}{|1-u|_r}\,du
   =-(\log r)\int_{\Z_r^\times\smallsetminus(1+r^N\Z_r)}|1-u|_r^{-1}\,du .
\]
Stratify $\Z_r^\times$ by $j:=v_r(1-u)$. With $du$ the probability Haar measure
of $\Z_r^\times$, the subgroup $1+r^j\Z_r$ has index $(r-1)r^{j-1}$ in
$\Z_r^\times$ for $j\ge1$, so $\vol(1+r^j\Z_r)=((r-1)r^{j-1})^{-1}$ and
\[
  \vol\{v_r(1-u)=j\}=\frac{1}{(r-1)r^{j-1}}-\frac{1}{(r-1)r^{j}}=r^{-j}
   \quad(j\ge1),
\]
\[
  \vol\{v_r(1-u)=0\}=1-\frac{1}{r-1}=\frac{r-2}{r-1},
\]
the two adding to $1$ as they must. On the stratum $j$ the integrand
$|1-u|_r^{-1}$ is $r^j$, constant. Hence
\[
  \int_{\Z_r^\times\smallsetminus(1+r^N\Z_r)}|1-u|_r^{-1}\,du
   =\frac{r-2}{r-1}\cdot1+\sum_{j=1}^{N-1}r^{-j}\cdot r^{j}
   =\frac{r-2}{r-1}+(N-1),
\]
which is \eqref{eq:unitpacket}. For $r=2$ the first term is $0$, correctly,
since $\Z_2^\times=1+2\Z_2$. Summing over $r\in S_f\smallsetminus P$ and adding
the vanishing terms of (i) and (ii) gives \eqref{eq:unitpackettotal}.

(iv) Put $b:=\Delta^{(n),[\varphi_\varepsilon,N]}\circ j_\infty$, so
$b(t)=\varphi_\varepsilon(t)$ in the logarithmic coordinate, because
$j_\infty(x)$ has every finite coordinate $1\in1+r^N\Z_r$. Here
$b(0)=\varphi_\varepsilon(0)=1\ne0$, so the identity class is in the support
and \eqref{eq:archterm} is used with its counterterm:
\[
  W_\infty=(\gamma+\log4\pi)+\int_0^\infty K(t)
    \bigl[b(t)+b(-t)-2e^{-t/2}\bigr]dt .
\]
Split at $t=\varepsilon$. On $(0,\varepsilon)$ we have
$K(t)=\tfrac1{2t}+O(1)$ and, since $\varphi$ is even and smooth with
$\varphi(0)=1$, $\varphi(t/\varepsilon)-e^{-t/2}=O((t/\varepsilon)^2)+O(t)$, so
$\int_0^\varepsilon K(t)\cdot2[\varphi(t/\varepsilon)-e^{-t/2}]dt=O(1)$; indeed
$\int_0^\varepsilon\tfrac1{2t}O(t^2/\varepsilon^2)dt=O(1)$ and
$\int_0^\varepsilon\tfrac1{2t}O(t)dt=O(\varepsilon)$. On $(\varepsilon,\infty)$
we have $b(\pm t)=0$, so the integrand is $-2K(t)e^{-t/2}$, and
$K(t)e^{-t/2}=\tfrac1{2t}+O(1)$ near $0$ gives
$-2\int_\varepsilon^\infty K(t)e^{-t/2}dt=-\log\tfrac1\varepsilon+O(1)$. Adding
the constant $\gamma+\log4\pi$ gives \eqref{eq:archdiverges}.

The last sentence: \eqref{eq:unitpackettotal} is exact and linear in $N$ with
negative slope $-\sum_{r\in S_f\smallsetminus P}\log r$, and
\eqref{eq:archdiverges} tends to $-\infty$; both are negative and neither has a
limit.
\end{proof}

\begin{corollary}\label{cor:gramdiagconv}
Let $F:=\Frob_n^{\mathrm{sm}}$ be the smeared Frobenius with archimedean
profile $h\not\equiv0$ and level $N$, in the bare convention, with
$\varepsilon<\tfrac14\log2$. Put $g:=h\star h^\dagger$ (convolution on
$\R^\times$ for $d^*x_\infty$) and $\psi(t):=g(e^t)$. Then $\psi\in
C_c^\infty([-2\varepsilon,2\varepsilon])$, $\psi(-t)=\overline{\psi(t)}$, and
\[
  \psi(0)=\|h\|^2:=\int_{\R^\times}|h|^2\,d^*x_\infty>0 .
\]
Let $G^{(n)}_{\psi,N}\in\DA$ be
\[
  G^{(n)}_{\psi,N}(x):=\mathbf 1[x_\infty>0]\,\psi(\log x_\infty)\cdot
     \prod_{r\in S_f\smallsetminus P(n)}\mathbf 1_{U_r}(x_r),\qquad
  U_r:=1+r^N\Z_r .
\]
Then
\begin{equation}\label{eq:FFdagger}
  F\star F^\dagger=V_N(n)\cdot G^{(n)}_{\psi,N},
\end{equation}
and
\begin{equation}\label{eq:diagvalue}
  I_A(F,F)=V_N(n)\,\|h\|^2\Bigl[-\sum_{r\in S_f\smallsetminus P(n)}(\log r)
    \Bigl(\tfrac{r-2}{r-1}+N-1\Bigr)-\log\tfrac1\varepsilon
    +O\Bigl(1+\varepsilon^2\,\frac{\|\psi''\|_\infty}{\|h\|^2}\Bigr)\Bigr],
\end{equation}
with an absolute implied constant. If $h$ is a rescaled profile,
$h(x)=c_\varepsilon\,\eta(\log(nx)/\varepsilon)$ for a fixed
$\eta\in C_c^\infty((-1,1))$ and any constants $c_\varepsilon\neq0$, then
$\varepsilon^2\|\psi''\|_\infty/\|h\|^2$ is independent of $\varepsilon$, and
the error term is $O_\eta(1)$. Consequently, for such rescaled profiles:
the bracket is negative for $N$ large or $\varepsilon$ small; in the bare convention $V_N(n)$ decays
geometrically while the bracket grows linearly, so $I_A(F,F)/\|h\|^2\to0$ as
$N\to\infty$ for fixed $\varepsilon$; in the flat convention
$I_A(F^\flat,F^\flat)/\|h\|^2\to-\infty$ geometrically; and in both the
archimedean part $-\|h\|^2\log(1/\varepsilon)$ of the normalized bracket
diverges as $\varepsilon\to0$ for fixed $N$.
\end{corollary}

\begin{proof}
\emph{The profile.} $\supp h\subset(n^{-1}e^{-\varepsilon},n^{-1}e^{\varepsilon})$
gives $\supp g\subset(e^{-2\varepsilon},e^{2\varepsilon})$; $g$ is smooth as
the convolution of two $C_c^\infty$ functions, and it is supported in
$x_\infty>0$ because both factors are. From $F^\dagger(x)=\overline{F(x^{-1})}$,
$g(x^{-1})=\overline{g(x)}$, i.e.\ $\psi(-t)=\overline{\psi(t)}$; and
$g(1)=\int h(y)\overline{h(y)}\,d^*y=\|h\|^2$.

\emph{The convolution \eqref{eq:FFdagger}.} The Haar measure of $C_S$ is the
product of $d^*x_\infty$ and the probability Haar measures of the
$\Z_r^\times$ (Proposition~\ref{prop:splitting}), so $F\star F^\dagger$ is
the product of the coordinatewise convolutions. Archimedean: $g$. At
$r\in S_f\smallsetminus P(n)$:
$\mathbf 1_{n^{-1}U_r}\star\mathbf 1_{nU_r}=\vol(U_r)\mathbf 1_{U_r}$. At
$p\in P(n)$: $\mathbf 1_{\Z_p^\times}\star\mathbf 1_{\Z_p^\times}
=\mathbf 1_{\Z_p^\times}$. The product of the volumes is $V_N(n)$.

\emph{The trace of $G:=G^{(n)}_{\psi,N}$.} Repeat the proof of
Theorem~\ref{thm:gramdiag}, which used about $\varphi_\varepsilon$ only its
support and its value and smoothness near $0$.
(i) Shells: the archimedean coordinate of $j_q(q^kw)$ is $q^{-k}$, and
$|k|\log q\ge\log2>2\varepsilon$, so every shell term vanishes, for every
prime $q$ (also those outside $S_f$ under $T_\infty$).
(ii) At $p\in P(n)$, $a_p:=G\circ j_p\equiv\psi(0)$ on $\Z_p^\times$, so the
packet vanishes.
(iii) At $r\in S_f\smallsetminus P(n)$, $a_r(u)=\psi(0)\mathbf 1_{U_r}(u)$, and
the stratification of Theorem~\ref{thm:gramdiag}(iii), multiplied by
$\psi(0)$, gives the packet
$-\psi(0)(\log r)\bigl[\tfrac{r-2}{r-1}+N-1\bigr]$.
(iv) Archimedean. The profiles of $G$ are $b_+(t)=\psi(t)$ and $b_-\equiv0$
($G$ vanishes on $x_\infty<0$, the paragraph after Proposition~\ref{prop:trace}); the identity class is
in the support, so \eqref{eq:archterm} is used with its counterterm:
\[
  W_\infty(G)=(\gamma+\log4\pi)\psi(0)
   +\int_0^\infty \bigl[2K_s(t)\Re\psi(t)-2K(t)\psi(0)e^{-t/2}\bigr]dt .
\]
Since $\Re\psi$ is even and smooth, $|\Re\psi(t)-\psi(0)|\le\tfrac12
\|\psi''\|_\infty t^2$; and $2K_s(t)=t^{-1}+O(t)$,
$2K(t)e^{-t/2}=t^{-1}+O(1)$ near $0$, so $2K_s-2Ke^{-t/2}$ is bounded on
$(0,1]$. Hence
\[
 \int_0^{2\varepsilon}\bigl[2K_s(t)\Re\psi(t)-2K(t)\psi(0)e^{-t/2}\bigr]dt
 =\int_0^{2\varepsilon}2K_s(\Re\psi-\psi(0))\,dt
  +\psi(0)\int_0^{2\varepsilon}\bigl(2K_s-2Ke^{-t/2}\bigr)dt
 =O\bigl(\varepsilon^2\|\psi''\|_\infty+\varepsilon\,\psi(0)\bigr).
\]
On $(2\varepsilon,\infty)$, $\psi=0$ and
$-2\psi(0)\int_{2\varepsilon}^\infty K(t)e^{-t/2}dt
=-\psi(0)\log\tfrac1\varepsilon+O(\psi(0))$, as in
Theorem~\ref{thm:gramdiag}(iv). Summing,
$W_\infty(G)=-\psi(0)\log\tfrac1\varepsilon
+O(\psi(0)+\varepsilon^2\|\psi''\|_\infty)$.

Adding (i)--(iv), multiplying by $V_N(n)$ and using $\psi(0)=\|h\|^2$ gives
\eqref{eq:diagvalue}. For the rescaled profile, $\psi(t)=c\,\Psi(t/\varepsilon)$
with $\Psi$ fixed (a convolution of $\eta$ with its reflected conjugate) and
$c=|c_\varepsilon|^2\varepsilon$ up to an absolute factor; so $\psi(0)$ and
$\varepsilon^2\|\psi''\|_\infty$ both equal $c$ times quantities depending
only on $\eta$, and their ratio is independent of $\varepsilon$. The remaining
statements follow: the bracket in \eqref{eq:diagvalue} is
$-(\text{linear in }N)-\log(1/\varepsilon)+O(1)$.
\end{proof}

\begin{remark}\label{rem:nodiagonal}
The Frobenius point classes are not elements of $\DA$: they are multipliers
(Remark~\ref{rem:multiplier}, Theorem~\ref{thm:frobincompletion}). There is
therefore no reason for $I_A(\Frob_n,\Frob_n)$ to exist, and
Corollary~\ref{cor:gramdiagconv} says it does not: the regularization
parameters $(\varepsilon,N)$ cannot both be removed. This is the only
convention-free statement available about the diagonal, because by
\eqref{eq:diagvalue} the rate --- indeed the direction of the limit in $N$ ---
depends on the normalization of the level factors. It is also why the point
classes belong to $\mathcal N_A^{\mathrm{dist}}$ and not to $\DA$, and why the
pairing is not extended to $\mathcal N_A^{\mathrm{dist}}\times
\mathcal N_A^{\mathrm{dist}}$.
\end{remark}

\subsection{The unit-averaged form}

\begin{corollary}\label{cor:onlyfinite}
Let $A\subset\N_{\ge2}$ be finite, $G:=\sum_{n\in A}c_n\Frob_n^{\mathrm{sm}}$,
$\varepsilon<\varepsilon_0(A,S)$, and let $N$ be large enough that
$s^N\nmid(n-m)$ for some $s\in S_f\smallsetminus(P(m)\cup P(n))$ for every pair
$m\ne n$ in $A$; this is possible whenever
$S_f\supsetneq\bigcup_{n\in A}P(n)$, since $n-m\ne0$. Then
\begin{equation}\label{eq:gramtotal}
  I_A(G,G)=\sum_{m\ne n}c_mc_n\,V_N^{(m,n)}\,
    \Wfin\bigl(\{\log m-\log n\}\bigr)\,(h_m\star h_n^\dagger)(n/m)
   +\sum_{n\in A}c_n^2\,I_A(\Frob_n^{\mathrm{sm}},\Frob_n^{\mathrm{sm}}),
\end{equation}
with $V_N^{(m,n)}:=c_N(m,n)=\prod_{s\in S_f\smallsetminus(P(m)\cap P(n))}
\vol(1+s^N\Z_s)$ of \eqref{eq:gramconst} --- which for $m=n$ is exactly the
$V_N(n)$ of \eqref{eq:VN}, so that one formula gives both the off-diagonal and
the diagonal constants --- and the diagonal given by
Corollary~\ref{cor:gramdiagconv}, and with \emph{no archimedean off-diagonal
block at all}. The Weil kernel is $\Wfin+\omega_\infty$; the Frobenius family
carries $\Wfin$ and not $\omega_\infty$.
\end{corollary}

\begin{proof}
Theorem~\ref{thm:gramoff} gives the off-diagonal finite part, with its profile
value and convolution constant carried; \eqref{eq:gramarch} with the stated
choice of $N$ gives the vanishing of the off-diagonal archimedean part; and
Corollary~\ref{cor:gramdiagconv} gives the diagonal.
\end{proof}

Two factors must travel with each entry: the atom mass is multiplied by the
value of the convolved profile at the index and by the constant
\eqref{eq:gramconst}. With unit-height profiles one has
$(h_m\star h_n^\dagger)(n/m)=1$; and in the flat convention the constant
becomes $\prod_{s\in S_f\smallsetminus F}\vol(1+s^N\Z_s)^{-1}$, which is $1$
only when $F=S_f$. The entry reduces to the bare atom mass exactly under those
two conditions.

\begin{theorem}\label{thm:fullweil}
Let $\Gamma_n^{\mathrm{rad}}=\Frob_n\star\Delta^{\mathrm{rad}}$
(Proposition~\ref{prop:avgtorad}) and let $\Gamma_{h_m},\Gamma_{h_n}$ be the
smeared radial classes of Definition~\ref{def:smeared}, with $h_m,h_n\in
C_c^\infty(\R)$ supported near $\log m$, $\log n$ and $0\notin
\supp(h_m\star h_n^\dagger)$. Assume, as in Theorem~\ref{thm:gramoff}, that
$\varepsilon<\varepsilon_0(A,S)$, so that $\supp(h_m\star h_n^\dagger)$ meets
$\{q^{-k}\}$ at most in $n/m$; and read $T$ as the colimit trace and $\Wfin$ as
the colimit kernel. Then, for \emph{every} level and with \emph{no indicators
and no arithmetic side conditions},
\begin{equation}\label{eq:radgram}
  I_A\bigl(\Gamma_{h_m},\Gamma_{h_n}\bigr)=T\bigl(\Gamma_{h_m\star h_n^\dagger}
   \bigr)=\int_\R(h_m\star h_n^\dagger)\,d\bigl(\Wfin+\omega_\infty\bigr),
\end{equation}
and in the limit of concentrated profiles the entry is
\begin{equation}\label{eq:radgramentry}
  \Wfin\bigl(\{\log m-\log n\}\bigr)+K\bigl(|\log m-\log n|\bigr):
\end{equation}
the full Weil kernel, both halves, for every pair $m\ne n$.
\end{theorem}

\begin{proof}
$\Gamma_m^{\mathrm{rad}}\star(\Gamma_n^{\mathrm{rad}})^\dagger
=\Gamma_{m/n}^{\mathrm{rad}}$ by Theorem~\ref{thm:radmonoidal}(i)--(ii), the
composition of level-set Haar measures being the multiplication of
$(\R_{>0},\times)$; and on profiles $\Gamma_h\star\Gamma_g^\dagger
=\Gamma_{h\star g^\dagger}$ by Lemma~\ref{lem:smeared}. Then \eqref{eq:radgram}
is Theorem~\ref{thm:radtrace} applied to $h_m\star h_n^\dagger$, i.e. the
identity $T(\Gamma_h)=\int h\,d(\Wfin+\omega_\infty)$ for
$0\notin\supp h$ of Theorem~\ref{thm:repmeasures}. There are no indicators
because a radial class is $\Knat$-invariant: in the notation of
\eqref{eq:smearedfrob} the factors $\prod_r\mathbf 1(x_r)$ are replaced by the
constant $1$ on the whole compact part, so the evaluation at $j_\infty(x)$ ---
whose finite coordinates are $1$ --- is unobstructed, and the arithmetic
indicator of \eqref{eq:frobarch} is identically $1$. Formula
\eqref{eq:radgramentry} is the atom-plus-density reading of \eqref{eq:radgram}
as $h_m\star h_n^\dagger$ concentrates at $\log m-\log n\ne0$.
\end{proof}

\begin{theorem}\label{thm:twohalves}
\begin{enumerate}
\item[(i)] the Frobenius family $\{\Frob_n^{\mathrm{sm}}\}$ has Gram form
$\Wfin$ off the diagonal, with no archimedean block
(Corollary~\ref{cor:onlyfinite}), and a diagonal that diverges as
$-\log(1/\varepsilon)$ archimedeanly and linearly in $N$ at each place
$r\nmid n$ (Theorem~\ref{thm:gramdiag});
\item[(ii)] the unit-averaged family
$\{\Frob_n\star\Delta^{\mathrm{rad}}\}=\{\Gamma_n^{\mathrm{rad}}\}$ has Gram
form the full kernel $\Wfin+\omega_\infty$, exactly and at every level
(Theorem~\ref{thm:fullweil});
\item[(iii)] the mechanism is the complementary-slice dichotomy of
Corollary~\ref{cor:frobnoarch} and Remark~\ref{rem:twoflows}: $\supp\Frob_n$ is
disjoint from $\Iml\,j_\infty$ when $S_f\supsetneq P(n)$, so the Frobenius sees
no archimedean local term; convolving with the projector
$\Delta^{\mathrm{rad}}$ spreads the support onto the whole modulus level set,
which meets $\Iml\,j_\infty$ in the two points $j_\infty(\pm1/n)$ (the
level set is a coset of $\Knat$, and $\Knat\ni-1$), and that is
where $\omega_\infty$ comes from.
\end{enumerate}
Thus the averaging projector $\Delta^{\mathrm{rad}}$ converts the finite part
of the kernel into the whole of it, at the cost, by Lemma~\ref{lem:notideal},
of forgetting the factorization.
\end{theorem}

\begin{proof}
(i) is Corollary~\ref{cor:onlyfinite} and Theorem~\ref{thm:gramdiag}; (ii) is
Theorem~\ref{thm:fullweil}; (iii) is Corollary~\ref{cor:frobnoarch},
Proposition~\ref{prop:avgtorad} and the proof of
Theorem~\ref{thm:fullweil}, where the indicators become trivial precisely
because $\Delta^{\mathrm{rad}}$-invariance removes the finite-coordinate
conditions.
\end{proof}

\subsection{Primitive classes and the finite-place form}

\begin{lemma}\label{lem:complexprofiles}
For a complex $g\in C_c^\infty(\R;\C)$ the class $\Gamma_g$ and the pairing are
defined, and since the kernel $\Wfin+\omega_\infty$ is real and even,
\begin{equation}\label{eq:realimag}
  I_A(\Gamma_g,\Gamma_g)=I_A(\Gamma_{\Rel g},\Gamma_{\Rel g})
   +I_A(\Gamma_{\Iml g},\Gamma_{\Iml g}),
\end{equation}
with $\Gamma_g^\dagger=\Gamma_{g^\dagger}$,
$g^\dagger(t)=\overline{g(-t)}$.
\end{lemma}

\begin{proof}
Put $A:=\widehat{\Rel g}$ and $B:=\widehat{\Iml g}$, so $\widehat g=A+iB$ and
\[
  |\widehat g|^2=|A|^2+|B|^2+2\Iml\bigl(A\overline B\bigr).
\]
Since $\Rel g$ and $\Iml g$ are real, $A(-\tau)=\overline{A(\tau)}$ and
$B(-\tau)=\overline{B(\tau)}$, hence
$(A\overline B)(-\tau)=\overline{(A\overline B)(\tau)}$ and
$\Iml(A\overline B)$ is an \emph{odd} function of $\tau$. The symbol
$\widehat T$ is even on the principal component
(Proposition~\ref{prop:tate}), so the cross term integrates to $0$ in
\eqref{eq:planch}, leaving \eqref{eq:realimag}.
\end{proof}

\begin{definition}\label{def:primitive}
Call $\Gamma_g$ \emph{primitive} when its ruling component vanishes,
\begin{equation}\label{eq:primitive}
  M(\Gamma_g)=\Bigl(\int_\R g(s)e^{s/2}ds,\ \int_\R g(s)e^{-s/2}ds\Bigr)
   =(0,0)\iff\widehat g(\mp i/2)=0 ,
\end{equation}
the two Tate conditions. Write $I_\partial$ for $I_A$ restricted to primitive
radial classes.
\end{definition}

\begin{proposition}\label{prop:formsdiffer}
Let $Q_{\mathrm{fin}}(g):=\int_\R(g\star g^\dagger)\,d\Wfin$ and let $Q$ denote
the full form with kernel $\Wfin+\omega_\infty$. Then $Q-Q_{\mathrm{fin}}$ is
the form $g\mapsto\int|\widehat g(\tau)|^2\widehat\omega_\infty(\tau)\,
d\tau/2\pi$ with
$\widehat\omega_\infty(\tau)=\log\pi-\Rel\psi(1/4+i\tau/2)$, which is
\emph{indefinite}: it is $\approx5.372>0$ at $\tau=0$ and tends to $-\infty$
as $|\tau|\to\infty$. Consequently the two forms are distinct, and their
difference takes both signs.
\end{proposition}

\begin{proof}
The difference of the kernels is $\omega_\infty$ and Plancherel
\eqref{eq:planch} gives the displayed form. For the symbol, by
Proposition~\ref{prop:tate},
$\widehat\omega_\infty(\tau)=\widehat W_\infty(0,\tau)
=\log\pi-\Rel\psi(\tfrac14+\tfrac{i\tau}2)$, whose maximum is
$\log\pi-\psi(1/4)\approx5.372$ at $\tau=0$; and since
$\psi(z)=\log z+O(|z|^{-1})$ as $|z|\to\infty$ with $|\arg z|\le\pi/2$, and
$|\tfrac14+\tfrac{i\tau}2|=\tfrac{|\tau|}2+O(|\tau|^{-1})$,
\begin{equation}\label{eq:symbolasym}
  \widehat\omega_\infty(\tau)=-\log|\tau|+\log2\pi+O(|\tau|^{-1}),
  \qquad|\tau|\to\infty .
\end{equation}
Choosing $\widehat g$ concentrated where $\widehat\omega_\infty>0$ and where
$\widehat\omega_\infty<0$ shows the difference is a nonzero form of both signs.
\end{proof}

\begin{theorem}\label{thm:qfinpositive}
Fix a finite chart $S$ with $S_f\ne\emptyset$ and let
\begin{equation}\label{eq:qfin}
  Q_{\mathrm{fin}}(g)=\int_\R|\widehat g(\tau)|^2\,
    \widehat\omega_{\mathrm{fin}}(\tau)\,\frac{d\tau}{2\pi},\qquad
  \widehat\omega_{\mathrm{fin}}(\tau):=\sum_{p\in S_f}
    \widehat W_p(\mathfrak c_0,\tau),
\end{equation}
be the finite-place form on real radial profiles at that chart.  Then there
exists a \emph{primitive} real $g\in C_c^\infty(\R)$ with
$Q_{\mathrm{fin}}(g)>0$.  In particular $Q_{\mathrm{fin}}$ is not negative
semidefinite on primitive radial classes, at any finite chart.
\end{theorem}

\begin{proof}
\emph{Step 1: the symbol is strictly positive at $\tau=0$.} By
\eqref{eq:symbol} and \eqref{eq:nophase}, on the principal component
$\zeta_p=p^{-1/2+i\tau}$, so at $\tau=0$ we have $\zeta_p=p^{-1/2}\in(0,1)$
real and
\begin{equation}\label{eq:omegafin0}
  \widehat\omega_{\mathrm{fin}}(0)
   =\sum_{p\in S_f}(\log p)\frac{2p^{-1/2}}{1-p^{-1/2}}>0,
\end{equation}
strictly positive, each summand being positive and $S_f\ne\emptyset$. Moreover
$\widehat\omega_{\mathrm{fin}}$ is real-analytic and bounded,
\begin{equation}\label{eq:AS}
  |\widehat\omega_{\mathrm{fin}}|\le A_S
   :=\sum_{p\in S_f}(\log p)\frac{2p^{-1/2}}{1-p^{-1/2}}<\infty
\end{equation}
uniformly in $\tau$. Indeed, for $\zeta=\rho e^{i\theta}$ with $0<\rho<1$ one
computes
\[
  2\Rel\frac{\zeta}{1-\zeta}
   =\frac{2\rho\cos\theta-2\rho^2}{1-2\rho\cos\theta+\rho^2},
\]
whose derivative with respect to $\cos\theta$ is
$2\rho(1-\rho^2)/|1-\zeta|^4>0$; so the expression is maximal at $\theta=0$,
with value $2\rho/(1-\rho)$, and minimal at $\theta=\pi$, with value
$-2\rho/(1+\rho)$. With $\rho=p^{-1/2}$ and $\theta=\tau\log p$ by
\eqref{eq:nophase} this gives \eqref{eq:AS}, and moreover
$A_S=\widehat\omega_{\mathrm{fin}}(0)$. By continuity choose $\delta>0$ with
$\widehat\omega_{\mathrm{fin}}(\tau)\ge\tfrac12\widehat\omega_{\mathrm{fin}}(0)
=:c>0$ for $|\tau|\le\delta$.

\emph{Step 2: a primitive test.} Let $\varphi\in C_c^\infty(\R)$ be real, even,
$\varphi\not\equiv0$, and for $T>0$ set $g_{0,T}(t):=\varphi(t/T)$ and
\[
  g_T:=\Pi g_{0,T}=-g_{0,T}''+\tfrac14g_{0,T}\in C_c^\infty(\R),
\]
real, with $\supp g_T=\supp g_{0,T}$, the operator $\Pi$ being local. Then
$g_T$ is primitive:
\[
  \lambda_\mp(\Gamma_{g_T})
   =\int_\R\bigl(-g_{0,T}''+\tfrac14g_{0,T}\bigr)e^{\pm s/2}\,ds
   =-\tfrac14\int g_{0,T}e^{\pm s/2}ds+\tfrac14\int g_{0,T}e^{\pm s/2}ds=0,
\]
integrating by parts twice with no boundary terms and using
$\partial_s^2e^{\pm s/2}=\tfrac14e^{\pm s/2}$. In Fourier terms
$\widehat g_T(\tau)=(\tau^2+\tfrac14)\widehat g_{0,T}(\tau)
=(\tau^2+\tfrac14)T\widehat\varphi(T\tau)$, and the vanishing at
$\tau=\mp i/2$ is visible as the zero of $\tau^2+\tfrac14$.

\emph{Step 3: concentration.} Write
$M_{\mathrm{bulk}}(T):=\int_{|\tau|\le\delta}|\widehat g_T|^2\,d\tau/2\pi$ and
$M_{\mathrm{tail}}(T):=\int_{|\tau|>\delta}|\widehat g_T|^2\,d\tau/2\pi$.
Substituting $\sigma=T\tau$,
\[
  \int_{|\tau|>\delta}|\widehat g_T(\tau)|^2\frac{d\tau}{2\pi}
   =T\int_{|\sigma|>T\delta}\Bigl(\frac{\sigma^2}{T^2}+\frac14\Bigr)^2
     |\widehat\varphi(\sigma)|^2\frac{d\sigma}{2\pi},
\]
and since $\widehat\varphi$ is Schwartz and
$(\sigma^2/T^2+1/4)^2\le(\sigma^2+1)^2$ for $T\ge1$, the integral over
$|\sigma|>T\delta$ tends to $0$ faster than any power of $T$; so
$M_{\mathrm{tail}}(T)=o(T^{-k})\cdot T$ for every $k$. The same substitution
over $|\sigma|\le T\delta$ gives
$M_{\mathrm{bulk}}(T)=T\bigl(\tfrac1{16}\|\widehat\varphi\|_{L^2}^2/2\pi
+o(1)\bigr)\asymp T$. Hence
$M_{\mathrm{tail}}(T)/M_{\mathrm{bulk}}(T)\to0$.

\emph{Step 4: conclusion.} Using $\widehat\omega_{\mathrm{fin}}\ge c$ on
$|\tau|\le\delta$ and $|\widehat\omega_{\mathrm{fin}}|\le A_S$ everywhere,
\[
  Q_{\mathrm{fin}}(g_T)\ge c\,M_{\mathrm{bulk}}(T)-A_S\,M_{\mathrm{tail}}(T)
   =M_{\mathrm{bulk}}(T)\Bigl(c-A_S\frac{M_{\mathrm{tail}}(T)}
     {M_{\mathrm{bulk}}(T)}\Bigr),
\]
which is $>0$ for all $T$ large enough. Taking such a $T$ gives the theorem.
\end{proof}

\begin{remark}\label{rem:qfinscope}
Theorem~\ref{thm:qfinpositive} is a fixed-chart statement and is not asserted
in the colimit. The proof uses Plancherel against the \emph{bounded} symbol
\eqref{eq:AS} of a finite chart. In the colimit
$\widehat\omega_{\mathrm{fin}}^\infty=\sum_{\text{all }p}\widehat W_p$ is not a
bounded function --- indeed
$\widehat\omega_{\mathrm{fin}}^\infty(0)=\sum_p(\log p)2p^{-1/2}/(1-p^{-1/2})
=+\infty$ --- so the Plancherel step does not transfer, and the two scales in
the argument (the concentration width $\asymp1/T$ of $\widehat g_T$ and the
variation scale $\asymp1/\log p_{\max}$ of the symbol) become comparable.
\end{remark}

\section{The self-intersection of the diagonal}\label{sec:diagonal}

\begin{theorem}\label{thm:infgenus}
Fix a finite chart $S$ with $S_f\ne\emptyset$ and read $T=T_S$ (this is the
exception to the convention of Section~\ref{sec:gram}). Fix a real even $\varphi\in C_c^\infty(\R)$,
$\varphi\not\equiv0$, and for $\varepsilon>0$ let $F_\varepsilon\in\DA$ be the
$\Knat$-invariant test with archimedean profile
$\varepsilon^{-1}\varphi(t/\varepsilon)$ in the coordinate $t=\log|x|_S$, that
is, $F_\varepsilon=\Gamma_h$ with $h(t)=\varepsilon^{-1}\varphi(-t/\varepsilon)$.
Then $F_\varepsilon\star F_\varepsilon^\dagger\to\Delta$ weak-$*$, and
\begin{equation}\label{eq:diagdiv}
  I_A(F_\varepsilon,F_\varepsilon)=T(F_\varepsilon\star F_\varepsilon^\dagger)
   =-\frac{\|\varphi\|_{L^2}^2}{\varepsilon}
     \Bigl(\log\frac1\varepsilon+O(1)\Bigr)\longrightarrow-\infty
   \qquad(\varepsilon\to0^+).
\end{equation}
Equivalently, in normalization-free form: if
$G_\varepsilon:=\|F_\varepsilon\|_{L^2}^{-1}F_\varepsilon$, so that
$\int_{\widehat{C_S}}|\widehat G_\varepsilon|^2d\chi=1$, then
\begin{equation}\label{eq:diagdivnorm}
  I_A(G_\varepsilon,G_\varepsilon)=-\log\frac1\varepsilon+O(1).
\end{equation}
The proof gives a bounded remainder in \eqref{eq:diagdivnorm}, not a limiting
constant; the latter would need a finer expansion of the archimedean symbol
than \eqref{eq:symbolasym} and control of $\sum_p\widehat W_p$ as
$|\tau|\to\infty$.
\end{theorem}

\begin{proof}
$F_\varepsilon\in\DA$ by Lemma~\ref{lem:smeared}, and
$\widehat F_\varepsilon$ is supported on the principal component with
$\widehat F_\varepsilon(\tau)=\widehat\varphi(\varepsilon\tau)$, by the
computation of the radial Fourier transform in the proof of
Theorem~\ref{thm:radincompletion}(ii) together with the scaling
$\widehat{\varepsilon^{-1}\varphi(\cdot/\varepsilon)}(\tau)
=\widehat\varphi(\varepsilon\tau)$. Weak-$*$ convergence
$F_\varepsilon\star F_\varepsilon^\dagger\to\delta_1$ is
$\widehat{F_\varepsilon\star F_\varepsilon^\dagger}
=|\widehat\varphi(\varepsilon\tau)|^2\to|\widehat\varphi(0)|^2$; normalizing
$\widehat\varphi(0)=1$, possible after scaling since $\varphi\not\equiv0$ is
real even, makes $F_\varepsilon\star F_\varepsilon^\dagger$ an approximate
identity.

By \eqref{eq:planch} with $c=0$,
\[
  I_A(F_\varepsilon,F_\varepsilon)
   =\int_{\widehat{C_S}}|\widehat F_\varepsilon(\chi)|^2\widehat T(\chi)\,d\chi
   =\int_\R|\widehat\varphi(\varepsilon\tau)|^2\,
     \widehat T(\mathfrak c_0,\tau)\,\frac{d\tau}{2\pi}.
\]

\emph{The symbol.} The asymptotic \eqref{eq:symbolasym} holds as
$|\tau|\to\infty$, and is used below after the substitution
$\sigma=\varepsilon\tau$, where the region $|\sigma|\lesssim\varepsilon$
corresponds to \emph{bounded} $\tau$, in which it is false. The bounded region
is handled separately: $\widehat W_\infty(0,\cdot)$ is continuous on $\R$,
bounded above by its value $\log\pi-\psi(1/4)$ at $\tau=0$ and bounded below on
any compact set, so for every $R_0>0$
\begin{equation}\label{eq:boundedregion}
  \Bigl|\int_{|\tau|\le R_0}|\widehat\varphi(\varepsilon\tau)|^2
    \widehat T(\mathfrak c_0,\tau)\frac{d\tau}{2\pi}\Bigr|
   \le C(R_0)\|\widehat\varphi\|_\infty^2=O(1)
\end{equation}
uniformly in $\varepsilon$, so the $|\sigma|\le\varepsilon R_0$ part of the
scaled integral contributes $\varepsilon^{-1}\cdot O(\varepsilon)=O(1)$ and does
not affect \eqref{eq:diagdiv}. At the finite places, by \eqref{eq:symbol} and
\eqref{eq:AS}, $|\sum_{p\in S_f}\widehat W_p(\mathfrak c_0,\tau)|\le A_S$
uniformly in $\tau$. Hence
$\widehat T(\mathfrak c_0,\tau)=-\log|\tau|+O(1)$ with an $O(1)$ bounded by
$\log2\pi+A_S+o(1)$.

\emph{The scaling.} Substitute $\sigma=\varepsilon\tau$, using
\eqref{eq:symbolasym} on $|\tau|>R_0$ and \eqref{eq:boundedregion} on
$|\tau|\le R_0$:
\[
  I_A(F_\varepsilon,F_\varepsilon)=\frac1\varepsilon\int_\R
   |\widehat\varphi(\sigma)|^2\bigl[\log\varepsilon-\log|\sigma|+O(1)\bigr]
   \frac{d\sigma}{2\pi}.
\]
The integral $\int|\widehat\varphi(\sigma)|^2\log|\sigma|\,d\sigma$ converges
absolutely, $\widehat\varphi$ being Schwartz and $\log|\sigma|$ locally
integrable, and $\int|\widehat\varphi|^2d\sigma/2\pi=\|\varphi\|_{L^2}^2$ by
Plancherel. Therefore
$I_A(F_\varepsilon,F_\varepsilon)=\varepsilon^{-1}\|\varphi\|_2^2
(\log\varepsilon+O(1))$, which is \eqref{eq:diagdiv}; dividing by
$\|F_\varepsilon\|_{L^2}^2=\varepsilon^{-1}\|\varphi\|_2^2$ gives
\eqref{eq:diagdivnorm}.

Finally, the conclusion is unchanged, and easier, in the colimit: for
$\varepsilon<\tfrac12\log2$ the profile of
$F_\varepsilon\star F_\varepsilon^\dagger$ is supported in
$|t|\le2\varepsilon<\log2$, so it meets \emph{no} atom of $\Wfin$, whose
support consists of points of modulus at least $\log2$ by
Theorem~\ref{thm:contactlocus}(i). The finite part is therefore identically
zero and \eqref{eq:diagdiv} is purely archimedean, no appeal to the bound
\eqref{eq:AS} being needed.  So the theorem also holds on the colimit.
\end{proof}

\begin{remark}\label{rem:infinitegenus}
For a smooth projective curve of genus $g$ over a finite field one has
$\Delta\cdot\Delta=2-2g$ \cite{Weil1948}. Theorem~\ref{thm:infgenus} says that
the regularized self-intersection of the diagonal on $C_S$ diverges to
$-\infty$; under that dictionary this reads $g=\infty$. The
divergence comes entirely from the archimedean symbol
\eqref{eq:symbolasym}, the finite-place symbol being bounded at a fixed chart
by \eqref{eq:AS}, and the rate is $\log(1/\varepsilon)$ per unit of
$L^2$-normalized test.  The interpretation is not used in any proof.  Note also
that the coefficient of $\log(1/\varepsilon)$ is $1$ in the
$L^2$-normalization, and that in the mass-normalized normalization
($\int F_\varepsilon=1$) the divergence is
$-\varepsilon^{-1}\|\varphi\|_2^2\log(1/\varepsilon)$, one power of
$\varepsilon^{-1}$ faster; the rate is normalization-dependent and the
qualitative statement is not.
\end{remark}

\section{The normalization of the shell volumes}
\label{sec:normalization}

The identification $\Lambda(p^k)=\log p=\vol_{d^*x}(p^k\Z_p^\times)$ depends on
the normalization \eqref{eq:shellvol}.  That normalization is determined by
the local Tate residues.

\begin{theorem}\label{thm:tateresidue}
Let the local multiplicative measures be normalized by
\[
  \vol_{d^*x_p}(\Z_p^\times)=c_p>0\ (p\ \text{finite}),\qquad
  d^*x_\infty=dx/(2|x|),
\]
all other conventions unchanged. Define the \emph{local Tate residue} of a
place $v$ as the residue at $s=0$ of the local zeta integral
$\zeta_v(s):=\int_{|x|_v\le1}|x|_v^s\,d^*x_v$. Then:
\begin{enumerate}
\item[(i)] $\zeta_p(s)=c_p/(1-p^{-s})$, with a simple pole at $s=0$ of residue
$c_p/\log p$;
\item[(ii)] $\zeta_\infty(s)=1/s$, of residue $1$;
\item[(iii)] hence the local Tate residues agree at all places if and only if
$c_p=\log p$ for every finite $p$.
\end{enumerate}
\end{theorem}

\begin{proof}
(i) $\{|x|_p\le1\}=\bigsqcup_{k\ge0}p^k\Z_p^\times$, and $|x|_p^s=p^{-ks}$ on
the $k$-th shell, whose $d^*x_p$-volume is $c_p$ by translation invariance. So
$\zeta_p(s)=c_p\sum_{k\ge0}p^{-ks}=c_p/(1-p^{-s})$ for $\Rel s>0$,
meromorphically continued, and $1-p^{-s}=s\log p+O(s^2)$ gives residue
$c_p/\log p$.

(ii) $\int_{|x|\le1}|x|^s\,dx/(2|x|)=2\cdot\tfrac12\int_0^1x^{s-1}dx=1/s$.

(iii) is (i) together with (ii).
\end{proof}

\begin{corollary}\label{cor:uniquenorm}
Under the requirement of a \emph{universal Tate residue} --- one and the same
residue at every place of $S$ --- the normalization
$\vol_{d^*x}(\Z_p^\times)=\log p$ is unique, and
Corollary~\ref{cor:frobinvariant} then reads
\[
  \Lambda(p^k)=\log p=\vol_{d^*x}\bigl(\supp\Frob_{p^k}\bigr):
\]
the Haar volume of the support of the Frobenius object, in the unique local
normalization compatible with the Tate boundary. The only remaining freedom is
an overall positive scale on $T$, which rescales the pairing and changes no
sign, no vanishing and no ratio.
\end{corollary}

\begin{proof}
Theorem~\ref{thm:tateresidue}(iii). Rescaling all local measures by one
$\lambda>0$ gives $c_p=\lambda\log p$ and $\lambda\cdot dx/(2|x|)$ at $\infty$,
so the normalization is unique in the quotient by the global scale.
\end{proof}

\begin{remark}\label{rem:notchartinv}
Chart invariance alone does not force $c_p=\log p$. By Theorem~\ref{thm:tateresidue}(i)--(ii) each place
of $S$ carries one local Tate residue --- $c_p/\log p$ at a finite $p$ and $1$
at $\infty$ --- and there are $|S_f|+1$ of them; that count changes when a
place is added to the chart, so the collection of local residues is manifestly
chart-dependent, whatever the $c_p$. The colimit trace is chart-invariant on
tests of small support radius, but a residue is read off the local zeta
integral, not off values on a compact window.  What Theorem~\ref{thm:tateresidue}
shows is that, once one requires every place to carry the same Tate residue,
the normalization is unique; this is the classical normalization of the local
multiplicative measures \cite{Weil1952}, \cite{Tate}.
\end{remark}

\section{Approximation lemmas and a single-frequency computation}
\label{sec:approx}

This section is analytic and contains three independent results: a continuity bound for the
colimit trace on a moving window; a smooth partition-of-unity approximation of
a function on $[L/2,L]$ by a Dirichlet-polynomial profile with per-$n$ widths,
together with the $C^1$ control of the approximant; and an exact computation of
the contribution of a resonant quadruple of zeros to a Hermitian zero-sum
against a single-frequency test.

\begin{lemma}\label{lem:traceontinuity}
Let $R\ge1$ and let $h\in C_c^\infty(\R)$ be supported in $[-R,R]$. Then
\begin{equation}\label{eq:tracebound}
  |T_\infty(\Gamma_h)|\le C_0\|h\|_\infty e^{R/2}+C_1\|h\|_{C^1},
\end{equation}
with $C_0,C_1$ absolute constants. In particular
$h\mapsto T_\infty(\Gamma_h)$ is bounded by $O(e^{R/2})$ on
$\{h:\supp h\subseteq[-R,R],\ \|h\|_{C^1}\le1\}$.
\end{lemma}

\begin{proof}
Here $0$ may lie in $\supp h$, so Theorem~\ref{thm:radtrace} does not apply as
stated; what is used instead is the \emph{pointwise} vanishing of the unit
packets on radial sections established inside its proof, which needs no
hypothesis on $\supp h$, together with the local terms
\eqref{eq:localterms}--\eqref{eq:archterm} read on the colimit. This gives
$T_\infty(\Gamma_h)=\int h\,d\Wfin+\int h\,\omega_\infty+(\text{identity
term})$, all three finite. The finite part is
$\sum_k\Lambda(k)k^{-1/2}[h(\log k)+h(-\log k)]$, a sum over $k\le e^R$ bounded
by $2\|h\|_\infty\sum_{k\le e^R}\Lambda(k)k^{-1/2}$. By Chebyshev's bound
$\psi(X):=\sum_{k\le X}\Lambda(k)\le cX$ with an absolute $c$
\cite[Ch.~2]{IK}, partial summation gives
$\sum_{k\le X}\Lambda(k)k^{-1/2}\le C_0'X^{1/2}$ with $C_0'$ absolute, whence
the first term of \eqref{eq:tracebound}.  The archimedean part is
\eqref{eq:archterm}: away from $t=0$ the density $K$ is integrable against a
bounded compactly supported function; near $t=0$ one has $K(t)=1/(2t)+O(1)$ and
the bracket $h(t)+h(-t)-2h(0)e^{-t/2}$ is $O(t\|h\|_{C^1})$, so that term is
$O(\|h\|_{C^1})$.
\end{proof}

\begin{lemma}[Discrete approximation on a moving window]
\label{lem:discreteapprox}
Fix once and for all a real $\rho\in C_c^\infty([-1,1])$ with
\begin{equation}\label{eq:partition}
  \sum_{k\in\Z}\rho(x-k)\equiv1\qquad\text{for all }x\in\R,
\end{equation}
a smooth partition of unity subordinate to the unit grid; such a $\rho$ exists,
and $\int\rho=1$ follows by integrating \eqref{eq:partition} over a period. Let
$L\ge4$ and let $g\in C^1(\R)$ be supported in $[L/2,L]$. For $n$ with
$\log n\in[L/2,L]$ put $\mathrm{gap}_n:=\log(n+1)-\log n=\log(1+1/n)$, so that
$e^{-L}\le\mathrm{gap}_n\le2e^{-L/2}$, and set
\begin{equation}\label{eq:gD}
  g_D(t):=\sum_{\log n\in[L/2,L]}g(\log n)\,
    \rho\Bigl(\frac{t-\log n}{\mathrm{gap}_n}\Bigr).
\end{equation}
Then $g_D$ is a Dirichlet-polynomial profile with per-$n$ widths
$\varepsilon_n=\mathrm{gap}_n$, and
\begin{equation}\label{eq:gDapprox}
  \|g-g_D\|_\infty\le Ce^{-L/2}\|g\|_{C^1},\qquad
  \supp(g-g_D)\subseteq[L/2-1,L+1],
\end{equation}
with $C$ absolute.
\end{lemma}

\begin{proof}
\emph{Step 1: the mesh sums to $1+O(e^{-L/2})$, pointwise.} Fix
$t\in[L/2,L]$ and put $x:=e^t\in[e^{L/2},e^L]$. For $n$ with
$|t-\log n|\le\mathrm{gap}_n$ we have $|x/n-1|\le e^{\mathrm{gap}_n}-1\le2/n$,
so $|x-n|\le2$. For such $n$,
\[
  t-\log n=\log(x/n)=\frac{x-n}{n}\Bigl(1+O\bigl(\tfrac{|x-n|}{n}\bigr)\Bigr),
  \qquad
  \mathrm{gap}_n=\log(1+1/n)=\frac1n\bigl(1+O(1/n)\bigr),
\]
hence
\begin{equation}\label{eq:gapratio}
  \frac{t-\log n}{\mathrm{gap}_n}=(x-n)\bigl(1+O(1/n)\bigr).
\end{equation}
Therefore, with $K:=\|\rho'\|_\infty$ and using \eqref{eq:partition} on the
uniform grid,
\[
  \Bigl|\sum_n\rho\Bigl(\frac{t-\log n}{\mathrm{gap}_n}\Bigr)
   -\sum_n\rho(x-n)\Bigr|
   \le K\sum_{|x-n|\le2}|x-n|\cdot O(1/n)=O(K/n)=O\bigl(Ke^{-L/2}\bigr),
\]
because at most five integers $n$ satisfy $|x-n|\le2$ and each has
$n\ge e^{L/2}-2$. With $\sum_n\rho(x-n)=1$ this reads
\begin{equation}\label{eq:meshsum}
  \sum_{\log n\in[L/2,L]}\rho\Bigl(\frac{t-\log n}{\mathrm{gap}_n}\Bigr)
   =1+O(e^{-L/2})
\end{equation}
for $t$ in the interior of the window, at distance at least $2e^{-L/2}$ from
its endpoints; near the two endpoints the sum lies in $[0,1]$ and $g$ vanishes
there up to $O(e^{-L/2}\|g\|_{C^1})$, which is absorbed. This is a
\emph{pointwise} statement, and it is what the lemma needs; no tested estimate
is used.

\emph{Step 2: the approximation.} By \eqref{eq:meshsum},
\[
  g_D(t)-g(t)=\sum_n\rho\Bigl(\frac{t-\log n}{\mathrm{gap}_n}\Bigr)
   \bigl[g(\log n)-g(t)\bigr]+g(t)\cdot O(e^{-L/2}),
\]
and on the support of the $n$-th term
$|t-\log n|\le\mathrm{gap}_n\le2e^{-L/2}$, so
$|g(\log n)-g(t)|\le2e^{-L/2}\|g'\|_\infty$. Since at most five terms are
nonzero at any $t$ and $0\le\rho\le\|\rho\|_\infty$, the first sum is
$O(e^{-L/2}\|g\|_{C^1})$ and the second is $O(e^{-L/2}\|g\|_\infty)$. The
support statement is immediate from $\mathrm{gap}_n\le1$.
\end{proof}

\begin{lemma}[$C^1$ control of the approximant]\label{lem:C1control}
In the notation of Lemma~\ref{lem:discreteapprox},
\begin{equation}\label{eq:C1bound}
  \|(g_D)'\|_\infty\le C\|g\|_{C^1},\qquad\text{hence}\qquad
  \|(g-g_D)'\|_\infty\le C\|g\|_{C^1},
\end{equation}
with $C$ depending only on $\rho$. The bound is $O(1)$, not $o(1)$.
\end{lemma}

\begin{proof}
Write $a_n:=(t-\log n)/\mathrm{gap}_n$, so
$(g_D)'(t)=\sum_ng(\log n)\rho'(a_n)/\mathrm{gap}_n$. Split
$g(\log n)=g(t)+[g(\log n)-g(t)]$.

\emph{The coefficient-variation part.} On the support of the $n$-th term
$|t-\log n|\le\mathrm{gap}_n$, so
$|g(\log n)-g(t)|\le\|g'\|_\infty\mathrm{gap}_n$, and the term is bounded by
$\|g'\|_\infty\|\rho'\|_\infty$; at most five $n$ contribute, so this part is
$O(\|g'\|_\infty)$.

\emph{The mesh part, where the cancellation lives.} It is $g(t)S(t)$ with
$S(t):=\sum_n\rho'(a_n)/\mathrm{gap}_n$. Put $x:=e^t$. By
\eqref{eq:gapratio}, $a_n=(x-n)(1+\theta_n)$ with $|\theta_n|\le C/n$, and only
$n$ with $|x-n|\le2$ contribute. Taylor expansion gives
\[
  \sum_n\rho'(a_n)=\sum_n\rho'(x-n)+\sum_n\rho''(\xi_n)(x-n)\theta_n
   =0+O\bigl(\|\rho''\|_\infty\cdot C/n\bigr)=O(1/x),
\]
where $\sum_n\rho'(x-n)\equiv0$ because it is the derivative of the partition
of unity \eqref{eq:partition}, and at most five terms occur. This cancellation
is the whole point: without it the naive bound would be
$\|\rho'\|_\infty/\mathrm{gap}_n\asymp e^L$. Finally
$1/\mathrm{gap}_n=n(1+O(1/n))=x+O(1)$, so
\[
  S(t)=\sum_n\rho'(a_n)\bigl(x+O(1)\bigr)
   =x\cdot O(1/x)+O(\|\rho'\|_\infty)=O(1).
\]
Adding the two parts gives \eqref{eq:C1bound}.
\end{proof}

We close with an exact computation.  It is the only statement in this paper
about the zeros of the zeta function.

\begin{lemma}[The single-frequency test]\label{lem:singlefreq}
For $g\in C_c^\infty(\R)$ write
\begin{equation}\label{eq:zerosum}
  \widehat g(s):=\int_\R g(t)\,e^{(s-\frac12)t}\,dt,\qquad
  Q(g):=\sum_\rho\widehat g(\rho)\,\overline{\widehat g(1-\bar\rho)},
\end{equation}
the sum over the nontrivial zeros $\rho$ of the zeta function, with
multiplicity, which is
Hermitian after the involution $\rho\mapsto1-\bar\rho$ of the zero multiset.
Let $\psi\in C_c^\infty([1/2,1])$ with $\psi\ge0$, $\psi\not\equiv0$, let
$\gamma_0\in\R$, and put
\begin{equation}\label{eq:singlefreqtest}
  g_L(t):=\psi(t/L)\,e^{-i\gamma_0t}.
\end{equation}
Assume, for this lemma only, that $\rho_0=\beta+i\gamma_0$ is a zero with
$0<\delta:=\beta-\tfrac12<\tfrac12$; parts (ii)--(v) are vacuous if no such
zero exists. Write $a:=\inf\supp\psi$ and $b:=\sup\supp\psi$, so that
$\tfrac12\le a<b\le1$. Then:
\begin{enumerate}
\item[(i)] every zero on the critical line contributes a non-negative amount:
for $\rho=\tfrac12+i\gamma$ one has $1-\bar\rho=\rho$, so the term is
$|\widehat g_L(\rho)|^2\ge0$;
\item[(ii)] the resonant off-line pair contributes a \emph{positive} amount:
the two zeros $\rho_0$ and $1-\bar\rho_0=(1-\beta)+i\gamma_0$ together
contribute $2\widehat g_L(\rho_0)\widehat g_L(1-\bar\rho_0)$ with
\begin{equation}\label{eq:resonantpair}
  \widehat g_L(\rho_0)=\int\psi(t/L)e^{\delta t}dt>0,\qquad
  \widehat g_L(1-\bar\rho_0)=\int\psi(t/L)e^{-\delta t}dt>0,
\end{equation}
both factors being positive real numbers, whose exponential scales are
\begin{equation}\label{eq:laplacescales}
  \lim_{L\to\infty}\frac1L\log\widehat g_L(\rho_0)=\delta b,\qquad
  \lim_{L\to\infty}\frac1L\log\widehat g_L(1-\bar\rho_0)=-\delta a,
\end{equation}
so that the resonant pair contributes a positive quantity of exponential scale
\[
  \lim_{L\to\infty}\frac1L
   \log\bigl(\widehat g_L(\rho_0)\widehat g_L(1-\bar\rho_0)\bigr)
   =\delta(b-a)\ \le\ \frac\delta2 ,
\]
with equality exactly when $\supp\psi$ spans the whole of $[1/2,1]$. Only the
exponential scale is asserted: sub-exponential factors depend on the order of
vanishing of $\psi$ at its endpoints and are not claimed here;
\item[(iii)] the mismatched pair $\{\bar\rho_0,1-\rho_0\}$ contributes
$2\Rel[\widehat g_L(\bar\rho_0)\overline{\widehat g_L(1-\rho_0)}]$, which is
oscillatory and whose sign is not determined here, but which is bounded in
modulus by the resonant contribution, since
$|\widehat g_L(\bar\rho_0)|\le\widehat g_L(\rho_0)$ and
$|\widehat g_L(1-\rho_0)|\le\widehat g_L(1-\bar\rho_0)$;
\item[(iv)] the phenomenon is structural, not a matter of choosing $\psi$: for
any $g(t)=w(t)e^{-i\gamma_0t}$ with $w\ge0$,
\[
  \widehat g(\rho_0)=\int w\,e^{\delta t}dt\ \ge\ \int w\,e^{-\delta t}dt
   =\widehat g(1-\bar\rho_0)>0,
\]
so a non-negative amplitude can never make $\widehat g$ peak at
$1-\bar\rho_0$ rather than at $\rho_0$; more generally, for any real tilt
$w(t)=\psi(t/L)e^{\alpha t}$ the two exponents are $\alpha+\delta$ and
$\alpha-\delta$, and $\alpha+\delta>\alpha-\delta$ always. Peaking at
$1-\bar\rho_0$ therefore requires cancellation in $\int we^{\delta t}dt$, i.e.
an amplitude that changes sign or phase;
\item[(v)] the conclusion survives primitivization. With
$\widehat{\Pi g}(s)=s(1-s)\widehat g(s)$ and
$A:=\rho_0(1-\rho_0)=\tfrac14-(\delta+i\gamma_0)^2$, the factor at the second
resonant point $s=1-\bar\rho_0$ is $(1-\bar\rho_0)\bar\rho_0=\overline A$, so
\[
  \widehat{\Pi g_L}(\rho_0)\overline{\widehat{\Pi g_L}(1-\bar\rho_0)}
   +\mathrm{c.c.}=2\Rel\bigl(A^2\bigr)\,
   \widehat g_L(\rho_0)\widehat g_L(1-\bar\rho_0),
\]
and with $A=\tfrac14-\delta^2+\gamma_0^2-2i\delta\gamma_0$,
\[
  \Rel(A^2)=\bigl(\tfrac14-\delta^2+\gamma_0^2\bigr)^2-4\delta^2\gamma_0^2>0
  \qquad\text{whenever }|\gamma_0|>2\delta,
\]
in particular for every $|\gamma_0|\ge1$, and a fortiori for every nontrivial
zero of $\zeta$, all of which have $|\gamma|>14$; here $0<\delta<\tfrac12$ is
used. So the resonant pair of the primitivized test is again positive.
\end{enumerate}
\end{lemma}

\begin{proof}
(i) is $1-\bar\rho=\rho$ on the line.

(ii) Substituting \eqref{eq:singlefreqtest} into \eqref{eq:zerosum}: at
$s=\rho_0=\tfrac12+\delta+i\gamma_0$ the exponential is
$e^{-i\gamma_0t}e^{(\delta+i\gamma_0)t}=e^{\delta t}$, and at
$s=1-\bar\rho_0=\tfrac12-\delta+i\gamma_0$ it is $e^{-\delta t}$. Both
integrals have non-negative integrand, hence are positive reals, and the pair
$\{\rho_0,1-\bar\rho_0\}$ contributes
\[
  \widehat g_L(\rho_0)\overline{\widehat g_L(1-\bar\rho_0)}
   +\widehat g_L(1-\bar\rho_0)\overline{\widehat g_L(\rho_0)}
   =2\widehat g_L(\rho_0)\widehat g_L(1-\bar\rho_0)>0 .
\]
For the exponential scales \eqref{eq:laplacescales}, substitute $t=Lu$:
\[
  \widehat g_L(\rho_0)=L\int_{1/2}^{1}\psi(u)e^{\delta Lu}\,du,\qquad
  \widehat g_L(1-\bar\rho_0)=L\int_{1/2}^{1}\psi(u)e^{-\delta Lu}\,du .
\]
For the upper bound in the first,
$L\int\psi(u)e^{\delta Lu}du\le L\|\psi\|_\infty e^{\delta Lb}$, since
$\psi$ vanishes for $u>b$; taking $\frac1L\log$ gives
$\le\delta b+\frac{\log(L\|\psi\|_\infty)}L$. For the lower bound, fix
$\eta>0$; since $b=\sup\supp\psi$ there is $u_0\in(b-\eta,b)$ with
$\psi(u_0)>0$, and by continuity an interval $I\ni u_0$ of length $\ell>0$ on
which $\psi\ge\psi(u_0)/2$, which we may take inside $(b-\eta,b)$. Then
$L\int\psi e^{\delta Lu}du\ge L\tfrac12\psi(u_0)\ell\,e^{\delta L(b-\eta)}$,
so $\liminf\frac1L\log\widehat g_L(\rho_0)\ge\delta(b-\eta)$ for every
$\eta>0$. The two bounds give the first limit; the second is identical with
$u\mapsto-u$, the supremum of $-\delta u$ on $\supp\psi$ being $-\delta a$.

(iii) $\widehat g_L(\bar\rho_0)=\int\psi(t/L)e^{\delta t}e^{-2i\gamma_0t}dt$
and $\widehat g_L(1-\rho_0)=\int\psi(t/L)e^{-\delta t}e^{-2i\gamma_0t}dt$;
bounding the integrands in modulus gives the stated bounds, since $\psi\ge0$.
Non-stationary phase with a fixed $\gamma_0$ gains only a constant per
integration by parts, so no better bound is claimed.

(iv) is the displayed inequality $e^{\delta t}\ge e^{-\delta t}$ for $t>0$ and
the monotonicity of $\alpha\mapsto\alpha\pm\delta$.

(v) $\Pi=-\partial_t^2+\tfrac14$ has symbol $s(1-s)$ in the dictionary
\eqref{eq:zerosum}, so $\widehat{\Pi g}(s)=s(1-s)\widehat g(s)$; at $s=\rho_0$
the factor is $A$ and at $s=1-\bar\rho_0$ it is
$(1-\bar\rho_0)\bar\rho_0=\overline A$. Conjugating the second factor in the
Hermitian pairing gives $A\cdot A=A^2$, so the pair term acquires the real
coefficient $2\Rel(A^2)$. Writing
$A=\tfrac14-\delta^2+\gamma_0^2-2i\delta\gamma_0$, we get
\[
  \Rel(A^2)=\bigl(\tfrac14-\delta^2+\gamma_0^2\bigr)^2-4\delta^2\gamma_0^2
   =\bigl(\tfrac14-\delta^2+\gamma_0^2-2\delta|\gamma_0|\bigr)
    \bigl(\tfrac14-\delta^2+\gamma_0^2+2\delta|\gamma_0|\bigr),
\]
and the first factor equals $(|\gamma_0|-\delta)^2+\tfrac14-2\delta^2>0$ for
$|\gamma_0|>2\delta$ with $0<\delta<\tfrac12$, while the second is then
positive as well.
\end{proof}

\begin{remark}\label{rem:singlefreqscope}
Lemma~\ref{lem:singlefreq} computes the contribution of one resonant quadruple to a Hermitian
zero-sum against a single-frequency test on a moving window, and finds it to be
$O(e^{\delta L/2})$ with the resonant pair \emph{positive}. It decides nothing
about the location of any zero, and it makes no assertion about the sign of the
full sum. Part (iv) says that a non-negative amplitude, or any real
exponential tilt of one, cannot reverse the peaking; producing a test peaked at
$1-\bar\rho_0$ requires a sign- or phase-changing amplitude, which destroys the
sup-norm control that Lemmas \ref{lem:traceontinuity}--\ref{lem:C1control}
would need.
\end{remark}

\begin{thebibliography}{99}

\bibitem{Companion}
D.~A. Trejo Pizzo,
\emph{A correspondence algebra with a canonical trace on the idele class
group, its Tate symbol and its numerical form}, preprint.

\bibitem{Illusie}
L.~Illusie,
\emph{Complexe cotangent et d\'eformations I},
Lecture Notes in Mathematics 239, Springer, Berlin, 1971.

\bibitem{IK}
H.~Iwaniec and E.~Kowalski,
\emph{Analytic Number Theory},
American Mathematical Society Colloquium Publications 53, AMS, Providence, RI,
2004.

\bibitem{Lawson}
M.~V. Lawson,
\emph{Inverse Semigroups: The Theory of Partial Symmetries},
World Scientific, Singapore, 1998.

\bibitem{Stacks}
The Stacks Project Authors,
\emph{The Stacks Project}, \texttt{https://stacks.math.columbia.edu}, 2026;
Tag 08UQ.

\bibitem{Tate}
J.~T. Tate,
\emph{Fourier analysis in number fields and Hecke's zeta-functions},
thesis, Princeton University, 1950; reprinted in
J.~W.~S. Cassels and A.~Fr\"ohlich (eds.), \emph{Algebraic Number Theory},
Academic Press, London, 1967, 305--347.

\bibitem{Weil1948}
A.~Weil,
\emph{Sur les courbes alg\'ebriques et les vari\'et\'es qui s'en d\'eduisent},
Actualit\'es Sci. Ind. 1041, Hermann, Paris, 1948.

\bibitem{Weil1952}
A.~Weil,
\emph{Sur les ``formules explicites'' de la th\'eorie des nombres premiers},
Comm. S\'em. Math. Univ. Lund (1952), 252--265.

\bibitem{P45} D.~A.~Trejo Pizzo,
\emph{An explicit formula and a comparison theorem for the canonical trace on a correspondence algebra of the idele class group}, preprint.

\bibitem{P47} D.~A.~Trejo Pizzo,
\emph{Weil positivity on a finite window: duality, certificates and a finite reduction}, preprint.

\bibitem{P48} D.~A.~Trejo Pizzo,
\emph{A map of the routes to Weil positivity: eight arithmetic inputs, eight carriers, and the specification of a closing argument}, preprint.

\end{thebibliography}

\end{document}
